% !TEX program = pdflatex
\documentclass[11pt]{amsart}
% Omit volatile PDF metadata so clean pdfLaTeX builds can be compared exactly.
\ifdefined\pdfinfoomitdate\pdfinfoomitdate=1\fi
\ifdefined\pdftrailerid\pdftrailerid{}\fi
\ifdefined\pdfsuppressptexinfo\pdfsuppressptexinfo=15\fi
\usepackage[T1]{fontenc}
\usepackage{amsmath,amssymb,amsthm,mathtools}
\usepackage{enumitem}
\usepackage{etoolbox,needspace}
\usepackage{array,booktabs,tabularx,longtable}
\usepackage{microtype}
\usepackage[margin=1.05in]{geometry}
\usepackage{aliascnt}
\usepackage{xcolor}
\definecolor{KDLink}{RGB}{24,62,102}
\definecolor{KDCite}{RGB}{35,82,67}
\definecolor{KDUrl}{RGB}{38,74,88}
\usepackage[
  pdftitle={Two-term dg kernels and bounded profiles of universal localisations},
  pdfauthor={Marouane Lamharzi Alaoui},
  pdfsubject={Prescribed two-term dg kernels, t-exact optimality, and bounded profiles},
  colorlinks=true,
  linkcolor=KDLink,
  citecolor=KDCite,
  urlcolor=KDUrl
]{hyperref}
\usepackage[capitalize,noabbrev]{cleveref}
\setlength{\emergencystretch}{2em}
\clubpenalty=10000
\widowpenalty=10000
\displaywidowpenalty=10000
\brokenpenalty=10000
\raggedbottom

\makeatletter
\renewcommand{\tocsubsection}[3]{%
  \indentlabel{%
    \if\relax\detokenize{#2}\relax
      \hspace*{5pc}% align unnumbered entries with subsection titles
    \else
      \ignorespaces#1 #2.\quad
    \fi
  }#3}
\makeatother

\pretocmd{\section}{\Needspace{5\baselineskip}}{}{}
\pretocmd{\subsection}{\Needspace{4\baselineskip}}{}{}

\newtheorem{theorem}{Theorem}[section]
\newaliascnt{proposition}{theorem}
\newtheorem{proposition}[proposition]{Proposition}
\aliascntresetthe{proposition}
\newaliascnt{lemma}{theorem}
\newtheorem{lemma}[lemma]{Lemma}
\aliascntresetthe{lemma}
\newaliascnt{corollary}{theorem}
\newtheorem{corollary}[corollary]{Corollary}
\aliascntresetthe{corollary}
\newaliascnt{question}{theorem}
\newtheorem{question}[question]{Question}
\aliascntresetthe{question}
\theoremstyle{definition}
\newaliascnt{definition}{theorem}
\newtheorem{definition}[definition]{Definition}
\aliascntresetthe{definition}
\newaliascnt{example}{theorem}
\newtheorem{example}[example]{Example}
\aliascntresetthe{example}
\newaliascnt{notation}{theorem}
\newtheorem{notation}[notation]{Notation}
\aliascntresetthe{notation}
\theoremstyle{remark}
\newaliascnt{remark}{theorem}
\newtheorem{remark}[remark]{Remark}
\aliascntresetthe{remark}

\AtBeginEnvironment{theorem}{\Needspace{4\baselineskip}}
\AtBeginEnvironment{proposition}{\Needspace{4\baselineskip}}
\AtBeginEnvironment{lemma}{\Needspace{4\baselineskip}}
\AtBeginEnvironment{corollary}{\Needspace{4\baselineskip}}
\AtBeginEnvironment{question}{\Needspace{4\baselineskip}}
\AtBeginEnvironment{definition}{\Needspace{4\baselineskip}}
\AtBeginEnvironment{example}{\Needspace{4\baselineskip}}
\AtBeginEnvironment{notation}{\Needspace{4\baselineskip}}
\AtBeginEnvironment{remark}{\Needspace{4\baselineskip}}
\AtBeginEnvironment{proof}{\Needspace{3\baselineskip}}

\crefname{theorem}{theorem}{theorems}
\Crefname{theorem}{Theorem}{Theorems}
\crefname{lemma}{lemma}{lemmas}
\Crefname{lemma}{Lemma}{Lemmas}
\crefname{proposition}{proposition}{propositions}
\Crefname{proposition}{Proposition}{Propositions}
\crefname{corollary}{corollary}{corollaries}
\Crefname{corollary}{Corollary}{Corollaries}
\crefname{definition}{definition}{definitions}
\Crefname{definition}{Definition}{Definitions}
\crefname{remark}{remark}{remarks}
\Crefname{remark}{Remark}{Remarks}
\crefname{example}{example}{examples}
\Crefname{example}{Example}{Examples}
\crefname{question}{question}{questions}
\Crefname{question}{Question}{Questions}
\crefname{notation}{notation}{notations}
\Crefname{notation}{Notation}{Notations}

\DeclareMathOperator{\Loc}{Loc}
\DeclareMathOperator{\thick}{thick}
\DeclareMathOperator{\per}{per}
\DeclareMathOperator{\Cone}{Cone}
\DeclareMathOperator{\Tor}{Tor}
\DeclareMathOperator{\Ext}{Ext}
\DeclareMathOperator{\Hom}{Hom}
\DeclareMathOperator{\End}{End}
\DeclareMathOperator{\RHom}{RHom}
\DeclareMathOperator{\add}{add}
\DeclareMathOperator{\rad}{rad}
\DeclareMathOperator{\Kar}{Kar}
\DeclareMathOperator{\id}{id}
\DeclareMathOperator{\gl}{gl.dim}
\DeclareMathOperator{\im}{im}
\DeclareMathOperator{\coker}{coker}
\DeclareMathOperator{\Calk}{Calk}
\DeclareMathOperator{\Fact}{Fact}
\DeclareMathOperator{\Obs}{Obs}

\newcommand{\Db}{D^{\mathrm b}}
\newcommand{\Dsg}{D_{\mathrm{sg}}}
\newcommand{\Dfd}{D_{\mathrm{fd}}}
\newcommand{\Dfg}{D^{\mathrm b}_{\mathrm{fg}}}
\newcommand{\Kb}{K^{\mathrm b}}
\newcommand{\Kc}{K^c}
\newcommand{\R}{\mathbf R}
\newcommand{\Ltensor}{\otimes^{\mathbf L}}
\newcommand{\DeltaB}{\Delta_{\mathrm b}}
\newcommand{\Defsg}{\operatorname{Def}_{\mathrm{sg}}}
\newcommand{\modcat}{\operatorname{mod}}
\newcommand{\Mod}{\operatorname{Mod}}
\newcommand{\op}{\mathrm{op}}
\newcommand{\eps}{\varepsilon}
\newcommand{\cT}{\mathcal T}
\newcommand{\cP}{\mathcal P}
\newcommand{\cK}{\mathcal K}
\newcommand{\cL}{\mathcal L}
\newcommand{\cC}{\mathcal C}
\newcommand{\cH}{\mathcal H}
\newcommand{\cF}{\mathcal F}
\newcommand{\cA}{\mathcal A}
\newcommand{\cB}{\mathcal B}
\newcommand{\cE}{\mathcal E}
\newcommand{\cD}{\mathcal D}
\newcommand{\vect}{\operatorname{vect}}
\newcommand{\kerob}{\ker_{\mathrm{ob}}}


\DeclareMathOperator{\HMF}{\mathbf{HMF}}
\DeclareMathOperator{\pd}{pd}
\DeclareMathOperator{\ann}{ann}
\DeclareMathOperator{\Proj}{Proj}
\newcommand{\Dk}{D_k}
\newcommand{\cU}{\mathcal U}
\newcommand{\cQ}{\mathcal Q}
\newcommand{\cG}{\mathcal G}
\newcommand{\sD}{\mathscr D}
\newcommand{\VV}{\mathbb V}
\newcommand{\PP}{\mathbb P}
\newcommand{\Gam}{\Gamma}
\newcommand{\Prof}{\mathcal P}
\newcommand{\bbZ}{\mathbb Z}
\usepackage{mathrsfs}

\title[Two-term dg kernels and bounded profiles]{Two-term dg kernels and bounded profiles\\ of universal localisations}
\author{Marouane Lamharzi Alaoui}
\thanks{Independent researcher. Correspondence: \texttt{marouane@masonai.fr}.}
\date{}
\subjclass[2020]{18E30, 16E35, 18G80, 16S85, 16G10, 13D07}
\keywords{universal localisation, dg algebra, bounded profile, singularity category, one-point extension, support variety}


\begin{document}
\begin{abstract}
We study which dg algebras occur as kernels of universal localisations and which bounded subcategories their embeddings determine. Every proper dg algebra over a field with cohomology in degrees \(-1\) and \(0\) is realised as the dg endomorphism algebra of a perfect kernel generator for a stably flat universal localisation of finite-dimensional algebras at one map. The target is projective on both sides, and the marked kernel embedding is standard t-exact with full finite-dimensional profile. The construction uses a finite crossed extension and identifies the prescribed dg algebra by a strict dg map; a nineteen-dimensional source realises a non-formal kernel. Conversely, a standard t-exact marked realisation of a connective dg algebra by an ordinary homological ring epimorphism forces a two-term projective generator, without assuming target projectivity.

The bounded profile is detected by the coreflected injective cogenerator. For these realisations, one-point extensions preserve the kernel and realise every far-away orthogonal cut of the full profile. For complete-intersection and finite \(p\)-group kernels, these cuts are classified as marked support-avoidance subcategories; intermediate cuts can have no classical generator. We also compare the bounded defect with an ambient singular quotient, which identifies with the singularity-functor kernel under bounded extension and restriction hypotheses. Classical compact factorisation proves full faithfulness; a simultaneous negative-\(K\) kernel measures the possible discrepancy before idempotent completion.
\end{abstract}
\maketitle
\setcounter{tocdepth}{1}
\tableofcontents

\section{Introduction}\label{sec:introduction}
A universal localisation \(\lambda:R\to S\) forces specified maps between finitely generated projective modules to become invertible. When it is stably flat, its derived kernel is generated by their cones \cite{Schofield85,NeemanRanicki04,Neeman07}. We ask two related questions: which dg algebras can describe this kernel, and which of their modules have bounded finite-dimensional cohomology after embedding in \(D(R)\)? The second question depends on the embedding. The two module sides of the AHKLY dual-numbers example have the same intrinsic kernel algebra and different bounded defects (\Cref{thm:ahkly-left,thm:dual-numbers}).

For a finite-dimensional algebra \(R\), a compact kernel generator \(T\) gives a dg Morita equivalence
\[
 B=\RHom_R(T,T)^{\op},\qquad
 F=T\Ltensor_B-:D(B)\xrightarrow{\sim}\cK_\lambda.
\]
Its \emph{bounded profile} is \(\Prof(F)=F^{-1}\Db(R\text{-}\modcat)\). The bounded defect discards the perfect objects in this profile:
\[
 \DeltaB(\lambda)=\bigl(\cK_\lambda\cap\Db(R\text{-}\modcat)\bigr)/\Kc_\lambda
 \simeq\Prof(F)/\per B.
\]
All profile classifications below retain this specified identification, called a \emph{marking}.

\subsection*{Prescribed kernels and an optimality theorem}
The main construction and its converse give the following statement.
\begin{theorem}[Two-term realisation and t-exact optimality]\label{thm:intro-Q}
Let \(B\) be a proper connective dg algebra over a field. It admits a realisation as the dg endomorphism algebra of a perfect generator of the kernel of a homological ring epimorphism of finite-dimensional algebras, with target projective on both sides and standard t-exact marked embedding, if and only if \(H^{<-1}(B)=0\). The realisation can be chosen to be universal localisation at one map, with full bounded profile.

More generally, for an ordinary homological ring epimorphism, a perfect kernel generator with connective dg endomorphism algebra and standard t-exact marked embedding has a projective representative in degrees \([-1,0]\). The epimorphism is universal localisation at its differential. Target projectivity is not required for this implication.
\end{theorem}
\begin{proof}[Proof references]
Combine \Cref{thm:realisation,thm:finite-crossed,thm:texact-optimality,cor:texact-universal}.
\end{proof}

For a finite two-term dg model \(D=[M\xrightarrow{\delta}A]\), the source and target are explicit:
\[
 R_\delta=\begin{pmatrix}A&A\\M&A\end{pmatrix},\qquad
 \lambda_\delta(a,x,m,b)=\begin{pmatrix}a&x\\\delta(m)&b\end{pmatrix}\in M_2(A).
\]
The pairings in the Morita context are induced by \(\delta\). A strict dg map identifies \(D\) with the cone endomorphism algebra; the two vertex components prove t-exactness. Thus the construction retains higher structure, rather than only its cohomology algebra. Finite crossed extensions realise every prescribed class in \(HH^3(H^0B,H^{-1}B)\), and \Cref{ex:nonformal} gives a nineteen-dimensional source with nonzero \(m_3\). Finite-model existence itself is prior work of Raedschelders--Stevenson \cite[Appendix~A]{Goodbody24}. The converse follows from the cohomological amplitude of the ordinary localisation triangle. A negative Calabi--Yau obstruction and an explicit Schur example delimit its scope (\Cref{prop:negative-cy-obstruction,ex:sphere-schur}).

\subsection*{Changing the bounded profile}
For \(E\in\Dfd(B)\), put
\[
 \cU_E=\{Y\in\Dfd(B):\RHom_B(Y,E)\text{ has finite-dimensional total cohomology}\}.
\]
This is a far-away left orthogonal in the sense of \cite{KostasPsaroudakisVitoria25}. The canonical detector theorem identifies the bounded profile with \(\cU_{D_kT}\): intrinsically in the kernel, the detector is the coreflection of the injective cogenerator \(DR\). One-point extension by a finite right \(R\)-module \(W\) then gives
\[
 R[W]\longrightarrow S[W\otimes_RS],\qquad
 \Prof(\lambda')=\Prof(\lambda)\cap\cU_{\Theta(D_kW)},
 \quad\Theta=\RHom_R(T,-),
\]
provided \(\Tor_i^R(W,S)=0\) for every \(i>0\); that condition is also necessary for stable flatness of the extended localisation (\Cref{thm:ambient}). Applying this to the realisation above yields exactly all marked defects
\[
 \cU_E/\per B,\qquad E\in\Dfd(B),
\]
for every proper two-term \(B\) (\Cref{cor:thin-all-profiles}). Smaller cuts need not retain the seed's full profile or t-exactness.

For an ordinary kernel algebra \(C\), each detector can be replaced by a module. Complete-intersection and finite \(p\)-group support theory then identifies the defects with support-avoidance subcategories \(\sD_Z\) for closed \(Z\). If an irreducible support component meets both \(Z\) and its complement, \(\sD_Z\) has no classical generator. It therefore cannot be the singularity category of a finite-dimensional algebra or a proper connective one-object dg algebra (\Cref{thm:no-generator}). Distinct marked cuts can nevertheless be abstractly equivalent: exchanging the coordinates of \(k[x,y]/(x^2,y^2)\) gives an explicit example (\Cref{ex:support-symmetry}).

\subsection*{Classical comparison}
The ambient quotient \(\Defsg(\lambda)=\thick(\per R,\Kb(\lambda))/\per R\) is defined without assuming that derived localisation preserves bounded objects. Under \textup{(B1)--(B3)}, it is the object-kernel of the induced singularity functor (\Cref{thm:actual-kernel}). The fully faithful comparison \(\DeltaB(\lambda)\to\Defsg(\lambda)\) is a consequence of classical compact factorisation and Krause's localisation-of-subcategories criterion \cite[Lemma~4.7.1(1), Theorem~7.2.1(1)]{Krause10}. We give that short proof without algebraicity or enhancement hypotheses in \Cref{sec:rigidity}. The dg explanation using Efimov's Calkin exactness is retained in \Cref{sec:obstruction} as an enhanced account of this classical statement.

The only possible pre-completion discrepancy is a summand obstruction. \Cref{thm:negative-k-obstruction} identifies it with the kernel of the two maps from \(K_{-1}(\Kc_\lambda)\) to \(K_{-1}(\Kb(\lambda))\) and \(K_{-1}(R)\). It vanishes in every example constructed here; no example with nonzero obstruction is supplied. The heart-amplitude criterion, one-sided bounded dictionary, and explicit Nakayama calculations provide further tests for when the intrinsic dg category computes the ambient bounded profile.

\subsection*{Organisation}
After conventions, \Cref{sec:realisation} proves the realisation theorem and converse. \Cref{sec:detector,sec:ambient,sec:ordinary,sec:support} develop detectors, profile changes, and support applications. \Cref{sec:dg-models,sec:heart-amplitude,sec:dictionary,sec:periodic-ladder,sec:stratifying,sec:nakayama} give the intrinsic criteria and examples. \Cref{sec:actual-kernel,sec:rigidity} treat the singularity kernel and classical comparison; \Cref{sec:literature,sec:limitations} discuss prior work and the remaining questions. The appendices contain the auxiliary factorisation arguments, dg enhancement, negative-\(K\) proofs, bounded absorption, and computational protocol.


\section{Conventions and standing hypotheses}\label{sec:preliminaries}

For a ring \(R\), \(R\text{-}\Mod\) denotes all left \(R\)-modules and \(R\text{-}\modcat\) denotes finitely generated left \(R\)-modules.  We write \(D(R)=D(R\text{-}\Mod)\).  When \(R\) is left noetherian, \(\Db(R\text{-}\modcat)\) is the bounded derived category of finitely generated left modules.  If a right-module calculation is used, it is explicitly described as a calculation over the opposite ring.

The dg enhancement in \Cref{sec:obstruction} is over a fixed commutative ground ring \(\mathsf k\), taken to be \(\mathbb Z\) when none is specified. Its morphism complexes are replaced by \(K\)-flat models. Algebraic triangulated categories are stable categories of Frobenius exact categories. These hypotheses belong to the enhanced appendix, not to the classical comparison in \Cref{sec:rigidity}.

Throughout, \emph{before idempotent completion} means at the level of the displayed Verdier quotients themselves; otherwise the application of \(\Kar(-)\) is shown explicitly.  The idempotent completion carries its canonical triangulated structure by Balmer--Schlichting~\cite[Theorem~1.5]{BalmerSchlichting01}.

Complexes are cohomologically graded, \(X[1]^n=X^{n+1}\), with differential \(-d_X\).  If \(\sigma:P\to Q\) is a map of projectives, \(\Cone(\sigma)=[P\to Q]\) has \(P\) in degree \(-1\) and \(Q\) in degree zero.  Paths are composed from left to right in written quiver notation: if \(a:i\to j\) and \(b:j\to l\), then \(ab:i\to l\).

For finite \(\Sigma\), set
\[
 T_\Sigma=\bigoplus_{\sigma\in\Sigma}\Cone(\sigma),\qquad
 \Lambda=\RHom_R(T_\Sigma,T_\Sigma)^{\op}.
\]
Thus \(T_\Sigma\) is a left \(R\)-complex and a right dg \(\Lambda\)-module, and the realisation functor is
\[
 \Psi=T_\Sigma\Ltensor_\Lambda-:D(\Lambda)\to D(R).
\]
The opposite in the definition of \(\Lambda\) is fixed throughout.

The bounded hypotheses are:
\begin{enumerate}[label=\textup{(B\arabic*)}]
\item \(R\) and \(S\) are left noetherian;
\item \(S_R\) has finite Tor-amplitude;
\item restriction of scalars sends \(\Db(S\text{-}\modcat)\) to \(\Db(R\text{-}\modcat)\).
\end{enumerate}

\begin{definition}[Object-kernel]\label{def:object-kernel}
For an exact functor \(F:\cD\to\cE\), write
\[
 \kerob(F)=\{X\in\cD\mid F(X)\cong0\}
\]
for the full triangulated subcategory on the objects sent to zero.  It is thick.  Throughout, an \emph{actual kernel} means this embedded object-kernel in the displayed triangulated category; it does not mean an enhanced homotopy fibre.
\end{definition}

\begin{definition}[Bounded and singular defects]\label{def:singularity-defect}
Let \(\lambda:R\to S\) be a ring map with \(R\) left noetherian, and put
\[
 \cK_\lambda=\kerob(S\Ltensor_R-:D(R)\to D(S)),
 \qquad
 \Kb(\lambda)=\cK_\lambda\cap\Db(R\text{-}\modcat),
 \qquad
 \Kc_\lambda=\cK_\lambda\cap\per(R).
\]
The \emph{bounded defect} and \emph{singularity defect} are
\[
 \DeltaB(\lambda)=\Kb(\lambda)/\Kc_\lambda,
 \qquad
 \Defsg(\lambda)=\thick(\per(R),\Kb(\lambda))/\per(R).
\]
These Verdier quotients are defined without \textup{(B2)} or \textup{(B3)}.  Those hypotheses enter only when extension and restriction are required to preserve the bounded categories and when \(\Defsg(\lambda)\) is identified with an object-kernel.
\end{definition}

No noetherian hypothesis is used in \Cref{sec:obstruction,sec:rigidity} or in the thick-scale formula of \Cref{thm:relative-dg-formula}.  The bounded defects and the specialization in \Cref{cor:relative-dg-bounded} require \(R\) to be left noetherian.  Conditions \textup{(B1)--(B2)} give the bounded extension functor and induced singularity functor, while \textup{(B3)} is additionally used for bounded restriction.  The actual object-kernel theorem also requires \(\lambda\) to be a homological ring epimorphism; conditions \textup{(B1)--(B3)} alone do not imply homologicality.  Finiteness of \(\Sigma\) is needed only for a one-object kernel dga; the compact-kernel comparison permits a set of compact generators.

\subsection*{The AHKLY algebra in fixed bases}
Let
\[
 A=kQ/(\beta^2,\alpha\beta),\qquad
 Q:\;1\xrightarrow{\alpha}2,\quad \beta:2\to2,
\]
and let \(S=M_2(k)\), with
\[
 e_i\mapsto E_{ii},\qquad \alpha\mapsto E_{12},\qquad \beta\mapsto0.
\]
In left modules the inverted presentation is
\[
 \sigma_L:Ae_1\longrightarrow Ae_2,\qquad u\longmapsto u\alpha,
\]
with bases \(Ae_1=(e_1)\) and \(Ae_2=(e_2,\alpha,\beta)\).  In right modules the opposite presentation is
\[
 \sigma_R:e_2A\longrightarrow e_1A,\qquad v\longmapsto\alpha v,
\]
with \(e_2A=(e_2,\beta)\) and \(e_1A=(e_1,\alpha)\).  Every side-sensitive assertion below refers to these two maps.

\subsection*{Kernel dg algebras, profiles, and markings}
In \Cref{sec:detector,sec:ambient,sec:realisation,sec:ordinary,sec:support} the ground ring is a field \(k\), ordinary ambient algebras are finite-dimensional over \(k\) unless otherwise stated, and
\[
 \Db(R\text{-}\modcat)=\Bigl\{X\in D(R):\ \bigoplus_nH^n(X)\text{ is finite-dimensional}\Bigr\}.
\]
We write \(D_k=\Hom_k(-,k)\).  For a right \(R\)-module \(W\), the dual \(D_kW\) is a left \(R\)-module via \((rf)(w)=f(wr)\); the injective cogenerator of left modules is \(DR=D_k(R_R)\), with \((rf)(s)=f(sr)\).  The graded opposite of a dg algebra carries the Koszul sign, \(a\cdot^{\op}b=(-1)^{|a||b|}ba\).

If \(T\) is a compact generator of a compactly generated localising subcategory \(\cK\subseteq D(R)\), we put
\[
 B=\RHom_R(T,T)^{\op},\qquad F=T\Ltensor_B-:D(B)\xrightarrow{\sim}\cK,\qquad\Theta=\RHom_R(T,-):D(R)\to D(B);
\]
for \(T=T_\Sigma\) this is the raw kernel dga \(\Lambda\) and its realisation functor \(\Psi\) above, with notation suited to the profile problem.  Thus \(T\) is an \(R\)-\(B\) bimodule and \(\Theta X\) is a left dg \(B\)-module by precomposition.  When \(T\) is a bounded complex of finite-dimensional projective \(R\)-modules, \(B\) is proper, that is, \(\dim_kH^*(B)<\infty\).  For a proper dg algebra \(B\), \(\Dfd(B)\subseteq D(B)\) is the full subcategory of dg modules with finite-dimensional total cohomology and \(\Dsg(B)=\Dfd(B)/\per B\).

\begin{definition}[Far-away orthogonal, profile, marking]\label{def:profile}
For \(E\in\Dfd(B)\) put
\[
 \cU_E=\{Y\in\Dfd(B):\RHom_B(Y,E)\in\Db(k\text{-}\modcat)\}.
\]
This is a thick subcategory of \(\Dfd(B)\) containing \(\per B\): the condition is stable under shifts, cones, and summands in \(Y\), and for \(Y=B\) it says that \(E\) is proper.  For an embedding \(F:D(B)\simeq\cK\subseteq D(R)\) with \(R\) finite-dimensional, the \emph{bounded profile} is
\[
 \Prof(F)=\{Y\in D(B):FY\in\Db(R\text{-}\modcat)\},
\]
so that \(\Prof(\Psi)=\Db_\lambda(A_\lambda)\) in the notation of \Cref{cor:relative-dg-bounded}.  A \emph{marking} of a kernel is a specified dg Morita identification \(D(B)\simeq\cK\); the classifications of \Cref{sec:ambient,sec:realisation,sec:ordinary,sec:support} concern full subcategories of the fixed marked category \(\Dfd(B)\) or \(\Dsg(B)\).  They are not classifications up to abstract triangle equivalence.
\end{definition}

\section{Two-term dg kernels}\label{sec:realisation}

\Cref{thm:ambient} will show that profiles can be cut by detectors: once a kernel dg algebra \(B\) has one realisation with full profile and projective localisation, every far-away orthogonal \(\cU_E\) is realised.  The substantive question is which dg algebras admit such a realisation at all.  This section gives an explicit realisation for every proper dg algebra with cohomology in degrees \(-1,0\), and proves that these are exactly the proper connective dg algebras admitting a standard t-exact marked kernel realisation.  The converse holds for ordinary homological ring epimorphisms without target projectivity.  Full profile alone has a wider scope, as \Cref{ex:marking-boundary} shows.

\begin{question}[Kernel realisation]\label{q:realisation}
Which proper cohomologically non-positive dg algebras \(B\) are, up to a specified dg Morita identification, kernel dg algebras of stably flat universal localisations of finite-dimensional algebras at finite \(\Sigma\) with full bounded profile?
\end{question}

By \Cref{thm:heart-amplitude}, for a silting marking full profile is equivalent to fixed finite Tor-amplitude of the heart bimodule; it is a property of the embedding, not of \(B\) alone.

\subsection{Necessary conditions}

\begin{proposition}[Thin generators]\label{prop:thin}
In every kernel as in \Cref{q:realisation}, the transported sum \(C_0\) of the presentation cones is a classical generator of \(\per B\) with \(\Hom(C_0,C_0[n])=0\) for \(n\notin[-1,1]\).  If \(C_0\) is silting, \(H^n\RHom(C_0,C_0)=0\) for \(n\notin[-1,0]\).
\end{proposition}

\begin{proof}
The ambient representative of \(C_0\) is a two-term complex of finite projectives in degrees \(-1,0\); its mapping complex has terms only in degrees \(-1,0,1\) and computes derived Hom.  The cones classically generate the compact kernel by \Cref{lem:compact-intersection}.
\end{proof}

The statement concerns the existence of a generator with this property; a silting generator obtained by shifting summands, as in \Cref{sec:nakayama}, need not have two-term amplitude.  The silting hypothesis cannot be dropped: for \(R=k\times k\), \(p=(1,0)\), the universal localisation at the zero map \(Rp\to Rp\) is the projection onto the second factor, stably flat with projective target, and its cone \(Rp[1]\oplus Rp\) has \(H^1\) of its endomorphism dg algebra of dimension one.

\begin{proposition}[Higher multiplications]\label{prop:higher-mult}
Let \(H\) be a minimal \(A_\infty\)-algebra with underlying graded space \(H^{-1}\oplus H^0\).  Then the only possible higher multiplication is
\[
 m_3:(H^0)^{\otimes3}\longrightarrow H^{-1},
\]
and in a strictly unital model it is a normalised Hochschild \(3\)-cocycle of the algebra \(\Gam=H^0\) with values in the bimodule \(U=H^{-1}\).
\end{proposition}

\begin{proof}
The operation \(m_n\) has degree \(2-n\).  For \(n\ge4\) all inputs have non-positive degree and the output has degree at most \(-2\), where \(H\) vanishes; for \(n=3\) a negative input forces the output below \(-1\).  The \(A_\infty\) relation in arity \(3\) makes \(m_2\) associative, and the relation in arity \(4\) on degree-zero inputs is exactly the Hochschild cocycle condition; compositions of two \(m_3\) vanish for degree reasons.  Minimal models exist over a field by the transfer theorem, see \cite[Section~3.3]{Keller01Ainf}.
\end{proof}

With the graded algebra and identity linear marking fixed, the class of \(m_3\) is the remaining invariant of such a structure.

\begin{lemma}[Gauge classification]\label{lem:gauge}
Two strictly unital minimal \(A_\infty\)-structures on \(\Gam\oplus U[1]\) with the same associative product on \(\Gam\) and the same bimodule structure on \(U\) are \(A_\infty\)-isomorphic by an isomorphism with identity linear part if and only if their \(m_3\) differ by a Hochschild coboundary.  Thus such structures are classified by \(HH^3(\Gam,U)\).
\end{lemma}

\begin{proof}
An \(A_\infty\)-morphism \(f\) with \(f_1=\id\) has components \(f_n\) of degree \(1-n\); on \(\Gam\oplus U[1]\) the only possibly nonzero higher component is \(f_2:\Gam^{\otimes2}\to U\).  The arity-\(2\) morphism equation is vacuous since \(m_1=0\), so \(m_2\) is unchanged; the arity-\(3\) equation reads \(m_3'-m_3=\pm d_{\mathrm{Hoch}}(f_2)\); every higher equation involves a composite with two negative-degree outputs and holds automatically.  Conversely any normalised \(2\)-cochain \(f_2\) defines such an isomorphism.
\end{proof}

\subsection{Two-term dg algebras as crossed bimodules}

\begin{lemma}[Two-term dg algebras]\label{lem:crossed}
A unital dg algebra \(D\) with \(D^i=0\) for \(i\notin\{-1,0\}\) is the same as an algebra \(A=D^0\), an \(A\)-bimodule \(M=D^{-1}\), and a bimodule map \(\delta:M\to A\) satisfying
\begin{equation}\label{eq:peiffer}
 \delta(m)m'=m\delta(m')\qquad(m,m'\in M);
\end{equation}
the product of two elements of \(M\) is zero.  Put \(\Gam=A/\delta(M)\) and \(U=\ker\delta\).  Then \(\delta(M)\) annihilates \(U\) on both sides, \(U\) is a \(\Gam\)-bimodule, \(H^0(D)=\Gam\), \(H^{-1}(D)=U\), and
\begin{equation}\label{eq:crossed-extension}
 0\longrightarrow U\longrightarrow M\xrightarrow{\ \delta\ }A\longrightarrow\Gam\longrightarrow0
\end{equation}
is a crossed extension of \(\Gam\) by \(U\) in the sense of Baues--Minian \cite{BauesMinian02}.
\end{lemma}

\begin{proof}
The products in degrees \((0,0)\), \((0,-1)\), \((-1,0)\) give the algebra and bimodule structures, and products in degree \(-2\) vanish.  The Leibniz rule with one factor of degree zero says that \(\delta\) is a bimodule map; applied to the zero product \(mm'\) it gives \(0=\delta(m)m'-m\delta(m')\).  Conversely these identities are the Leibniz rule in every degree pair, and the bimodule axioms give associativity.  For \(u\in U\) and \(m\in M\), \eqref{eq:peiffer} gives \(\delta(m)u=m\delta(u)=0\) and \(u\delta(m)=\delta(u)m=0\).
\end{proof}

\begin{lemma}[The minimal model of a two-term dg algebra]\label{lem:two-term-minimal}
Let \(D=[M\xrightarrow{\delta}A]\) be as in \Cref{lem:crossed}, with \(\Gam\), \(U\) finite-dimensional.  Choose a \(k\)-linear section \(s:\Gam\to A\) of the projection with \(s(1)=1\), a complement \(M=U\oplus M'\), and for \(x,y\in\Gam\) let \(\beta(x,y)\in M'\) be the unique element with
\begin{equation}\label{eq:beta}
 s(x)s(y)-s(xy)=\delta\beta(x,y).
\end{equation}
Then the transferred minimal \(A_\infty\)-structure on \(\Gam\oplus U[1]\) has \(m_2\) the product of \(\Gam\) and the bimodule structure of \(U\), and
\begin{equation}\label{eq:m3-cocycle}
 m_3(x,y,z)=x\beta(y,z)-\beta(xy,z)+\beta(x,yz)-\beta(x,y)z\in U,
\end{equation}
the characteristic cocycle of the crossed extension \eqref{eq:crossed-extension}.  Its class in \(HH^3(\Gam,U)\) is independent of all choices and is the characteristic class of \eqref{eq:crossed-extension}.
\end{lemma}

\begin{proof}
Let \(i:\Gam\oplus U[1]\to D\) be \(s\) on \(\Gam\) and the inclusion on \(U\); let \(p:D\to\Gam\oplus U[1]\) be the projection \(A\to\Gam\) on \(A\) and the projection to \(U\) along \(M'\) on \(M\); let \(h:D\to D[-1]\) vanish on \(s(\Gam)\) and on \(M\), and on \(\delta(M)\) let \(h(\delta m')=-m'\) for \(m'\in M'\).  Then \(pi=\id\) and \(ip-\id=dh+hd\): on \(a=s(x)+\delta(m')\) both sides equal \(-\delta(m')\), and on \(m=u+m'\) both sides equal \(-m'\).  The transfer formula gives
\[
 m_3(x,y,z)=p\bigl(h(s(x)s(y))\,s(z)\bigr)-p\bigl(s(x)\,h(s(y)s(z))\bigr)
 =p\bigl(s(x)\beta(y,z)-\beta(x,y)s(z)\bigr),
\]
using \(h(s(x)s(y))=h(s(xy))+h(\delta\beta(x,y))=-\beta(x,y)\).  Applying \(\delta\) to the right-hand side of \eqref{eq:m3-cocycle} and expanding with \eqref{eq:beta} gives zero, so that cocycle lies in \(U\); hence the \(M'\)-component of \(s(x)\beta(y,z)-\beta(x,y)s(z)\) equals \(\beta(xy,z)-\beta(x,yz)\), and projecting to \(U\) along \(M'\) yields \eqref{eq:m3-cocycle}.  Different choices of \(s\), \(M'\), or of the transfer data change the minimal structure by an \(A_\infty\)-isomorphism with identity linear part \cite[Section~3.3]{Keller01Ainf}, hence \(m_3\) by a coboundary (\Cref{lem:gauge}).  That \eqref{eq:m3-cocycle} is the characteristic cocycle of the crossed extension is the construction of \cite[Section~3]{BauesMinian02}.
\end{proof}

\subsection{Realisation of a strict two-term dg algebra}

\begin{theorem}[Realisation of two-term dg algebras]\label{thm:realisation}
Let \(D=[M\xrightarrow{\delta}A]\) be a finite-dimensional unital dg algebra concentrated in degrees \(-1,0\), and let
\begin{equation}\label{eq:R-delta}
 R_\delta=\begin{pmatrix}A&A\\M&A\end{pmatrix},
 \qquad
 (a,v,m,b)(a',v',m',b')=\bigl(aa'+v\delta(m'),\ av'+vb',\ ma'+bm',\ bb'+\delta(m)v'\bigr),
\end{equation}
the Morita context ring of the pairings \(v\otimes m\mapsto v\delta(m)\) and \(m\otimes v\mapsto\delta(m)v\).  Let \(p=e_{11}\), \(e=e_{22}\), let \(s\) be the upper-right copy of \(1_A\), and \(\sigma:R_\delta p\to R_\delta e\), \(x\mapsto xs\).  Then:
\begin{enumerate}[label=\textup{(\roman*)}]
\item\label{re:i} the ring map \(\lambda_\delta:R_\delta\to S=M_2(A)\), \((a,v,m,b)\mapsto\begin{pmatrix}a&v\\\delta(m)&b\end{pmatrix}\), is the universal localisation at \(\sigma\);
\item\label{re:ii} \(S_{R_\delta}\cong(pR_\delta)^{\oplus2}\) and \({}_{R_\delta}S\cong(R_\delta e)^{\oplus2}\); in particular \(\lambda_\delta\) is stably flat;
\item\label{re:iii} for \(T=\Cone(\sigma)\) there is a strict quasi-isomorphism of dg algebras \(D\to\End^\bullet_{R_\delta}(T)^{\op}\);
\item\label{re:iv} the Morita functor \(F=T\Ltensor_D-:D(D)\xrightarrow{\sim}\cK_{\lambda_\delta}\) is t-exact for the standard t-structures, with \(H^n(FY)\cong(0,H^nY)\) as \(R_\delta\)-modules through \(R_\delta\to\Gam\); it preserves and reflects finite-dimensional total cohomology, so the bounded profile is all of \(\Dfd(D)\); and the heart bimodule is \(T\Ltensor_D\Gam\simeq h\Gam\), where \(hN=(0,N)\) denotes the \(R_\delta\)-module supported on the lower-right corner, of Tor-amplitude \([0,0]\).
\end{enumerate}
Moreover \(\dim_kR_\delta=3\dim_kA+\dim_kM\) and \(\dim_kS=4\dim_kA\).
\end{theorem}

\begin{proof}
\emph{Associativity and the ring map.}  The diagonal products are products in \(A\); each off-diagonal associativity identity is one of the bimodule identities for the regular bimodule \(A\) or for \(M\), except the identity \(\bigl((m\cdot v')\cdot m''\bigr)=m\cdot(v'\cdot m'')\), which reads \(\delta(m)v'm''=mv'\delta(m'')\) and follows from \eqref{eq:peiffer} applied to \(m\) and \(v'm''\) since \(\delta\) is a bimodule map; and the identity \((v\cdot m')\cdot v''=v\cdot(m'\cdot v'')\), which is associativity in \(A\).  The unit is \((1,0,0,1)\).  The map \(\lambda_\delta\) is multiplicative because \(\delta(ma'+bm')=\delta(m)a'+b\delta(m')\); its kernel is \(\{(0,0,u,0):u\in U\}\).

\ref{re:i}.  Since \(\lambda_\delta(s)=E_{12}\) has the two-sided inverse \(E_{21}\) between the corner projectives of \(M_2(A)\), \(\lambda_\delta\) inverts \(\sigma\).  Let \(f:R_\delta\to D'\) invert \(\sigma\), put \(p'=f(p)\), \(e'=f(e)\), \(s'=f(s)\), and let \(t\in e'D'p'\) satisfy \(s't=p'\), \(ts'=e'\).  For \(m\in M\) one has \(ms=(0,0,0,\delta(m))\), so \(f(m)=f(m)p'=f(m)s't=f(\delta(m)_e)\,t\), where \(\gamma_e\) denotes \(\gamma\in A\) in the lower-right corner; thus \(f\) is determined on \(M\) by its restriction to the diagonal and by \(t\).  Also \(\gamma_ps=s\gamma_e\) for \(\gamma\in A\), so \(f(\gamma_p)s'=s'f(\gamma_e)\) and the two diagonal copies of \(A\) are identified through \(s'\) and \(t\).  The elements \(p',e',s',t\) are matrix units, and \(\gamma\mapsto f(\gamma_p)+f(\gamma_e)\), \(E_{21}\mapsto t\) defines a ring map \(M_2(A)\to D'\).  It extends \(f\): on \(M\) it sends \(\lambda_\delta(m)=\delta(m)E_{21}\) to \(f(\delta(m)_e)t=f(m)\).  Uniqueness holds because \(t\) is the unique inverse of the corner isomorphism.

\ref{re:ii}.  The right module \(pR_\delta\) is the row \((A,A)\) with action \((x,y)(a,v,m,b)=(xa+y\delta(m),\ xv+yb)\), and each row of \(M_2(A)\) restricted along \(\lambda_\delta\) has this action, so \(S_{R_\delta}\cong(pR_\delta)^{\oplus2}\).  The left module \(R_\delta e\) is the column \((A,A)^{\mathsf T}\) with action \((a,v,m,b)\binom xy=\binom{ax+vy}{\delta(m)x+by}\), and each column of \(S\) has this action, so \({}_{R_\delta}S\cong(R_\delta e)^{\oplus2}\).  A ring epimorphism with projective target on one side has \(\Tor^{R_\delta}_i(S,S)=0\) for \(i>0\); universal localisations are ring epimorphisms, so \(\lambda_\delta\) is stably flat.

\ref{re:iii}.  The endomorphism complex of \(T=[R_\delta p\xrightarrow{s}R_\delta e]\), written with coefficients of right multiplications, is
\[
 M\xrightarrow{\ m\mapsto(\delta(m),\delta(m))\ }A\oplus A\xrightarrow{\ (a,b)\mapsto a-b\ }A
\]
in degrees \(-1,0,1\): the degree \(-1\) term is \(\Hom_{R_\delta}(R_\delta e,R_\delta p)=eR_\delta p=M\), and its differential consists of the two composites \(x\mapsto xsm\) and \(x\mapsto xms\), which are right multiplications by \(sm=\delta(m)_p\) and \(ms=\delta(m)_e\); the degree \(0\) term is \(pR_\delta p\oplus eR_\delta e\), with differential the commutator with \(s\); the degree \(1\) term is \(pR_\delta e=A\).  Its cohomology is \(U\) in degree \(-1\), \(\{(a,a)\}/\{(\delta m,\delta m)\}\cong\Gam\) in degree \(0\), and \(0\) in degree \(1\).  Send \(a\in A\) to simultaneous right multiplication by \(a\) on \(R_\delta p\) and \(R_\delta e\), and \(m\in M\) to right multiplication by \(m\) as a map \(R_\delta e\to R_\delta p\).  This is compatible with the differentials, \(d(a,a)=0\) and \(d(m)=(\delta m,\delta m)\), and multiplicative in the opposite algebra: composing ``right \(a\), then right \(m\)'' is right multiplication by \(am\), so \(a\cdot^{\op}m=am\) and \(m\cdot^{\op}a=ma\) with the corner at \(e\) acting on the left and the corner at \(p\) on the right, which is why the opposite is required; products of two degree \(-1\) maps vanish because \(T\) has no degree \(-2\) endomorphisms.  The map induces the isomorphisms \(\Gam\to H^0\) and \(U\to H^{-1}\) just computed, so it is a quasi-isomorphism.  By stable flatness, \(\cK_{\lambda_\delta}=\Loc(T)\) and \(\Kc_{\lambda_\delta}=\thick(T)\) (\Cref{lem:compact-intersection}), and dg Morita theory gives \(F\).

\ref{re:iv}.  Restrict \(T\), with its right \(D\)-action through \ref{re:iii}, to the two idempotents.  The \(e\)-component is \(eT=[eR_\delta p\to eR_\delta e]=[M\xrightarrow{\delta}A]=D\) as a right dg \(D\)-module.  The \(p\)-component is \(pT=[A\xrightarrow{1}A]\), on which \(A\) acts by right multiplication and \(m\in M\) acts from degree \(0\) to degree \(-1\) by multiplication by \(\delta(m)\); the degree \(-1\) map \(h\) equal to the identity on \((pT)^0=A\) and zero on \((pT)^{-1}\) is right \(D\)-linear, since on degree \(-1\) elements both \(h(x\cdot m)\) and \(h(x)\cdot m\) vanish for degree reasons, and satisfies \(dh+hd=1\).  Thus \(pT\) is contractible as a right dg \(D\)-module.  For a semi-free \(Y\), \(FY=T\otimes_DY\) has \(e\)-component \(D\otimes_DY=Y\) and contractible \(p\)-component, so \(H^n(FY)\cong(0,H^nY)\) with \(R_\delta\) acting through \(R_\delta\to R_\delta/R_\delta pR_\delta=\Gam\).  This proves t-exactness and the profile assertion.  For \(Y=\Gam\) all identifications commute with the regular right \(\Gam\)-action, giving \(T\Ltensor_D\Gam\simeq h\Gam\).  Independently, \ref{re:ii} and \Cref{thm:kernel-dictionary} give full profile.  The dimension count is immediate.
\end{proof}

\begin{remark}\label{rem:realisation-benchmarks}
For \(\delta=0\), \eqref{eq:R-delta} is the Morita context ring with zero pairings of \cite{GreenPsaroudakis14,GaoPsaroudakis17}, and the theorem realises the formal dg algebra \(A\oplus M[1]\); for \(M=0\) it is the localisation \(T_2(A)\to M_2(A)\) of the upper triangular ring, with ordinary kernel \(A\); if \(\delta\) is an isomorphism the kernel is zero.  Matrix-ring universal localisations of exact contexts and their homologicality criterion are due to Chen--Xi \cite{ChenXi19}; the content here is the identification of the kernel of the cone with the prescribed two-term dg algebra, differential included.
\end{remark}

\subsection{Every thin dg algebra occurs}

Finite strict models for proper connective dg algebras are known over arbitrary fields by Raedschelders--Stevenson \cite[Appendix~A, Theorem~A.3]{Goodbody24}; see \Cref{rem:finite-strict-priority} below.  The next construction gives such a two-term model directly from a prescribed Hochschild class.

\begin{theorem}[Finite crossed models]\label{thm:finite-crossed}
Let \(\Gam\) be a finite-dimensional algebra, \(U\) a finite-dimensional \(\Gam\)-bimodule, and \(\mu\in HH^3(\Gam,U)\).  There is a finite-dimensional two-term dg algebra \(D=[M\to A]\) with \(H^0(D)=\Gam\), \(H^{-1}(D)=U\) and minimal-model class \([m_3]=\mu\).  Consequently every proper dg algebra \(B\) with \(H^i(B)=0\) for \(i\notin\{-1,0\}\) is quasi-isomorphic to a finite-dimensional two-term dg algebra, and hence, by \Cref{thm:realisation}, is the kernel dg algebra of a stably flat universal localisation \(R\to M_2(A)\) of finite-dimensional algebras at a single map, with target projective on both sides, t-exact kernel embedding, and full bounded profile.
\end{theorem}

\begin{proof}
Let \(\Gam^e=\Gam\otimes_k\Gam^{\op}\) and let \(I=\Hom_k(\Gam^e,U)\) with \((qf)(r)=f(rq)\); it is a finite-dimensional injective left \(\Gam^e\)-module, being coinduced from \(k\), and \(\nu:U\to I\), \(\nu(u)(r)=ru\), is an injective \(\Gam^e\)-linear map since \(\nu(u)(1)=u\).  Put \(J=I/U\) with quotient map \(q:I\to J\).  Since \(HH^n(\Gam,I)=\Ext^n_{\Gam^e}(\Gam,I)=0\) for \(n\ge1\), the long exact sequence gives an isomorphism \(\partial:HH^2(\Gam,J)\xrightarrow{\sim}HH^3(\Gam,U)\).  Choose a normalised Hochschild \(2\)-cocycle \(\alpha:\Gam^{\otimes2}\to J\) representing \(\partial^{-1}\mu\) and define
\[
 A=\Gam\oplus J,\qquad (x,j)(y,l)=(xy,\ xl+jy+\alpha(x,y)).
\]
The cocycle identity is the associativity of this product, normalisation makes \((1,0)\) the unit, and \(J\) is a square-zero ideal.  Let \(M=I\) with \(A\) acting through \(A\to\Gam\), and \(\delta(i)=(0,q(i))\).  Then \(\delta\) is a bimodule map, both sides of \eqref{eq:peiffer} vanish because \(J\) annihilates \(I\), and \Cref{lem:crossed} gives a finite-dimensional two-term dg algebra \(D\) with \(H^0(D)=A/(0\oplus J)=\Gam\) and \(H^{-1}(D)=\ker q=U\).

For the class, use the section \(s(x)=(x,0)\) and a linear section \(j:J\to I\) of \(q\) with \(M'=j(J)\); then \(s(x)s(y)-s(xy)=(0,\alpha(x,y))=\delta(j\alpha(x,y))\), so \(\beta=j\alpha\) in \eqref{eq:beta}, and \Cref{lem:two-term-minimal} gives
\[
 m_3(x,y,z)=x\,j\alpha(y,z)-j\alpha(xy,z)+j\alpha(x,yz)-j\alpha(x,y)\,z .
\]
This is the Hochschild differential of the lift \(j\alpha\) of \(\alpha\) to \(I\); it takes values in \(U\) and represents \(\partial[\alpha]=\mu\) by the definition of the connecting homomorphism.

For the consequence, let \(B\) be proper with cohomology in \(\{-1,0\}\); if \(B\) is acyclic it is quasi-isomorphic to \(0\), realised by the identity of \(k\) and its contractible cone.  Otherwise the transfer theorem \cite[Section~3.3]{Keller01Ainf} provides a strictly unital minimal \(A_\infty\)-model on \(H^*(B)=\Gam\oplus U[1]\), whose structure is \((m_2,m_3)\) by \Cref{prop:higher-mult}; let \(\mu\) be the class of \(m_3\).  The dg algebra \(D\) just constructed has a minimal model with the same \(\Gam\), \(U\) and class, hence \(A_\infty\)-isomorphic to that of \(B\) by \Cref{lem:gauge}.  Two dg algebras with \(A_\infty\)-isomorphic minimal models are connected by a zigzag of quasi-isomorphisms of dg algebras \cite[Section~3.3]{Keller01Ainf}, \cite[Section~1.3]{LefevreHasegawa03}.  Apply \Cref{thm:realisation} to \(D\).
\end{proof}

\begin{corollary}[All profiles for a thin kernel]\label{cor:thin-all-profiles}
For every proper dg algebra \(B\) with cohomology concentrated in degrees \(-1\) and \(0\), the marked bounded defects realised by stably flat universal localisations of finite-dimensional algebras with kernel dg algebra \(B\) are exactly
\[
 \{\cU_E/\per B:E\in\Dfd(B)\}.
\]
All of them are idempotent complete and coincide with the corresponding singular defects before completion.
\end{corollary}

\begin{proof}
Necessity is \Cref{cor:pair-invariant}.  \Cref{thm:finite-crossed} provides a full-profile realisation whose target is projective as a left module over the source, and \Cref{thm:ambient}\ref{amb:d} realises every \(\cU_E\).  The free module of a connective model is a silting compact generator with finite-dimensional endomorphism algebra, so \Cref{cor:finite-dimensional-precompletion,cor:negative-k-completeness} apply.
\end{proof}

\begin{example}[A non-formal kernel]\label{ex:nonformal}
Let \(A\) have basis \(1,t,a,b,c\), with unit \(1\) and the only nonzero products of radical basis elements
\[
 t^2=c,\qquad ta=b,\qquad at=c,
\]
so that \(\rad^3A=0\).  Let \(M\) have basis \(z,l,r,w\), with the only nonzero radical actions
\[
 tz=l,\qquad tr=w,\qquad zt=r,\qquad lt=w,
\]
and let \(\delta z=a\), \(\delta l=b\), \(\delta r=c\), \(\delta w=0\).  The bimodule laws hold, the only nontrivial mixed check being \((tz)t=w=t(zt)\); \(\delta\) is bimodule-linear, and \eqref{eq:peiffer} holds because the image of \(\delta\) annihilates \(M\).  Hence \(D=[M\to A]\) is a nine-dimensional two-term dg algebra with
\[
 \Gam=k[t]/(t^2),\qquad U=kw,\qquad tU=Ut=0 .
\]
Since \(t^2=c=\delta(r)\), \Cref{lem:two-term-minimal} with \(s(t)=t\) and \(\beta(t,t)=r\) gives
\[
 m_3(t,t,t)=t\,r-r\,t=w\ne0 .
\]
For any normalised \(2\)-cochain \(f:\Gam^{\otimes2}\to U\) the Hochschild coboundary satisfies \((d_{\mathrm{Hoch}}f)(t,t,t)=tf(t,t)-f(t^2,t)+f(t,t^2)-f(t,t)t=0\), so \(m_3\) is not a coboundary: \(D\) is not formal, and \(HH^3(\Gam,U)\cong k\) is generated by its class.  \Cref{thm:realisation} realises \(D\) as the kernel of a stably flat universal localisation \(R_\delta\to M_2(A)\) at a single map, with \(\dim_kR_\delta=19\), \(\dim_kM_2(A)=20\), target projective on both sides and full bounded profile.  Its kernel dg algebra therefore carries a nonzero \(m_3\): kernel dg algebras of finite universal localisations of finite-dimensional algebras need not be formal.  The dg algebra, the ring, the two projective decompositions, the strict dg identification and the value of \(m_3\), computed once on \(D\) and once on the presentation-cone algebra, are checked exactly by \texttt{verify\_nonformal.py} (\Cref{app:computation}).
\end{example}

\begin{remark}[Finite strictification and the crossed model]\label{rem:finite-strict-priority}
The existence of a finite two-term model also follows from \cite[Appendix~A, Theorem~A.3]{Goodbody24}.  Take a finite non-positive strict model \(C\), and quotient by the dg ideal equal to \(C^i\) in degrees \(i\le-2\), to \(d(C^{-2})\) in degree \(-1\), and to zero in degree zero.  Nonpositivity and the Leibniz rule show that this is a dg ideal.  If \(H^{<-1}(C)=0\), the projection is a quasi-isomorphism onto a finite two-term dg algebra.  Thus finite-model existence is prior work; the construction in \Cref{thm:finite-crossed} explicitly starts from \((\Gam,U,\mu)\) and computes the resulting minimal-model class, before applying the ring-epimorphism realisation of \Cref{thm:realisation}.
\end{remark}

\subsection{Optimality of the t-exact realisation}

The construction above has a converse when the marked kernel embedding is required to preserve the standard t-structures.  The converse applies to arbitrary ordinary homological ring epimorphisms, without a projectivity assumption on the target.

\begin{theorem}[T-exact kernel realisations are two-term]\label{thm:texact-optimality}
Let \(\lambda:R\to S\) be a homological ring epimorphism of ordinary rings and suppose
\[
 \cK_\lambda=\ker(S\Ltensor_R-)=\Loc(T),
 \qquad T\in K^b(R\text{-}\mathrm{proj}).
\]
If \(B=\RHom_R(T,T)^{\op}\) is connective and its Morita embedding
\(F:D(B)\to D(R)\) is t-exact for the standard t-structures, then \(T\) is homotopy equivalent to a complex of finitely generated projectives concentrated in degrees \([-1,0]\).  In particular \(H^i(B)=0\) for \(i<-1\).
\end{theorem}

\begin{proof}
A fully faithful t-exact functor reflects standard cohomological bounds.  Indeed, applying it to a truncation triangle gives a truncation triangle, whence
\(H^n(FY)\cong F_{\heartsuit}(H^nY)\); the fully faithful heart functor reflects zero objects.  This argument does not require \(Y\) to be bounded.

Homologicality gives the localisation triangle
\[
 K_R\longrightarrow R\xrightarrow{\lambda}S\longrightarrow K_R[1],
 \qquad K_R=\Cone(R\to S)[-1]\in\cK_\lambda.
\]
To check the last assertion in all degrees, apply \(S\Ltensor_R-\): the resulting map \(S\to S\Ltensor_RS\) is the inverse of the multiplication isomorphism.  With our cohomological convention, the displayed representative of \(K_R\) has terms in degrees \(0,1\).  Thus, writing \(K_R\cong FC\), reflection gives \(C\in D^{[0,1]}(B)\).  For every \(n\), adjunction gives
\[
 \Hom_{D(R)}(T,S[n])
 \cong\Hom_{D(S)}(S\Ltensor_RT,S[n])=0.
\]
Applying \(\RHom_R(T,-)\) to the triangle identifies the cohomology of \(T^\vee=\RHom_R(T,R)\) with that of \(\RHom_R(T,K_R)\).  Full faithfulness and \(F(B)=T\) identify the latter groups with \(\Hom_{D(B)}(B,C[n])=H^n(C)\).  Hence
\[
 T^\vee\in D^{[0,1]}(R^{\op}),
 \qquad T=F(B)\in D^{\le0}(R).
\]
Only cohomological bounds are asserted here; the two complexes have different additional module actions.

For a bounded projective complex whose cohomology vanishes above degree \(b\), the differential onto its highest term above \(b\) is a split surjection.  Removing the associated contractible summand and repeating eliminates all terms above \(b\).  Apply this first to \(T\), obtaining a projective representative with no positive terms.  Its finite-projective dual has no negative terms.  Apply the same cancellation to that dual above degree one; it is then supported in \([0,1]\).  Finite-projective biduality gives a representative of \(T\) in \([-1,0]\).  All steps preserve homotopy equivalence.  Its endomorphism complex has terms only in degrees \([-1,1]\), proving the vanishing below \(-1\).
\end{proof}

\begin{corollary}[Universality of a t-exact realisation]\label{cor:texact-universal}
Under \Cref{thm:texact-optimality}, write a two-term representative as \(T=[P^{-1}\xrightarrow{\sigma}P^0]\).  Then \(\lambda\) is the universal localisation at \(\sigma\).
\end{corollary}

\begin{proof}
The vanishing \(S\Ltensor_RT=0\) says that \(S\otimes_R\sigma\) is invertible.  The essential image of derived restriction is \(\cK_\lambda^\perp\).  A stalk \(R\)-module \(M\) belongs to this image exactly when \(\Hom_R(\sigma,M)\) is invertible.  Its preimage is a stalk \(S\)-module, because restriction is exact and reflects cohomological vanishing.  Its extending \(S\)-action is unique: full faithfulness applied to the identity of the underlying \(R\)-module identifies any two such actions.

Let \(f:R\to C\) invert \(\sigma\).  Tensor--Hom adjunction puts \({}_RC\) in \(\cK_\lambda^\perp\), giving its unique extending \(S\)-action.  Every right multiplication on \(C\) is \(R\)-linear and therefore \(S\)-linear.  Consequently \(s\cdot c=(s\cdot1)c\), and \(g(s)=s\cdot1\) is a unital ring homomorphism \(S\to C\) extending \(f\).  Any other extension defines the same \(S\)-action, hence equals \(g\).  This is the required universal property.
\end{proof}

\begin{corollary}[Classification of t-exact marked realisations]\label{cor:texact-classification}
For a proper connective dg algebra \(B\) over a field, the following are equivalent:
\begin{enumerate}[label=\textup{(\roman*)}]
\item \(B\) is quasi-isomorphic to the dg endomorphism algebra, with opposite, of a perfect generator of the kernel of a homological epimorphism of finite-dimensional algebras, with standard t-exact marked kernel embedding;
\item \(H^i(B)=0\) for \(i<-1\);
\item such a realisation exists as a universal localisation at one map, with target projective on both sides and full bounded profile.
\end{enumerate}
\end{corollary}

\begin{proof}
The implication (i)\(\Rightarrow\)(ii) is \Cref{thm:texact-optimality}; (ii)\(\Rightarrow\)(iii) is \Cref{thm:finite-crossed,thm:realisation}; and (iii)\(\Rightarrow\)(i) is immediate.  Quasi-isomorphisms of connective dg models preserve the standard t-structures.
\end{proof}

Thus allowing ordinary homological epimorphisms which are not assumed to be universal localisations does not enlarge this t-exact realisation range.  The full-profile seed of \Cref{thm:finite-crossed} still realises every cut by \Cref{thm:ambient}; an extension realising a proper cut need not retain the seed's projectivity or t-exactness.

\subsection{Projective targets and the Schur boundary}

\begin{corollary}[Kernel categories with projective right target]\label{cor:projective-target-classes}
The dg Morita classes of compact kernel categories of homological epimorphisms of finite-dimensional algebras with projective right target are exactly the classes admitting a proper connective dg model with cohomology in \([-1,0]\).  Every such class has a realisation with target projective on both sides, standard t-exact marking by a thin model, and full bounded profile.
\end{corollary}

\begin{proof}
Let \(S_R\) be projective.  By \Cref{prop:left-projection}, the kernel has the perfect generator
\[
 G=\Cone\!\left(\Hom_{R^{\op}}(S,R)\longrightarrow R\right),
\]
with terms in \([-1,0]\) and \(\Hom(G,G[n])\cong H^n(G)\).  Thus \(\RHom_R(G,G)^{\op}\) is proper with cohomology in \([-1,0]\), and models the compact kernel.  Conversely apply \Cref{thm:finite-crossed,thm:realisation} to a thin model.  A dg Morita equivalence preserves \(\Dfd\), by its intrinsic finite-total-Hom characterisation, so the full-profile assertion is independent of the compact-generator marking.  The t-exact assertion concerns the thin marking supplied by the construction.
\end{proof}

\begin{proposition}[Negative Calabi--Yau obstruction]\label{prop:negative-cy-obstruction}
Let \(B\) be a proper dg algebra such that \(\per B\ne0\) has Serre functor \([-d]\), where \(d\ge2\).  Then \(\per B\) is not equivalent to the compact kernel of either
\begin{enumerate}[label=\textup{(\roman*)}]
\item a homological epimorphism of finite-dimensional algebras with \(S_R\in\per(R^{\op})\); or
\item a finite stably flat universal localisation of finite-dimensional algebras, without any perfection assumption on the target.
\end{enumerate}
Neither exclusion requires t-exactness or full profile.
\end{proposition}

\begin{proof}
Every nonzero \(G\in\per B\) satisfies
\[
 \Hom(G,G[-d])\cong D_k\End(G)\ne0
\]
by Serre duality.  In (i), \Cref{prop:left-projection} supplies a nonzero perfect generator with \(\Hom(G,G[n])=0\) for \(n<-1\), a contradiction.  In (ii), the transported sum of the presentation cones has the same vanishing by \Cref{prop:thin}, again a contradiction.
\end{proof}

\begin{example}[Two excluded formal algebras]\label{ex:negative-cy-algebras}
The dg algebras
\[
 B_2=k[u]/(u^2),\quad |u|=-2,
 \qquad
 B_{\wedge}=\Lambda_k(u,v),\quad |u|=|v|=-1,
\]
with zero differential satisfy \Cref{prop:negative-cy-obstruction} with \(d=2\), over every field.  For \(B_2\), take the trace with value one on \(u\) and zero on \(1\); for \(B_{\wedge}\), take the coefficient of \(uv\).  The Gram matrices of \((a,b)\mapsto\operatorname{tr}(ab)\), in the respective bases \((1,u)\) and \((1,u,v,uv)\), are
\[
 \begin{pmatrix}0&1\\1&0\end{pmatrix},
 \qquad
 \begin{pmatrix}0&0&0&1\\0&0&1&0\\0&-1&0&0\\1&0&0&0\end{pmatrix}.
\]
These pairings are nondegenerate and graded symmetric, including in characteristic two, and give dg bimodule isomorphisms \(D_kB\cong B[-2]\).  Finite-projective tensor--Hom duality gives
\[
 D_k\RHom_B(X,Y)\simeq
 \RHom_B\bigl(Y,D_kB\Ltensor_BX\bigr)
 \qquad(X,Y\in\per B).
\]
The formula holds on finite free modules and extends by shifts, cones and direct summands; it identifies \([-2]\) as the Serre functor.
\end{example}

\begin{example}[A Schur realisation outside the epimorphism range]\label{ex:sphere-schur}
Let \(A\) be the radical-square-zero algebra of the oriented three-cycle
\[
 1\xrightarrow{a_1}2\xrightarrow{a_2}3\xrightarrow{a_3}1.
\]
Its basis is \(e_1,e_2,e_3,a_1,a_2,a_3\), with orthogonal vertex idempotents, the source and target identities of the arrows, and all products of two arrows zero.  Put \(e=e_2+e_3\), so \(eAe=k(2\to3)\), and take
\[
 T=[Ae_2\xrightarrow{a_2}Ae_3\xrightarrow{a_3}Ae_1]
 \quad\text{in degrees }[-2,-1,0].
\]
Then there is a recollement
\[
 D(B_2)\longrightarrow D(A)\xrightarrow{e(-)}D(eAe)
\]
whose kernel embedding is t-exact and has full bounded profile.  The source has dimension six and the quotient corner has dimension three.  The quotient functor here is the Schur functor, not extension along an ordinary ring epimorphism.

Indeed, the endomorphism complex of \(T\) has dimensions \(1,0,3,2\) in degrees \(-2,-1,0,1\), and zero in every other degree.  Its only nonzero differential sends the three diagonal scalars to their two consecutive differences and has rank two over every field.  Thus its cohomology is \(k\) in degrees \(-2,0\).  Sending \(u\) to right multiplication by \(a_1\), from \(Ae_1\) to \(Ae_2\), and \(1\) to the identity gives a strict quasi-isomorphism
\[
 B_2\longrightarrow\End_A^\bullet(T)^{\op}.
\]
As right dg \(B_2\)-modules, \(e_1T=B_2\); the other two vertex components are identity-contractible two-term complexes, on which \(u\) acts trivially.  Consequently \(T\Ltensor_{B_2}-\) preserves cohomological degrees and finite-dimensional total cohomology, and reflects the latter.

The object \(T\) belongs to \(\ker e(-)\).  If \(X\) belongs to this kernel, the finite filtration of \(\Hom_A^\bullet(T,X)\) by the three terms of \(T\) has two acyclic associated factors, \(e_2X\) and \(e_3X\), leaving \(e_1X\).  Hence \(\RHom_A(T,X)=0\) implies \(X=0\).  The \(\Loc(T)\)-approximation triangle inside \(D(A)\) now gives \(\ker e(-)=\Loc(T)\): its residual object belongs to the kernel and is right-orthogonal to \(T\), so is zero.  Dg Morita theory and the idempotent recollement of \cite[Proposition~6.1]{KalckYangII} yield the displayed recollement.  The complete finite computation is included in the supplement.  By \Cref{thm:texact-optimality,prop:negative-cy-obstruction}, this example cannot be replaced by any of the ordinary-epimorphism realisations excluded there.
\end{example}

\begin{example}[Full profile does not force a thin marking]\label{ex:marking-boundary}
Let \(A=k(1\to2)\), \(R=A\times k\), and \(\lambda:R\to k\) be projection.  The target is projective on both sides and multiplication \(k\otimes_Rk\to k\) is an isomorphism, so all positive Tor groups vanish and the kernel is \(D(A)\).  Put \(P_i=Ae_i\) and \(T=P_1\oplus P_2[2]\).  The complete Hom calculation is
\[
 \End_A(P_i)=k,\qquad
 \Hom_A(P_1,P_2)=k,\qquad\Hom_A(P_2,P_1)=0;
\]
all higher or lower Hom groups of these stalk projectives vanish.  Therefore \(B=\End_R^\bullet(T)^{\op}\) has zero differential,
\[
 H^0(B)=k\times k,\qquad H^{-2}(B)=k,
\]
and no other cohomology.  Since \(T\) and \(A\) classically generate the same perfect category, the finite-total-Hom characterisation gives full profile for this marking.  It is not t-exact by \Cref{thm:texact-optimality}.  Thus \Cref{cor:projective-target-classes} concerns existence of a thin model in a dg Morita class; it does not assert that every connective marked model is thin.  In particular, the broader \Cref{q:realisation} is not settled by the t-exact classification.
\end{example}


\subsection{Open problems}

The standard t-exact version of \Cref{q:realisation} is settled by \Cref{cor:texact-classification}, and \Cref{prop:negative-cy-obstruction} excludes a class of dg Morita types even without t-exactness.  The general full-profile question remains open, as does realisation by unrestricted ordinary homological epimorphisms after dropping t-exactness and allowing infinite right target projective dimension.  The Schur construction in \Cref{ex:sphere-schur} is a different realisation problem.

\begin{question}\label{q:profile-m3}
For \Cref{ex:nonformal}, do the posets of far-away orthogonals \(\{\cU_E/\per B\}\) of the non-formal dg algebra and of the formal dg algebra with the same cohomology differ?  The two have different Ext-algebras of the residue field, but eventual orthogonality is not determined by low-degree Ext.
\end{question}

\begin{question}\label{q:nonsilting}
The realisation theorem concerns kernels whose cone generator is silting.  Which dg algebras with cohomology in degrees \(-1,0,1\) occur as the endomorphism dg algebra of a presentation cone of a stably flat universal localisation of finite-dimensional algebras, as in the example \(k\times k\) after \Cref{prop:thin}?
\end{question}

\section{The canonical detector}\label{sec:detector}

\begin{theorem}[Canonical detector]\label{thm:detector}
Let \(R\) be a finite-dimensional \(k\)-algebra, let \(i:\cK\hookrightarrow D(R)\) be the inclusion of a compactly generated localising subcategory, and let \(\Gamma\) be its right adjoint.  Put \(I_{\cK}=\Gamma(DR)\in\cK\).  Then
\begin{equation}\label{eq:detector}
 \cK\cap\Db(R\text{-}\modcat)
 =\Bigl\{X\in\cK:\ \bigoplus_{n\in\bbZ}\Hom_{\cK}\bigl(X,I_{\cK}[n]\bigr)\text{ is finite-dimensional}\Bigr\}.
\end{equation}
If \(F:D(B)\xrightarrow{\sim}\cK\) is the dg Morita equivalence attached to a compact generator \(T\), then \(F^{-1}(I_{\cK})\simeq D_kT\), with the left \(B\)-action dual to the right \(B\)-action on \(T\); equivalently
\begin{equation}\label{eq:tensor-duality}
 \RHom_B(Y,D_kT)\simeq D_k\bigl(T\Ltensor_BY\bigr)\qquad(Y\in D(B)).
\end{equation}
\end{theorem}

\begin{proof}
For every \(X\in D(R)\), possibly unbounded,
\begin{equation}\label{eq:ambient-duality}
 \RHom_R(X,DR)\simeq D_kX,\qquad \Hom_{D(R)}(X,DR[n])\cong D_kH^{-n}(X).
\end{equation}
Indeed, on modules evaluation at \(1\) identifies \(\Hom_R(X,DR)\) with \(D_kX\); the inverse sends a functional \(f\) to \(x\mapsto(r\mapsto f(rx))\), and the left \(R\)-action fixed in \Cref{sec:preliminaries} makes this \(R\)-linear.  The identification holds on complexes with their Hom differentials.  Since \(D_k\) is exact, \(\Hom_R(-,DR)\) sends acyclic complexes to acyclic complexes, so \(DR\) is K-injective and the complex identification computes derived Hom.  Taking cohomology gives \eqref{eq:ambient-duality}.

For \(X\in\cK\), adjunction gives \(\Hom_{\cK}(X,\Gamma DR[n])\cong\Hom_{D(R)}(iX,DR[n])\).  By \eqref{eq:ambient-duality} the sum of these spaces is finite-dimensional exactly when only finitely many cohomology spaces of \(X\) are nonzero and each is finite-dimensional, which over a finite-dimensional algebra is membership in \(\Db(R\text{-}\modcat)\).  This proves \eqref{eq:detector}.

Since \(F\) is fully faithful with right adjoint \(\Theta=\RHom_R(T,-)\), the composite \(F\Theta\) is the coreflection into \(\cK\).  Hence
\[
 \Theta(\Gamma DR)\simeq\Theta(DR)=\RHom_R(T,DR)\simeq D_kT,
\]
the last isomorphism being \eqref{eq:ambient-duality} for \(X=T\); precomposition on the Hom complex is carried to the dual of the right \(B\)-action on \(T\).  Formula \eqref{eq:tensor-duality} is tensor--Hom adjunction over the field \(k\), computed on a K-projective replacement of \(Y\).
\end{proof}

\begin{corollary}[The bounded defect as an invariant of a pair]\label{cor:pair-invariant}
Let \(\lambda:R\to S\) be a stably flat universal localisation of finite-dimensional algebras at a finite set \(\Sigma\), with \(T=\bigoplus_\sigma\Cone(\sigma)\) and \(B\) as in \Cref{sec:preliminaries}.  Then \(D_kT\in\Dfd(B)\),
\begin{equation}\label{eq:profile-as-cut}
 \Prof(F)=\cU_{D_kT},
 \qquad
 \DeltaB(\lambda)\simeq\cU_{D_kT}/\per B,
\end{equation}
and \(\DeltaB(\lambda)\) is determined by the \(k\)-linear pair \((\cK_\lambda,\Gamma_\lambda DR)\): a \(k\)-linear triangle equivalence of such pairs commuting with coproducts induces an equivalence of bounded defects.
\end{corollary}

\begin{proof}
The complex \(T\) is a bounded complex of finite-dimensional projectives, so \(D_kT\) has finite-dimensional total cohomology.  For \(Y\in D(B)\) with \(FY\) bounded, \(Y\simeq\Theta FY\) has finite-dimensional total cohomology because \(T\) is a bounded complex of finitely generated projectives; hence \(\Prof(F)\subseteq\Dfd(B)\).  On \(\Dfd(B)\), \eqref{eq:detector} and \eqref{eq:tensor-duality} give \(\Prof(F)=\cU_{D_kT}\).  The Morita equivalence sends \(\per B\) onto \(\thick(T)=\Kc_\lambda\) by the Thomason--Neeman identification of compacts used in \Cref{lem:compact-intersection}, so \Cref{cor:relative-dg-bounded} gives the quotient formula.  A \(k\)-linear coproduct-preserving equivalence of pairs preserves compact objects and the finite-total-Hom condition in \eqref{eq:detector}, hence the numerator and denominator of the quotient.
\end{proof}

\begin{remark}[Reconstruction of the intrinsic subcategory]\label{rem:intrinsic}
Put
\[
 \cF(\cK)=\Bigl\{X\in\cK:\ \bigoplus_n\Hom(C,X[n])\text{ is finite-dimensional for all }C\in\cK^c\Bigr\}.
\]
Under Morita equivalence \(\bigoplus_n\Hom(T,FY[n])\cong\bigoplus_nH^n(Y)\), and for fixed \(X\) the compacts satisfying the condition form a thick subcategory, so testing \(T\) tests \(\thick(T)=\cK^c\).  This recovers \(F(\Dfd B)=\cF(\cK)\) and its generator independence \Cref{prop:finite-total-hom}, which is the Gao--Psaroudakis criterion \cite[Lemma~4.1]{GaoPsaroudakis18} in Morita coordinates.  Full bounded compatibility in the sense of \Cref{def:r-bounded} is the equality of the intrinsic category \(\cF(\cK)\) with the cut \eqref{eq:detector}; both are independent of the generator, which recovers \Cref{cor:rbounded-generator-independent} without reference to the silting graph.  The two objects \(D_kT\) and \(D_kB\) must not be conflated: the first records the embedding, the second is the intrinsic dual of the free module, and \(\cU_{D_kB}=\Dfd(B)\).
\end{remark}

\begin{remark}
The detector need not itself have bounded ambient realisation.  On the right-module side of the AHKLY localisation \Cref{prop:dual-dga}, \(T\simeq k\oplus k[1]\) as a left module over \(B=k[\eps]/(\eps^2)\), so \(D_kT\simeq k\oplus k[-1]\) is not perfect and \(F(D_kT)\) is unbounded.  Brown--Comenetz duals as detectors and far-away orthogonality already appear in \cite[Section~2]{KostasPsaroudakisVitoria25}; no novelty is claimed for the duality \eqref{eq:tensor-duality}.
\end{remark}

We record the truncation lemma that converts a complex detector into a module detector.

\begin{lemma}[Module detectors]\label{lem:module-detector}
Let \(R\) be finite-dimensional and \(X\in\Db(R\text{-}\modcat)\).  There are a finite left \(R\)-module \(N\), an integer \(m\), and a bounded complex \(J\) of finite-dimensional injective left \(R\)-modules with a triangle
\begin{equation}\label{eq:truncation-triangle}
 N[-m]\longrightarrow X\longrightarrow J\longrightarrow N[-m+1].
\end{equation}
Consequently, for every \(Z\in\Db(R\text{-}\modcat)\),
\begin{equation}\label{eq:module-detector}
 \RHom_R(Z,X)\in\Db(k\text{-}\modcat)\iff\RHom_R(Z,N)\in\Db(k\text{-}\modcat)\iff\Ext^{\gg0}_R(Z,N)=0 .
\end{equation}
\end{lemma}

\begin{proof}
Injective envelopes of finite modules over a finite-dimensional algebra are finite, so \(X\) has a bounded-below injective resolution \(I\) with finite-dimensional terms.  Choose \(m\) above the highest nonzero cohomology degree of \(X\).  The good truncation \(\tau_{\le m}I\) is quasi-isomorphic to \(I\) and has top term \(N=\ker(I^m\to I^{m+1})\) in degree \(m\); the short exact sequence of complexes \(0\to N[-m]\to\tau_{\le m}I\to\sigma_{\le m-1}I\to0\), with \(\sigma_{\le m-1}I\) the brutal truncation, gives \eqref{eq:truncation-triangle} with \(J=\sigma_{\le m-1}I\), a bounded complex of finite-dimensional injectives.  For bounded finite-dimensional \(Z\), \(\RHom_R(Z,J)\) is computed by the ordinary Hom complex from a bounded finite-dimensional representative of \(Z\) into \(J\), hence is bounded and finite-dimensional.  Applying \(\RHom_R(Z,-)\) to \eqref{eq:truncation-triangle} gives the first equivalence in \eqref{eq:module-detector}.  For the second, each \(\Ext^n_R(Z,N)\) is finite-dimensional and vanishes for \(n\) sufficiently negative because \(Z\) is bounded and \(N\) is a module.
\end{proof}

\section{Ambient one-point extensions}\label{sec:ambient}

\subsection{The hereditary theorem}
Let \(\lambda:R\to S=R_\Sigma\) be a stably flat universal localisation of finite-dimensional algebras at a finite set \(\Sigma\) of maps between finitely generated projective left modules, with \(T=\bigoplus_\sigma\Cone(\sigma)\), \(B\), \(F\), \(\Theta\) as in \Cref{sec:preliminaries}, and bounded profile \(\Prof(\lambda)=\{Y\in D(B):FY\in\Db(R\text{-}\modcat)\}\subseteq\Dfd(B)\).  Let \(W\) be a finite-dimensional right \(R\)-module and form the one-point extension
\[
 A'=R[W]=\begin{pmatrix}k&W\\0&R\end{pmatrix},\qquad p=e_{11},\qquad e=e_{22}.
\]
Left \(A'\)-modules are triples \((V,Z,\varphi)\) with \(\varphi:W\otimes_RZ\to V\); the corner functor \(j^*=e(-)\) has derived left adjoint \(j_!Z=A'e\Ltensor_RZ=(W\Ltensor_RZ,Z,\id)\), which is fully faithful, coproduct-preserving and compact-preserving, and \(i_*:D(k)\to D(A')\) is restriction along \(A'\to A'/A'eA'=k\).  Put
\[
 \Sigma'=\{A'e\otimes_R\sigma:\sigma\in\Sigma\}=j_!\Sigma .
\]

\begin{theorem}[Ambient one-point extension]\label{thm:ambient}
With the notation above:
\begin{enumerate}[label=\textup{(\alph*)}]
\item\label{amb:a} The universal localisation of \(A'\) at \(\Sigma'\) is
\begin{equation}\label{eq:ambient-target}
 S'=A'_{\Sigma'}\cong\begin{pmatrix}k&W\otimes_RS\\0&S\end{pmatrix},
\end{equation}
with \(\lambda':A'\to S'\) given by \(\lambda\) on \(R\), \(w\mapsto w\otimes1\) on \(W\), and the identity on \(k\).  It is stably flat if and only if
\begin{equation}\label{eq:tor-condition}
 \Tor_i^R(W,{}_RS)=0\qquad(i\ge1),
\end{equation}
which holds for every \(W\) when \({}_RS\) is projective.
\item\label{amb:b} Under \eqref{eq:tor-condition}, \(\cK_{\lambda'}=j_!(\cK_\lambda)\), \(\Kc_{\lambda'}=\thick(j_!T)\), and \(\RHom_R(T,T)\to\RHom_{A'}(j_!T,j_!T)\) is a quasi-isomorphism of dg algebras; the kernel dg algebra of \(\lambda'\), marked by \(j_!T\), is \(B\).
\item\label{amb:c} Under \eqref{eq:tor-condition}, in the marking of \ref{amb:b},
\begin{equation}\label{eq:profile-intersection}
 \Prof(\lambda')=\Prof(\lambda)\cap\cU_{\Theta(D_kW)}=\{Y\in\Prof(\lambda):W\Ltensor_RFY\in\Db(k\text{-}\modcat)\},
\end{equation}
so that \(\DeltaB(\lambda')\simeq\bigl(\Prof(\lambda)\cap\cU_{\Theta(D_kW)}\bigr)/\per B\).  At the level of detectors, \(D_k(j_!T)\simeq D_kT\oplus\Theta(D_kW)\) in \(D(B)\).
\item\label{amb:d} If \(\Prof(\lambda)=\Dfd(B)\) and \({}_RS\) is projective, then for every \(E\in\Dfd(B)\) there is a finite right \(R\)-module \(W\) with \(\Prof(\lambda')=\cU_E\).
\end{enumerate}
\end{theorem}

\begin{proof}
\ref{amb:a}, \emph{universal property.}  The displayed map inverts \(\Sigma'\), since extension along it restricted to the corner is extension along \(\lambda\).  Let \(f:A'\to D\) be a ring map inverting \(\Sigma'\), and put \(p_D=f(p)\), \(e_D=f(e)\).  Applying the exact functor \(e_D(-)\) to the inverted maps shows that the corner map \(R\to e_DDe_D\) inverts \(\Sigma\), so it extends uniquely to \(g:S\to e_DDe_D\).  The restriction of \(f\) to \(W\subseteq p_DDe_D\) extends to \(W\otimes_RS\to p_DDe_D\), \(w\otimes s\mapsto f(w)g(s)\), which is well defined because \(f(wr)=f(w)f(r)\) and \(g\) extends \(f|_R\).  Together with \(f\) on the corner \(k\), these define a ring map \(S'\to D\) extending \(f\); balancedness is the only nontrivial multiplicativity check, and uniqueness holds because \(k\), \(W\) and \(R\) generate \(A'\).  This proves \eqref{eq:ambient-target} independently of \eqref{eq:tor-condition}.

\ref{amb:a}, \emph{stable flatness.}  Let \(L=S\Ltensor_R-\) be the Bousfield localisation of \(D(R)\) with kernel \(\cK_\lambda\), so that \(LR=S\) as an \(R\)-module.  Since \(j_!\) preserves coproducts, cones and summands and is fully faithful, the localising subcategory generated by the new cones is
\begin{equation}\label{eq:new-kernel}
 \cC'=\Loc(\Cone\Sigma')=j_!(\cK_\lambda),
\end{equation}
which is generated by compact objects; let \(L'\) be the associated smashing localisation of \(D(A')\).  For \(Z\in D(R)\) apply \(j_!\) to the localisation triangle \(K_Z\to Z\to LZ\): the first term lies in \(\cC'\), and \(j_!LZ\) is \(\cC'\)-local because \(\Hom_{D(A')}(j_!K,j_!LZ[n])\cong\Hom_{D(R)}(K,LZ[n])=0\) for \(K\in\cK_\lambda\), by \(j^*j_!\cong\id\).  Hence
\begin{equation}\label{eq:L-prime-j}
 L'(j_!Z)\simeq j_!(LZ).
\end{equation}
Similarly \(i_*k=(k,0,0)\) is \(\cC'\)-local, since \(\Hom(j_!K,i_*k[n])\cong\Hom(K,j^*i_*k[n])=0\), so \(L'(i_*k)=i_*k\).  As \(A'=A'p\oplus A'e=i_*k\oplus j_!R\),
\begin{equation}\label{eq:L-prime-A}
 L'(A')\simeq i_*k\oplus j_!(S)=\bigl(k\oplus(W\Ltensor_RS),\ S,\ \ldots\bigr),
\end{equation}
whose cohomology is \(H^0(L'A')={}_{A'}S'\) and \(H^{-i}(L'A')=i_*\Tor_i^R(W,S)\) for \(i\ge1\), with \(H^{>0}=0\).

A smashing localisation of the derived category of an algebra \(A\) is induced by a homological ring epimorphism \(A\to A''\) with \(A''\) an ordinary algebra if and only if the localised regular object \(LA\) has cohomology concentrated in degree zero; in that case \(A''=\End_{D(A)}(LA)^{\op}\) in the left-module convention.  Indeed, Nicolás--Saorín \cite[Theorem~5]{NicolasSaorin09} model the smashing localisation by a homological epimorphism to a dg algebra \(\mathcal C\) whose restricted regular module is \(LA\); if \(H^{i}(LA)=0\) for \(i\ne0\), the maps of dg algebras \(\mathcal C\leftarrow\tau_{\le0}\mathcal C\to H^0(\mathcal C)\) under \(A\) are quasi-isomorphisms, replacing the dg target by an ordinary algebra with the same fully faithful restriction; conversely the restriction of the regular module of an ordinary homological target is a stalk complex.  The same reduction appears in the commutative setting in \cite[proof of Theorem~4.9]{AMSTV20}.

Under \eqref{eq:tor-condition}, \eqref{eq:L-prime-A} is concentrated in degree \(0\).  Write \(I=i_*k\) and \(J=j_!S\).  Then \(\End(I)=k\), \(\End(J)=\End_{D(R)}(S)=S^{\op}\) by homologicality of \(\lambda\), \(\Hom(J,I)=\Hom_{D(R)}(S,j^*i_*k)=0\), and \(\Hom(I,J)=pJ=W\otimes_RS\) because \(I=A'p\) is projective.  Hence \(\End_{D(A')}(L'A')^{\op}\cong S'\), and applying \(L'\) to right multiplications on \(A'\) recovers the map \(\lambda'\) of \eqref{eq:ambient-target}.  Thus \(\lambda'\) is a homological ring epimorphism, and being a universal localisation it is stably flat.  Conversely, if \(\lambda'\) were stably flat while some \(\Tor_i^R(W,S)\ne0\), Neeman--Ranicki \cite{NeemanRanicki04} would identify its kernel with \eqref{eq:new-kernel} and its localised regular object with the stalk module \(S'\), contradicting \eqref{eq:L-prime-A}.

\ref{amb:b}.  By stable flatness and \Cref{lem:compact-intersection}, \(\cK_{\lambda'}=\cC'\) and \(\Kc_{\lambda'}=\thick\{j_!\Cone\sigma\}=\thick(j_!T)\).  On K-projective representatives, derived adjunction and \(j^*j_!\cong\id\) give a quasi-isomorphism of dg Hom complexes \(\RHom_R(T,T)\to\RHom_{A'}(j_!T,j_!T)\) compatible with composition; taking graded opposites identifies the kernel dg algebras.

\ref{amb:c}.  A complex of \(A'\)-modules has bounded finite-dimensional cohomology iff both corner complexes do, since the two exact idempotent functors \(p(-)\), \(e(-)\) have underlying vector spaces summing to that of the complex.  By the formula for \(j_!\), \(j_!FY\) is bounded iff \(FY\) and \(W\Ltensor_RFY\) are, giving the second description in \eqref{eq:profile-intersection}.  Tensor duality and adjunction give, with all sides fixed as in \Cref{sec:preliminaries},
\[
 D_k\bigl(W\Ltensor_RFY\bigr)\simeq\RHom_R(FY,D_kW)\simeq\RHom_B\bigl(Y,\Theta(D_kW)\bigr),
\]
and \(\Theta(D_kW)\) is proper because \(T\) is a bounded complex of finitely generated projectives.  This proves \eqref{eq:profile-intersection}; the quotient formula is \Cref{cor:relative-dg-bounded} for \(\lambda'\), and the detector formula follows from the two components of \(j_!T\) and \eqref{eq:tensor-duality}.

\ref{amb:d}.  Let \(E\in\Dfd(B)\) and \(X=FE\), which lies in \(\Db(R\text{-}\modcat)\) by fullness.  \Cref{lem:module-detector} gives a module \(N\) with \(\RHom_R(FY,X)\) bounded iff \(\RHom_R(FY,N)\) bounded, for all \(Y\in\Dfd(B)\).  Put \(W=D_kN\), so \(D_kW=N\); condition \eqref{eq:tor-condition} holds because \({}_RS\) is projective.  Full faithfulness of \(F\) gives \(\RHom_B(Y,E)\simeq\RHom_R(FY,X)\), hence \(\cU_E=\cU_{\Theta(D_kW)}\), and \ref{amb:c} with \(\Prof(\lambda)=\Dfd(B)\) gives \(\Prof(\lambda')=\cU_E\).
\end{proof}

\begin{remark}
Formula \eqref{eq:profile-intersection} is an intersection: ambient one-point extensions can only shrink a profile, and \(W=0\) preserves it.  The projectivity hypothesis in \ref{amb:d} is on \({}_RS\), the side paired with \(W\) in \eqref{eq:tor-condition}; projectivity of \(S_R\) would not suffice.  When \eqref{eq:tor-condition} fails, \(\lambda'\) is a universal localisation of finite-dimensional algebras that is not stably flat, a source of such examples.
\end{remark}

\subsection{Marks--Nakayama kernels}
Let \(k\) be algebraically closed, \(d\ge3\), \(A_d=k\widetilde Q_{2d}/\mathfrak r^d\), and let \(J\) be Marks-admissible as in (\Cref{prop:marks-projectivity,thm:nakayama-intrinsic}).  Following \Cref{sec:nakayama} we work with right \(A_d\)-modules, that is, with \(R=A_d^{\op}\); the kernel dg algebra is \(B_J\), the profile is full, and \(S_J\) is projective on both sides.

\begin{corollary}[Profiles with Nakayama kernel]\label{cor:nakayama-profiles}
For every Marks-admissible \(J\), the marked bounded defects realised by stably flat universal localisations of finite-dimensional algebras with kernel dg algebra \(B_J\) are exactly
\[
 \{\cU_E/\per B_J:E\in\Dfd(B_J)\}.
\]
\end{corollary}

\begin{proof}
Necessity is \Cref{cor:pair-invariant}.  Sufficiency is \Cref{thm:ambient}\ref{amb:d}, whose hypotheses are supplied by (\Cref{prop:marks-projectivity,thm:nakayama-intrinsic}).
\end{proof}

We compute these defects explicitly in small rank.  Write \(M(i,l)=e_iA_d/\rad^l(e_iA_d)\), \(1\le l<d\), for the uniserial right module with composition factors \(S_i,\dots,S_{i+l-1}\), indices modulo \(2d\).

\begin{lemma}[Stable Hom in \(\underline{\modcat}\,A_d\)]\label{lem:stable-hom}
\begin{enumerate}[label=\textup{(\roman*)}]
\item \(\Omega M(i,l)=M(i+l,d-l)\), \(M(i,l)[1]=M(i+l-d,d-l)\), and \(\Omega^4M(i,l)\cong M(i,l)\).
\item For \(X=M(i,l)\), \(Y=M(j,m)\) and \(r\equiv i-j\pmod{2d}\) with \(0\le r\le 2d-1\),
\[
 \dim_k\underline{\Hom}(X,Y)=\begin{cases}1,&r<m\le r+l<d,\\0,&\text{otherwise.}\end{cases}
\]
\item For \(X,Y\in\Db(A_d^{\op}\text{-}\modcat)\), \(\Hom_{D}(X,Y[n])\cong\Hom_{\Dsg}(qX,qY[n])\) for \(n\gg0\).  Hence, for a finite module \(N\),
\begin{equation}\label{eq:profile-stable}
 \Ext^{\gg0}(X,N)=0\iff\underline{\Hom}(qX,qN[t])=0\ \text{for all }t\in\bbZ .
\end{equation}
\end{enumerate}
\end{lemma}

\begin{proof}
(i) The kernel of the projective cover of \(M(i,l)\) is the uniserial module starting at \(i+l\) of length \(d-l\); solving for the inverse syzygy gives the shift formula; two syzygies add \(d\) to the top index and four add \(2d\).

(ii) A map \(X\to Y\) is determined by the image of the top of \(X\), which can only lie in the layer of \(Y\) indexed \(r\); such a layer exists iff \(r<m\), and the image is annihilated by paths of length \(l\) iff \(m\le r+l\).  The map factors through the projective cover \(P_j\to Y\) exactly when its image lifts to layer \(r\) of \(P_j\) with the same annihilation relation, that is, iff \(r+l\ge d\); every map through a projective factors through \(P_j\).  Removing these maps gives the formula, over every field.

(iii) Replace \(X\) and \(Y\), modulo bounded projective complexes, by shifted finite modules; since \(A_d\) is self-injective these complexes are also bounded injective, and Hom into or out of them vanishes in large positive degree.  For modules, \(\Ext^n(M,N)\cong\underline\Hom(\Omega^nM,N)=\Hom_{\Dsg}(qM,qN[n])\) for \(n\ge1\).  Periodicity (i) gives \eqref{eq:profile-stable}.
\end{proof}

\begin{proposition}[Small lattices]\label{prop:m3}
Let \(\cQ_J=\Dfd(B_J)/\per B_J\simeq\thick_{\underline{\modcat}A_d}\{S_i:i\in J\cup(J+d)\}\) \textup{(\Cref{prop:marks-projectivity,thm:nakayama-intrinsic})}.
\begin{enumerate}[label=\textup{(\roman*)}]
\item\label{m3:single} If \(J=\{a\}\) is a singleton (all admissible \(J\) for \(d=3\)), \(\cQ_J\) is the split four-periodic category on the shifts of \(S_a\), its thick subcategories are \(0\) and \(\cQ_J\), and both are realised: by \(W=0\) and by \(W=D_kS_a\).
\item\label{m3:opposite} If \(d=4\) and \(J=\{a,a+2\}\), then \(\cQ_J=O_a\oplus O_{a+2}\) with \(O_i\) the split four-periodic category on the shifts of \(S_i\); its thick subcategories form the four-element Boolean lattice and all are realised, the subcategory omitting \(O_i\) by \(W=D_kS_i\) and zero by \(W=D_k(S_a\oplus S_{a+2})\).
\item\label{m3:adjacent} If \(d=4\) and \(J=\{a,a+1\}\), put \(A_0=S_a\), \(A_1=S_{a+1}\), \(A_2=M(a,2)\) and \(O_r=\add\{A_r[t]:t\in\bbZ\}\).  The uniserial extension gives a non-split triangle
\begin{equation}\label{eq:m3-triangle}
 A_1\longrightarrow A_2\longrightarrow A_0\longrightarrow A_1[1],
\end{equation}
\(\cQ_J=\add(O_0,O_1,O_2)\), and the thick subcategories of \(\cQ_J\) are exactly
\[
 0,\quad O_0,\quad O_1,\quad O_2,\quad\cQ_J,
\]
forming the lattice \(M_3\).  All five are realised: \(\cQ_J\) by \(W=0\), \(O_0\) by \(W=D_kM(a,2)\), \(O_1\) by \(W=D_kS_a\), \(O_2\) by \(W=D_kS_{a+1}\), and \(0\) by \(W=D_k(S_a\oplus S_{a+1})\).
\end{enumerate}
In each case the poset of realisable bounded defects with kernel \(B_J\) is the full lattice of thick subcategories of \(\cQ_J\).
\end{proposition}

\begin{proof}
By \Cref{cor:nakayama-profiles} and \eqref{eq:profile-stable}, the realisable defects are the images of the subcategories \(\{X:\underline\Hom(qX,qN[t])=0\ \forall t\}\) for finite modules \(N\), and every detector \(E\) may be replaced by such an \(N\) through \Cref{lem:module-detector}.  In all cases every indecomposable has endomorphism ring \(k\) in the stable category and each orbit has four objects, by \Cref{lem:stable-hom}(i),(ii).

\ref{m3:single}: all off-diagonal stable Homs between the four shifts of \(S_a\) vanish by \Cref{lem:stable-hom}(ii) (the case \(d=3\), \(a=1\) is \Cref{lem:a63-stable}); every morphism between finite sums reduces to identities and zeros, cones stay in the additive category, and a nonzero thick subcategory contains an indecomposable and hence, by shift closure, all four.  \(W=0\) realises \(\cQ_J\); \(N=S_a\) excludes every indecomposable since \(\underline\Hom(S_a[t],S_a[t])\ne0\), realising \(0\).

\ref{m3:opposite}: \Cref{lem:stable-hom}(ii) gives \(\underline\Hom(O_a,O_{a+2}[t])=0=\underline\Hom(O_{a+2},O_a[t])\) for all \(t\), so \(\cQ_J\) is the orthogonal sum of two split categories; thick subcategories are sums of a subset of the two orbits.  For \(N\) the sum of the simples of the omitted orbits, \eqref{eq:profile-stable} picks out exactly the remaining orbits.

\ref{m3:adjacent}: \Cref{lem:stable-hom}(ii) gives the stable Hom table below, rows indexed by the source, columns by the target, entries listing dimensions for target shifts \(0,1,2,3\):
\[
\begin{array}{c|ccc}
 & A_0 & A_1 & A_2\\ \hline
 A_0 & (1,0,0,0) & (0,1,0,0) & (0,0,0,0)\\
 A_1 & (0,0,0,0) & (1,0,0,0) & (1,0,0,0)\\
 A_2 & (1,0,0,0) & (0,0,0,0) & (1,0,0,0)
\end{array}
\]
Thus among the twelve indecomposables the only non-isomorphisms are scalar multiples of the arrows
\(A_1[t]\to A_2[t]\to A_0[t]\to A_1[t+1]\), the rotations of \eqref{eq:m3-triangle}, and consecutive compositions vanish because they are consecutive maps in a triangle; the radical of the additive category \(\add(O_0,O_1,O_2)\) has square zero.  Given a morphism between finite sums, split off identities by block row and column operations until all entries are radical; then for each source type the only nonzero block maps onto its unique successor, and independent base changes reduce it to rank normal form.  Every morphism is therefore a direct sum of identities, zeros, and copies of the arrows, whose cones are the third terms of rotations of \eqref{eq:m3-triangle}.  Hence \(\add(O_0,O_1,O_2)\) is thick, contains \(S_a,S_{a+1}\), and equals \(\cQ_J\).  A nonzero thick subcategory contains a complete orbit; any two orbits force the third by \eqref{eq:m3-triangle}; each single orbit is split and thick by the table.  This is the lattice \(M_3\).  The realising modules follow from the columns of the table via \eqref{eq:profile-stable}: \(N=S_a\) excludes \(O_0,O_2\) and keeps \(O_1\); \(N=S_{a+1}\) keeps \(O_2\); \(N=M(a,2)\) keeps \(O_0\); \(N=S_a\oplus S_{a+1}\) excludes everything.  The stable Hom tables for \(d=3,4\) are checked exactly by \texttt{verify\_nakayama\_profiles.py} (\Cref{app:computation}).
\end{proof}

\begin{corollary}[No support space for adjacent pairs]\label{cor:non-distributive}
For \(d=4\) and \(J=\{a,a+1\}\) the poset of realisable bounded defects with kernel \(B_J\) is the non-distributive lattice \(M_3\).  It is therefore neither order-isomorphic nor order-anti-isomorphic to the lattice of closed subsets of any topological space.  In particular no analogue of the support classification of \Cref{thm:support-classification} exists for this kernel.
\end{corollary}

\begin{proof}
Closed subsets of a space form a distributive lattice under union and intersection; distributivity is preserved by order isomorphisms and anti-isomorphisms; \(M_3\) is not distributive.
\end{proof}

\subsection{A genuinely dg boundary}

\begin{theorem}[The twenty-dimensional boundary]\label{thm:dg-boundary}
Let \(k\) be algebraically closed, \(A=A_3=k\widetilde Q_6/\mathfrak r^3\) with right modules as in \Cref{subsec:a63}, \(R=A^{\op}\), and \(\lambda\) the Marks localisation at \(\{S_1,S_4\}\) with kernel dg algebra
\[
 B^0=ke_1\oplus ke_4,\qquad B^{-1}=ka\oplus kb,\qquad a=e_4ae_1,\quad b=e_1be_4,\quad ab=ba=0,\quad d=0,
\]
and simple dg modules \(L_1,L_4\) \textup{(\Cref{prop:marks-projectivity,thm:nakayama-intrinsic})}.  Let \(W=D_kS_1\) and \(A'=R[W]\), an algebra of dimension \(1+1+18=20\), and let \(\lambda'\) be the induced localisation.  Then \(\lambda'\) is stably flat with kernel dg algebra \(B\),
\begin{equation}\label{eq:ext-degrees}
 \Ext^q_B(L_1,L_1)=\begin{cases}k,&q\equiv0\pmod4,\ q\ge0,\\0,&\text{otherwise},\end{cases}
 \qquad
 \Ext^q_B(L_4,L_1)=\begin{cases}k,&q\equiv2\pmod4,\ q\ge0,\\0,&\text{otherwise},\end{cases}
\end{equation}
\(\Theta(S_1)=L_1\), \(\cU_{L_1}=\per B\), and
\begin{equation}\label{eq:dg-boundary}
 \DeltaB(\lambda')=0,\qquad \Dfd(B)/\per B\simeq\vect_k^{\bbZ/4}\ne0 .
\end{equation}
Thus the intrinsic connective formula \(\DeltaB\simeq\Dfd(B)/\per B\) fails for a marked kernel dg algebra with nonzero negative cohomology.
\end{theorem}

\begin{proof}
The target of \(\lambda\) is projective on both sides and the profile is full (\Cref{prop:marks-projectivity,thm:nakayama-intrinsic}), so \Cref{thm:ambient} applies with the marked kernel preserved.  The identifications \(F(L_1)=S_1\), \(F(L_4)=S_4\) are \Cref{prop:nakayama-cancellation,prop:a63-bounded-kernel}; full faithfulness gives \(\Theta(S_1)=L_1\), so the new detector is \(\Theta(D_kW)=L_1\).

The short exact sequences \(0\to L_4[1]\to Be_1\to L_1\to0\) and \(0\to L_1[1]\to Be_4\to L_4\to0\) of \Cref{prop:a63-dga} splice to a semi-free resolution of \(L_1\) whose underlying graded module is
\[
 P^{\#}=Be_1\oplus Be_4[2]\oplus Be_1[4]\oplus Be_4[6]\oplus\cdots,
\]
with free generators \(x_n\) of degree \(-2n\) at vertex \(1\) for \(n\) even and at vertex \(4\) for \(n\) odd, and differential \(dx_0=0\), \(dx_n=ax_{n-1}\) for \(n\) odd, \(dx_n=bx_{n-1}\) for \(n\) even.  Each connecting shift is by \(2\), one from the internal degree of the radical generator and one from the cone.  Since \(ab=ba=0\), \(d^2=0\), the cycles are \(x_0\) and the elements \(ax_n,bx_n\), all of which except \(x_0\) are boundaries, so the augmentation to \(L_1\) is a quasi-isomorphism.  The radical acts by zero on \(L_1\), so \(\Hom_B(P,L_1)\) has zero differential, and \(\Hom_B(Be_v[2n],L_1)=e_vL_1\) in degree \(2n\).  Generators at vertex \(1\) occur for \(n=2m\), giving degrees \(4m\); this proves the first formula in \eqref{eq:ext-degrees}.  The resolution of \(L_4\) is obtained by interchanging the vertices and arrows, and its vertex-\(1\) generators occur for \(n=2m+1\), giving degrees \(4m+2\).

For \(\cU_{L_1}=\per B\), let \(Y\in\Dfd(B)\).  Full faithfulness gives \(\RHom_B(Y,L_1)\simeq\RHom_R(FY,S_1)\), and by \Cref{lem:stable-hom}(iii) its cohomology in large degrees is \(\underline\Hom(qFY,qS_1[n])\).  By \Cref{thm:a63-comparison}, \(q\Kb(\lambda)\) is the split four-periodic category on \(qS_1\) with \(\dim\Hom(Y_r,Y_s[n])=1\) iff \(s+n-r\equiv0\pmod4\); every nonzero object has a nonzero map to \(qS_1[n]\) for some \(n\), recurring in all degrees \(n+4m\).  Hence \(\RHom_B(Y,L_1)\) is bounded iff \(qFY=0\) iff \(FY\in\per R\) iff \(Y\in\per B\).  \Cref{thm:ambient}\ref{amb:c} gives \(\Prof(\lambda')=\per B\) and \eqref{eq:dg-boundary}; the intrinsic quotient is the displayed category by \Cref{thm:a63-comparison}.  Both formulas in \eqref{eq:ext-degrees} and the four-periodic Hom pattern are checked exactly by \texttt{verify\_dg\_boundary.py} (\Cref{app:computation}).
\end{proof}

\begin{remark}
The zero defect is reached without the kernel algebra being ordinary: this is the phenomenon of \Cref{thm:dual-numbers} for a kernel dg algebra with \(H^{-1}\ne0\).  The marked dg algebra \(B\) is not quasi-isomorphic to an ordinary algebra; no assertion is made about other objects of its dg Morita class.  As in \Cref{sec:support}, the failure lies inside the finite-dimensional silting range of \Cref{cor:finite-dimensional-precompletion}, so \(\DeltaB(\lambda')=\Defsg(\lambda')=0\) and no summand obstruction is involved.
\end{remark}

\section{Ordinary kernels: normal form and two-sidedness}\label{sec:ordinary}

Throughout this section \(C\) is a finite-dimensional \(k\)-algebra, viewed as a dg algebra concentrated in degree \(0\); then \(\Dfd(C)=\Db(C\text{-}\modcat)\) and \(\Dsg(C)\) is the usual singularity category.  For a finite left \(C\)-module \(N\) put
\begin{equation}\label{eq:UN}
 \cU_N=\{Y\in\Db(C\text{-}\modcat):\Ext^{\gg0}_C(Y,N)=0\},
\end{equation}
where the bound may depend on \(Y\).  By \Cref{lem:module-detector} this agrees with \Cref{def:profile} for \(E=N\).

\subsection{One-point extensions}
For a finite right \(C\)-module \(M\), the \emph{one-point extension} of Ringel \cite{Ringel84} is
\[
 C[M]=\begin{pmatrix}k&M\\0&C\end{pmatrix},\qquad p=e_{11},\qquad e=e_{22}.
\]
A left \(C[M]\)-module is a triple \((V,Y,\varphi)\) with \(V\) a \(k\)-space, \(Y\) a left \(C\)-module and \(\varphi:M\otimes_CY\to V\); a right module is a triple \((V,Y',\psi)\) with \(Y'\) a right \(C\)-module and \(\psi:V\otimes_kM\to Y'\).  The corner functor \(j^*=e(-)\) has derived left adjoint \(j_!=C[M]e\Ltensor_C-\) and right adjoint \(j_*=\RHom_C(eC[M],-)\); explicitly
\begin{equation}\label{eq:ope-functors}
 j_!Y=(M\Ltensor_CY,\ Y,\ \id),\qquad j_*Y=(0,Y,0).
\end{equation}

\begin{lemma}\label{lem:ope-stratifying}
The ideal \(C[M]eC[M]\) is stratifying.  The quotient map \(\lambda_M:C[M]\to k\) is a stably flat universal localisation at \(0\to C[M]e\), its left kernel is \(j_!(D(C))\) with compact part \(\thick(C[M]e)\simeq\per C\), and \(D(k)\to D(C[M])\to D(C)\) is a stratifying recollement.
\end{lemma}

\begin{proof}
As a left \(C\)-module \(eC[M]=(0,C)\cong C\) is free, so multiplication \(C[M]e\otimes_CeC[M]\to C[M]eC[M]\) is an isomorphism onto the sum of the two right-hand corners and the higher Tor terms vanish.  The remaining assertions are \Cref{thm:stratifying-delta}.
\end{proof}

\begin{theorem}[Normal form for ordinary kernels]\label{thm:normal-form}
Fix a finite-dimensional algebra \(C\).  Among coproduct-preserving, compact-preserving dg Morita embeddings \(F:D(C)\hookrightarrow D(R)\) with \(R\) finite-dimensional, in particular among the kernels of stably flat universal localisations of finite-dimensional algebras at finite \(\Sigma\) whose kernel dg algebra is \(C\), the bounded profiles are exactly the subcategories \(\cU_N\), \(N\in C\text{-}\modcat\).  Every \(\cU_N\) is the profile of the localisation \(\lambda_{D_kN}:C[D_kN]\to k\), whose compact kernel is marked by \(\per C\), and then
\begin{equation}\label{eq:normal-form-defect}
 \DeltaB(\lambda_{D_kN})\simeq\cU_N/\per C .
\end{equation}
Two modules give the same profile exactly when their far-away left orthogonals \eqref{eq:UN} agree.
\end{theorem}

\begin{proof}
Let \(T=F(C)\).  Compact preservation makes \(T\) a bounded complex of finite-dimensional projective \(R\)-modules, so \Cref{cor:pair-invariant} applies to the embedding: the profile is \(\cU_E\) with \(E=D_kT\in\Db(C\text{-}\modcat)\).  \Cref{lem:module-detector} over \(C\) produces a module \(N\) with \(\cU_E=\cU_N\).

Conversely let \(N\) be a finite left \(C\)-module and \(M=D_kN\).  By \Cref{lem:ope-stratifying} and \eqref{eq:ope-functors}, \(Y\in D(C)\) has bounded realisation \(j_!Y\) exactly when both \(Y\) and \(M\Ltensor_CY\) have finite-dimensional total cohomology.  Tensor duality gives
\[
 D_k(M\Ltensor_CY)\simeq\RHom_C(Y,D_kM)\simeq\RHom_C(Y,N),
\]
so the profile is \(\cU_N\), and \Cref{thm:stratifying-delta} gives \eqref{eq:normal-form-defect}.  The final assertion is a restatement.
\end{proof}

\begin{remark}
The theorem is a normal form for the profile, not for the embedding, the recollement, or the ambient algebra; it does not identify different modules with the same orthogonal.  No Gorenstein, commutativity, or finite-global-dimension hypothesis on \(C\) is used.  The passage from a complex detector to a module uses injective truncation; the projective direction would not give a left-orthogonality statement.  In the coordinates of \Cref{thm:detector} the canonical detector of \(\lambda_{D_kN}\) is \(D_k(C[M]e)\cong N\oplus D_kC\); the summand \(D_kC\) imposes no condition on bounded \(Y\) since \(\RHom_C(Y,D_kC)=D_kY\), which is why the reduced detector \(N\) suffices.
\end{remark}

\subsection{Two-sidedness}

\begin{theorem}[Two-sidedness of one-point extensions]\label{thm:two-sided}
Let \(N\) be a finite left \(C\)-module, \(M=D_kN\), \(A_N=C[M]\) and \(\lambda_N:A_N\to k\).  Then:
\begin{enumerate}[label=\textup{(\roman*)}]
\item on the left-module side, \(\Prof=\cU_N\) and \(\DeltaB^{\mathrm{left}}(\lambda_N)\simeq\cU_N/\per C\);
\item on the right-module side, the profile is all of \(\Db(C^{\op}\text{-}\modcat)\) and \(\DeltaB^{\mathrm{right}}(\lambda_N)\simeq\Dsg(C^{\op})\);
\item \(\cU_N=\Db(C\text{-}\modcat)\iff\pd_{C^{\op}}M<\infty\iff\pd_{A_N^{\op}}k<\infty\), that is, exactly when the bounded-extension hypothesis \textup{(B2)} of \Cref{sec:preliminaries} holds for \(\lambda_N\) on the left; the corresponding right-module hypothesis always holds, since \({}_{A_N}k=A_Np\) is projective.
\end{enumerate}
\end{theorem}

\begin{proof}
(i) is \Cref{thm:normal-form}.  For (ii), the right kernel functor is \(Y'\mapsto Y'\Ltensor_CeA_N=(0,Y')\), underived because \({}_C(eA_N)\) is free; it preserves bounded finite-dimensional complexes and sends \(C_C\) to the projective right module \(eA_N\).  The right-module version of \Cref{thm:stratifying-delta} gives (ii).

For (iii), if \(\pd(M_C)<\infty\) then \(M\Ltensor_C-\) preserves bounded finite-dimensional complexes, so \(\cU_N\) is everything.  Conversely, if the profile is full then \(M\Ltensor_C(C/\rad C)\) is bounded; choose a minimal projective resolution \(P_\bullet\to M\) by finite right \(C\)-modules.  Its differentials have entries in the radical, so tensoring with \(C/\rad C\) kills them and \(\Tor_n^C(M,C/\rad C)=P_n/P_n\rad C\).  Eventual vanishing and Nakayama's lemma force \(P_n=0\) for \(n\gg0\).  For the second equivalence use the exact sequence of right \(A_N\)-modules
\[
 0\longrightarrow(0,M)\longrightarrow pA_N\longrightarrow k\longrightarrow0 .
\]
The minimal right \(C\)-projective resolution of \(M\) inflates termwise to the minimal right \(A_N\)-projective resolution of \((0,M)\): its terms are sums of the old-corner projective \(eA_N\), its differentials are unchanged, and no new-corner projective can occur because the top of \((0,M)\) has no \(p\)-constituent.  Hence \(\pd_{A_N^{\op}}(0,M)=\pd_{C^{\op}}M\), and the sequence gives the equivalence with \(\pd_{A_N^{\op}}k<\infty\).
\end{proof}

\begin{example}[The AHKLY algebra]\label{ex:ahkly}
For \(C=E=k[\eps]/(\eps^2)\) and \(N=k\), the algebra \(A_N=E[k]\) is the AHKLY algebra \(kQ/(\beta^2,\alpha\beta)\) of \Cref{sec:preliminaries}; it appears as a triangular algebra in Chen's Example~4.3(2) \cite{ChenSchur09}.  Over the self-injective local algebra \(E\), \(\Ext^{\gg0}(Y,k)=0\) forces \(Y\) to have finite projective dimension, so \(\cU_k=\per E\): the left defect of the corner quotient \(E[k]\to k\) is zero, while its right defect is \(\Dsg(E)\simeq\vect_k\).  In the recollement of \Cref{lem:ope-stratifying}, the left-module kernel of the corner quotient is \(\ker(i^*)=j_!(D(E))\).  The arrow-inversion map \(E[k]\to M_2(k)\) of \Cref{thm:ahkly-left} is a different localisation of the same algebra: since \(M_2(k)\cong(pE[k])^{\oplus2}\) as a right module, its left-module kernel is \(\{X:pX=0\}=\ker(i^!)=j_*(D(E))\), whose profile is full by \eqref{eq:ope-functors}, in agreement with the nonzero defect found there.  The zero defect of \Cref{thm:dual-numbers} is the same phenomenon as (i) here, seen on the opposite module side of that arrow inversion.  Thus the two boundaries of \Cref{sec:periodic-ladder} are the two recollement embeddings \(j_!\) and \(j_*\) of \(D(E)\) into \(D(E[k])\).
\end{example}

\section{Support classification and non-generation}\label{sec:support}

Throughout this section \(k\) is algebraically closed and \(C\) is a finite-dimensional self-injective \(k\)-algebra, so that \(\Dsg(C)\simeq\underline{\modcat}\,C\) is the stable module category, by Buchweitz's theorem \cite{Buchweitz} in the form used in \Cref{sec:nakayama}.  We write \(q:\Db(C\text{-}\modcat)\to\Dsg(C)\) for the quotient.

\subsection{Support theories}
\begin{definition}[Support theory]\label{def:support-theory}
A \emph{support theory} for \(C\) consists of a projective \(k\)-variety \(\PP_C\) and closed subsets \(\VV_C(M)\subseteq\PP_C\) for finite \(C\)-modules \(M\), such that:
\begin{enumerate}[label=\textup{(S\arabic*)}]
\item\label{S:basic} \(\VV_C(M)=\varnothing\) iff \(M\) is projective; \(\VV_C(\Omega M)=\VV_C(M)\); \(\VV_C(M\oplus M')=\VV_C(M)\cup\VV_C(M')\); and \(\VV_C(M_h)\subseteq\VV_C(M_i)\cup\VV_C(M_j)\) for every short exact sequence \(0\to M_1\to M_2\to M_3\to0\) and \(\{h,i,j\}=\{1,2,3\}\);
\item\label{S:pair} for finite modules \(M,N\): \(\Ext^{\gg0}_C(M,N)=0\) iff \(\VV_C(M)\cap\VV_C(N)=\varnothing\);
\item\label{S:realise} every closed subset of \(\PP_C\) is \(\VV_C(N)\) for some finite module \(N\);
\item\label{S:full} \(\VV_C(C/\rad C)=\PP_C\).
\end{enumerate}
\end{definition}

By \ref{S:basic}, \(\VV_C\) is constant on isomorphism classes of the stable category and descends to \(\Dsg(C)\); for \(X\in\Db(C\text{-}\modcat)\) we put \(\VV_C(X)=\VV_C(M)\) for any finite module \(M\) with \(qM\cong qX[t]\) for some \(t\).  Then \(\VV_C(X)=\varnothing\) iff \(X\in\per C\), support is invariant under shifts, decreases under summands, and satisfies the union property for triangles.  For a closed subset \(Z\subseteq\PP_C\) put
\begin{equation}\label{eq:DZ}
 \sD_Z=\{X\in\Dsg(C):\VV_C(X)\cap Z=\varnothing\}.
\end{equation}
By \ref{S:basic} this is a thick subcategory; as a thick subcategory of the Krull--Schmidt category \(\underline{\modcat}\,C\) it is idempotent complete.

\begin{theorem}[Support classification]\label{thm:support-classification}
Let \(C\) admit a support theory.  The marked subcategories of \(\Dsg(C)\) occurring as bounded defects of stably flat universal localisations of finite-dimensional algebras at finite \(\Sigma\) with kernel dg algebra \(C\) are exactly the \(\sD_Z\), \(Z\subseteq\PP_C\) closed.  Distinct closed subsets give distinct subcategories.  Every \(\sD_Z\) is realised by \(\lambda_{D_kN}:C[D_kN]\to k\) for any finite module \(N\) with \(\VV_C(N)=Z\); in that family
\begin{equation}\label{eq:fixed-data}
 \cK_{\lambda}\simeq D(C),\qquad \Kc_\lambda\simeq\per C,\qquad C[D_kN]/C[D_kN]eC[D_kN]=k,\qquad \Dsg(C[D_kN])\simeq\Dsg(C),
\end{equation}
and every \(\DeltaB(\lambda_{D_kN})\) is idempotent complete and coincides with \(\Defsg(\lambda_{D_kN})\) before completion.
\end{theorem}

\begin{proof}
By \Cref{thm:normal-form} every profile is \(\cU_N\) for a finite module \(N\), and every \(\cU_N\) occurs, realised by \(\lambda_{D_kN}\).  For \(Y\in\Db(C\text{-}\modcat)\) choose a finite module \(M\) and \(t\) with \(qM\cong qY[t]\); the difference is measured by bounded complexes of projectives, which are also bounded complexes of injectives because \(C\) is self-injective, so \(\Ext^{\gg0}(Y,N)=0\) iff \(\Ext^{\gg0}(M,N)=0\).  By \ref{S:pair},
\begin{equation}\label{eq:UN-support}
 \cU_N=\{Y\in\Db(C\text{-}\modcat):\VV_C(Y)\cap\VV_C(N)=\varnothing\}.
\end{equation}
The functor \(\cU_N/\per C\to\Dsg(C)\) is fully faithful by \Cref{lem:subcategory-quotient-embedding}, since \(\per C\subseteq\cU_N\) and \(\cU_N\) is triangulated; its essential image consists of the objects of \(\Dsg(C)\) represented by modules in \(\cU_N\), which by \eqref{eq:UN-support} is \(\sD_{\VV_C(N)}\).  With \ref{S:realise} every closed \(Z\) occurs.

If \(Z_1\ne Z_2\), one of the locally closed sets \(Z_1\setminus Z_2\), \(Z_2\setminus Z_1\) is nonempty and, \(\PP_C\) being a variety over an algebraically closed field, contains a closed point \(p\).  By \ref{S:realise} there is a module \(X_p\) with \(\VV_C(X_p)=\{p\}\); its image lies in exactly one of \(\sD_{Z_1},\sD_{Z_2}\).

For \eqref{eq:fixed-data}: the kernel statements are \Cref{lem:ope-stratifying}; the singularity equivalence is Chen's Theorem~4.1(1) \cite{ChenSchur09}, applicable because the upper-left ring \(k\) has global dimension zero and \(D_kN\) is finite over \(k\), with no hypothesis on \((D_kN)_C\); its Schur functor is the exact corner restriction \(e(-)\), which sends \(j_!Y\) to \(Y\) and hence respects the marking.  Finally the compact kernel has the silting object \(C\) with artinian endomorphism algebra, so \Cref{thm:unconditional-comparison,cor:finite-dimensional-precompletion} gives idempotent completeness and the equality with the singular defect; alternatively \(\sD_Z\) is idempotent complete as a thick subcategory of a Krull--Schmidt category.
\end{proof}

\begin{theorem}[Non-generation]\label{thm:no-generator}
Let \(C\) admit a support theory and let \(Z\subseteq\PP_C\) be closed.  Suppose some irreducible component \(Y\) of \(\PP_C\) meets both \(Z\) and \(\PP_C\setminus Z\).  Then \(\sD_Z\) is a nonzero proper idempotent-complete triangulated category with no classical generator.  Consequently \(\sD_Z\) is not triangle equivalent to \(\Dsg(\Lambda)\) for any finite-dimensional algebra \(\Lambda\), nor to \(\Dfd(B')/\per B'\) for any proper cohomologically non-positive dg algebra \(B'\).
\end{theorem}

\begin{proof}
For \(G\in\Dsg(C)\), the objects \(X\) with \(\VV_C(X)\subseteq\VV_C(G)\) form a thick subcategory containing \(G\), by shift invariance, the union property, and monotonicity on summands; hence
\begin{equation}\label{eq:generator-support}
 X\in\thick(G)\implies\VV_C(X)\subseteq\VV_C(G).
\end{equation}
The set \(Y\cap(\PP_C\setminus Z)\) is a nonempty open subset of the irreducible variety \(Y\), hence dense in \(Y\), and its closed points are dense in it.  For each such closed point \(p\), \ref{S:realise} gives \(X_p\) with \(\VV_C(X_p)=\{p\}\); it lies in \(\sD_Z\), so \(\sD_Z\ne0\).  If \(G\) classically generated \(\sD_Z\), \eqref{eq:generator-support} would force every such \(p\) into the closed set \(\VV_C(G)\), hence \(Y\subseteq\VV_C(G)\); but \(Y\cap Z\ne\varnothing\) contradicts \(G\in\sD_Z\).  Properness: by \ref{S:full}, \(C/\rad C\notin\sD_Z\) because \(Z\ne\varnothing\).

For a finite-dimensional algebra \(\Lambda\), the finitely many simple modules classically generate \(\Db(\Lambda\text{-}\modcat)\) and hence \(\Dsg(\Lambda)\).  For proper non-positive \(B'\), the standard truncation triangles express every object of \(\Dfd(B')\) through finitely many finite-dimensional \(H^0(B')\)-modules \Cref{lem:d-bfg-heart}, which are built from the finitely many simple \(H^0(B')\)-modules; properness places \(\per B'\) in \(\Dfd(B')\), and the image of the sum of the simples generates the quotient.  A category with a classical generator cannot be equivalent to \(\sD_Z\).
\end{proof}

\begin{remark}
When \(Z=\varnothing\) the category \(\sD_Z=\Dsg(C)\) has a generator; when \(Z=\PP_C\) it is zero; the hypothesis excludes both.  The obstruction is independent of the summand obstruction \(\operatorname{SObs}_0\) of \Cref{subsec:summand-negative-k}, which vanishes throughout the one-point-extension family by \Cref{thm:support-classification}.  The statement excludes one-object proper connective models only; nonproper or multi-object dg descriptions are not excluded.
\end{remark}

\subsection{Complete intersections}

\begin{theorem}[Complete intersections]\label{thm:ci-classification}
Let \(C\) be a finite-dimensional commutative local complete intersection over \(k\) of codimension \(c\ge1\), and let \(\VV_C\) be the projectivised cohomological support of Avramov--Buchweitz \cite{AvramovBuchweitz00} in \(\PP_C=\PP^{c-1}_k\).  Then \(\VV_C\) is a support theory in the sense of \Cref{def:support-theory}.  Consequently the marked bounded defects with kernel \(\per C\) are exactly the \(\sD_Z\), \(Z\subseteq\PP^{c-1}_k\) closed, and for \(c\ge2\) and \(\varnothing\ne Z\subsetneq\PP^{c-1}_k\) the category \(\sD_Z\) has no classical generator.
\end{theorem}

\begin{proof}
An artinian local complete intersection is Gorenstein, hence self-injective.  Property \ref{S:basic} and \ref{S:full} are the basic properties of support varieties over complete intersections, \cite[Theorem~5.6]{AvramovBuchweitz00}, in the formulation \cite[Theorem~2.2]{BerghJorgensenEasy}.  Property \ref{S:pair} is the pair-support criterion: the support of the pair equals the intersection of the supports \cite[Theorem~2.2(3)]{BerghJorgensenEasy}, and it is trivial exactly when \(\Ext^{\gg0}_C(M,N)=0\) \cite[Theorem~5.6]{AvramovBuchweitz00}.  Property \ref{S:realise} is the realisability of every cone by a module of finite length \cite[Theorem~4.2 and the following Remark]{BerghJorgensenEasy}.  Since \(\PP^{c-1}_k\) is irreducible, the component hypothesis of \Cref{thm:no-generator} reduces to \(\varnothing\ne Z\subsetneq\PP^{c-1}_k\), which requires \(c\ge2\).
\end{proof}

\begin{proposition}[Hyperplanes and matrix factorisations]\label{prop:hmf}
Let \(C=Q/(f_1,\dots,f_c)\) be a minimal presentation of a complete intersection as above and let \(W\subseteq\PP^{c-1}_k\) be the hyperplane spanned by the classes of \(f_1,\dots,f_{c-1}\).  Then, by Bergh--Jorgensen \cite[Theorem~3.5]{BerghJorgensenMF},
\begin{equation}\label{eq:hmf}
 \sD_W\simeq\HMF\bigl(Q/(f_1,\dots,f_{c-1}),\,f_c\bigr),
\end{equation}
the homotopy category of matrix factorisations of \(f_c\) over the codimension-\((c-1)\) complete intersection \(Q/(f_1,\dots,f_{c-1})\), the functor being the cokernel functor of that theorem.  Hence every such matrix-factorisation category is the bounded defect of a one-point extension of \(C\).
\end{proposition}

\begin{proof}
Theorem~3.5 of \cite{BerghJorgensenMF} identifies \(\HMF(Q/(f_1,\dots,f_{c-1}),f_c)\) with the full subcategory of \(\Dsg(C)\) of objects whose affine support meets the hyperplane only in \(0\), which is \(\sD_W\).  \Cref{thm:ci-classification} realises \(\sD_W\).
\end{proof}

\begin{example}[A seven-dimensional algebra]\label{ex:seven}
Let \(C=k[x,y]/(x^2,y^2)\), a local complete intersection of codimension \(2\) with \(k\)-basis \(1,x,y,xy\).  Let \(M=C/(x)\), regarded as a right \(C\)-module, so that \(N=D_kM\cong C/(x)\) is the detector, and let \(X=C/(y)\) be a test object.  Put \(A=C[M]\); then \(\dim_kA=1+2+4=7\).

Since \(\ker(y\colon C\to C)=yC\) and \(\ker(x\colon C\to C)=xC\), the modules \(X\) and \(M\) have one-periodic minimal resolutions with all differentials \(y\), respectively \(x\).  Tensoring the resolution of \(X\) with \(M=k[y]/(y^2)\) gives the repeated map \(y\) on \(M\), whose kernel equals its image, so
\[
 \Tor_0^C(M,X)=k,\qquad\Tor^C_i(M,X)=0\quad(i>0);
\]
tensoring the resolution of \(M\) with \(M\) kills every differential, so \(\Tor_i^C(M,M)=M\) for all \(i\ge0\).  Hence \(X\in\cU_N\) and \(M\notin\cU_N\).  The bounded \(A\)-module realising \(X\) is \(j_!X=(k,X,\varphi)\) with \(\varphi:M\otimes_CX\xrightarrow{\sim}k\).  Its endomorphisms in the defect are the stable endomorphisms of \(X\): \(\End_C(X)=k[x]/(x^2)\), and every map \(X\to C\) has image in \(\ann_C(y)=yC\), which every map \(C\to X\) kills, so no nonzero endomorphism factors through a projective and
\[
 \End_{\DeltaB(\lambda_M)}(j_!X)\cong\underline{\End}_C(X)\cong k[x]/(x^2).
\]
The support of \(N\) is the single point of \(\PP^1_k\) represented by the relation \(x^2\): over the hypersurface \(Q/(ax^2+by^2)\) with \(b\ne0\), \(N\) is the quotient by a non-zerodivisor and has projective dimension one, while for \(b=0\) it has infinite projective dimension.  By \Cref{prop:hmf} with \(W=\{[x^2]\}\),
\[
 \DeltaB(\lambda_M)\simeq\HMF\bigl(k[[x,y]]/(x^2),\,y^2\bigr),
\]
and the rank-one factorisation \((y,y)\) corresponds to \(X\).  By \Cref{thm:no-generator} this matrix-factorisation category has no classical generator.  Since \(M\) is a nonprojective module over the self-injective algebra \(C\), \(\pd(M_C)=\infty\), so \(\pd(k_A)=\infty\) by \Cref{thm:two-sided}: the bounded-extension hypothesis (B2) fails on the left and the defect is not the kernel of a functor induced by derived extension.  The calculations are checked exactly over \(\mathbb Q\) by \texttt{check\_7dim.py} (\Cref{app:computation}).
\end{example}

\begin{example}[Distinct marked cuts related by a coordinate symmetry]\label{ex:support-symmetry}
Let \(C=k[x,y]/(x^2,y^2)\), \(M=C/(x)\), and \(N=C/(y)\). The algebra automorphism \(\tau\) exchanging \(x\) and \(y\) gives an exact pullback autoequivalence \(\tau^*\) of \(D(C)\), preserving perfect objects and exchanging \(M\) and \(N\). Since it preserves derived Hom, it induces
\[
 \cU_M/\per C\xrightarrow{\sim}\cU_N/\per C.
\]
These marked subcategories are distinct. Indeed, the free resolution of \(M\) has every differential multiplication by \(x\): exactness follows from \(\ker(x:C\to C)=xC\). Applying \(\Hom_C(-,M)\) gives zero differentials, so \(\Ext_C^n(M,M)\cong M\ne0\) for every \(n\ge0\). Applying \(\Hom_C(-,N)\), with \(N\cong k[x]/(x^2)\), gives multiplication by \(x\), whose kernel and image are both \(kx\). Thus \(\Ext_C^n(M,N)=0\) for every \(n>0\). Consequently \(M\in\cU_N\) but \(M\notin\cU_M\).

The two detector supports are the distinct points \([x^2]\) and \([y^2]\), exchanged by \(\tau\). Thus different closed subsets can yield abstractly equivalent avoidance categories, even though the marked classification distinguishes them. The remaining problem is to classify abstract equivalences, taking these symmetries into account. The identities and the single repeating resolution step are checked in \texttt{verify\_coordinate\_swap.py}; exactness of the full periodic resolution is the argument above.
\end{example}

\subsection{Finite \texorpdfstring{\(p\)}{p}-groups}

\begin{theorem}[\(p\)-groups]\label{thm:pgroup}
Let \(k\) have characteristic \(p>0\), let \(P\) be a finite \(p\)-group and \(C=kP\).  Let \(\PP_C=\Proj H^*(P,k)\) and let \(\VV_C(M)\subseteq\PP_C\) be the projectivised cohomological support variety of Carlson \cite{Carlson83}.  Then \(\VV_C\) is a support theory for \(kP\).  Consequently the marked bounded defects with kernel \(\per kP\) are exactly the \(\sD_Z\), \(Z\subseteq\Proj H^*(P,k)\) closed, and \(\sD_Z\) has no classical generator whenever some irreducible component of \(\Proj H^*(P,k)\) meets both \(Z\) and its complement.
\end{theorem}

\begin{proof}
The algebra \(kP\) is local and self-injective, with unique simple module \(k\).  The cohomology ring \(H^*(P,k)\) is a finitely generated graded-commutative \(k\)-algebra by Evens--Venkov, so \(\PP_C\) is a projective variety; its irreducible components correspond to conjugacy classes of maximal elementary abelian subgroups by Quillen \cite{Quillen71}.  Property \ref{S:basic} is standard: \(\VV_C(M)=\varnothing\) iff \(M\) is projective, support is invariant under \(\Omega\), additive on sums, and subadditive on short exact sequences \cite{Carlson83}, \cite[Chapter~5]{BensonII}.  Property \ref{S:full} holds because \(\VV_C(k)\) is the whole variety.  Property \ref{S:realise} is Carlson's realisability theorem: every closed homogeneous subvariety of the affine cone is the support of a module \cite{Carlson83,Carlson84}, see also \cite[Chapter~5]{BensonII}.

For \ref{S:pair}, use \(\Ext^n_{kP}(M,N)\cong H^n(P,M^*\otimes_kN)\) and the tensor product theorem of Avrunin--Scott \cite{AvruninScott82}, \(\VV_C(M^*\otimes_kN)=\VV_C(M^*)\cap\VV_C(N)=\VV_C(M)\cap\VV_C(N)\).  It remains to see that \(H^{\gg0}(P,L)=0\) iff \(L\) is projective.  Since \(k\) is the only simple \(kP\)-module, \(\Ext^n_{kP}(k,L)=0\) for all \(n>n_0\) means \(L\) has injective dimension at most \(n_0\); over the self-injective algebra \(kP\) a module of finite injective dimension is injective, hence projective.  The converse is clear.  \Cref{thm:support-classification,thm:no-generator} apply.
\end{proof}

\begin{remark}
For an elementary abelian \(p\)-group, \(kP\cong k[x_1,\dots,x_r]/(x_1^p,\dots,x_r^p)\) is a complete intersection and \Cref{thm:ci-classification} already applies; the new cases are the non-abelian \(p\)-groups, where \(kP\) is noncommutative and \(\PP_C\) may be reducible.  For the dihedral group of order \(8\) in characteristic \(2\), \(\Proj H^*(P,k)\) is two projective lines meeting in a point, and \(Z\) equal to one of the lines satisfies the hypothesis of \Cref{thm:no-generator}.  The theorem does not extend to arbitrary finite groups: for \(G=C_3\times C_2\) in characteristic \(3\), \(kG\cong k[t]/(t^3)\times k[t]/(t^3)\), the two simple modules have the same one-point support but \(\Ext^n_{kG}(M,N)=0\) for all \(n\), so \ref{S:pair} fails across blocks and \(\cU_N/\per\) is a nonzero proper defect not of the form \(\sD_Z\).  Any extension beyond \(p\)-groups must be formulated blockwise or with pair supports.
\end{remark}


\section{Relative dg models and the finite-total-Hom subcategory}\label{sec:dg-models}

For a homological universal localisation satisfying \textup{(B1)--(B3)}, the compact-kernel comparison identifies the compact-kernel quotient with the actual singularity kernel up to direct summands.  We now compute the compact-kernel quotient itself.

Let
\[
 T_\Sigma=\bigoplus_{\sigma\in\Sigma}\Cone(\sigma),
 \qquad
 A_\lambda=\End_R^\bullet(T_\Sigma)^{\op}.
\]
The opposite convention makes \(T_\Sigma\) a left \(R\)-complex and a right dg \(A_\lambda\)-module.  Put
\[
 \Psi_\Sigma=T_\Sigma\Ltensor_{A_\lambda}-:D(A_\lambda)\to D(R).
\]

\begin{theorem}[Relative dg formula at a thick scale]\label{thm:relative-dg-formula}
Assume that \(\Sigma\) is finite and that \(\lambda\) is stably flat.  Let \(\cB\subseteq D(R)\) be an essentially small thick subcategory containing \(\Kc_\lambda\), and define
\[
 K_{\cB}(\lambda)=\cK_\lambda\cap\cB
\]
and
\[
 D^{\cB}_\lambda(A_\lambda)
 =\{Y\in D(A_\lambda)\mid\Psi_\Sigma(Y)\in\cB\}.
\]
There is an exact equivalence
\[
 K_{\cB}(\lambda)/\Kc_\lambda
 \xrightarrow{\sim}
 D^{\cB}_\lambda(A_\lambda)/\per(A_\lambda).
\]
No noetherian hypothesis is required.
\end{theorem}

\begin{proof}
Neeman--Ranicki gives
\(\cK_\lambda=\Loc(T_\Sigma)\).  The object \(T_\Sigma\) is a compact generator of this localising category.  Keller--Morita theory therefore yields mutually inverse equivalences
\[
 \Psi_\Sigma:D(A_\lambda)\xrightarrow{\sim}\cK_\lambda,
 \qquad
 \Theta_\Sigma=\RHom_R(T_\Sigma,-)
\]
by Keller's dg Morita theorem~\cite[Theorem~8.2]{Keller94}.
For completeness, the unit and counit are isomorphisms on the generators
\(A_\lambda\) and \(T_\Sigma\); the full subcategories on which they are isomorphisms are localising, because \(T_\Sigma\) is compact.  Hence they are isomorphisms everywhere.

The equivalence restricts, by definition, to
\[
 D^{\cB}_\lambda(A_\lambda)\xrightarrow{\sim}K_{\cB}(\lambda).
\]
It also sends
\[
 \per(A_\lambda)=\thick(A_\lambda)
 \quad\text{to}\quad
 \thick(T_\Sigma)=\Kc_\lambda.
\]
Passing to Verdier quotients gives the statement.
\end{proof}

\begin{corollary}[Bounded finitely generated specialization]\label{cor:relative-dg-bounded}
Assume in addition that \(R\) is left noetherian, and put
\[
 \Db_\lambda(A_\lambda)
 =\{Y\in D(A_\lambda)\mid
   \Psi_\Sigma(Y)\in\Db(R\text{-}\modcat)\}.
\]
Then
\[
 \DeltaB(\lambda)
 \xrightarrow{\sim}
 \Db_\lambda(A_\lambda)/\per(A_\lambda).
\]
No noetherian assumption on \(S\), and neither \textup{(B2)} nor \textup{(B3)}, is required.
\end{corollary}

\begin{proof}
Apply \Cref{thm:relative-dg-formula} with \(\cB=\Db(R\text{-}\modcat)\).
\end{proof}

\begin{remark}[Why the noetherian specialization is necessary]\label{rem:nonnoetherian-bounded-scale}
Let \(R=k[x_1,x_2,\ldots]\) and \(I=(x_1,x_2,\ldots)\).  Both \(R\) and \(R/I\) are finitely generated \(R\)-modules, but the kernel \(I\) of \(R\twoheadrightarrow R/I\) is not finitely generated.  Hence \(R\text{-}\modcat\) is not abelian.  In \(D(R)\), the cone of \(R\to R/I\) has \(I\) as a cohomology module, so the objects with bounded finitely generated cohomology are not closed under cones.  Thus \(\Db(R\text{-}\modcat)\) cannot be used as a thick scale in this generality.
\end{remark}

The relative subcategory is external to the dg algebra.  Gao--Psaroudakis~\cite[Lemma~4.1]{GaoPsaroudakis18} characterize \(\Dfd(E)\) by finite total Hom from every perfect object.  Kuznetsov--Shinder~\cite{KuznetsovShinder25} develop the category of homologically finite-dimensional dg modules intrinsically for small dg-enhanced triangulated categories.  Our use concerns the image of this intrinsic category under a fixed dg Morita embedding of a localisation kernel into \(D(R)\), together with the additional question whether that image agrees with the ambient bounded kernel.  Transporting the Gao--Psaroudakis characterization through dg Morita equivalence gives the following reformulation.

\begin{proposition}[Generator-independent finite-total-Hom subcategory]\label{prop:finite-total-hom}
Let \(k\) be a field and let \(\mathcal K\) be a \(k\)-linear algebraic triangulated category with coproducts.  For any compact generator \(T\), put
\[
 E_T=\RHom(T,T)^{\op},
\]
and let \(F_T:D(E_T)\xrightarrow{\sim}\mathcal K\) be the corresponding dg Morita equivalence.  Then
\[
 F_T(\Dfd(E_T))=\mathcal F(\mathcal K),
\]
where
\begin{equation}\label{eq:finite-hom-subcategory}
 \mathcal F(\mathcal K)
 =\left\{
 X\in\mathcal K\ \middle|\
 \bigoplus_{n\in\mathbb Z}\Hom_{\mathcal K}(C,X[n])
 \text{ is finite-dimensional for every }C\in\mathcal K^c
 \right\}.
\end{equation}
In particular, \(F_T(\Dfd(E_T))\) is independent of the compact generator \(T\).
\end{proposition}

\begin{proof}
Under dg Morita equivalence, \(\per(E_T)\) corresponds to \(\mathcal K^c\).  The assertion is precisely Gao--Psaroudakis's characterization transported from \(D(E_T)\) to \(\mathcal K\).
\end{proof}

\begin{corollary}[Generator-independence of bounded compatibility]\label{cor:rbounded-generator-independent}
Let \(R\) be a finite-dimensional \(k\)-algebra and let \(\lambda\) be a finite homological universal localisation whose compact kernel admits a compact silting generator.  For any two compact silting generators \(T,T'\),
\[
 F_T(\Dfd(B_T))=F_{T'}(\Dfd(B_{T'}))
\]
inside \(\cK_\lambda\).  In particular,
\[
 F_T(\Dfd(B_T))=\Kb(\lambda)
 \quad\Longleftrightarrow\quad
 F_{T'}(\Dfd(B_{T'}))=\Kb(\lambda).
\]
Thus bounded compatibility is independent of the silting generator, even if the two generators lie in different components of the silting graph.
\end{corollary}

\begin{proof}
For a compact silting generator \(T\), the dg algebra
\(B_T=\tau_{\le0}\RHom_R(T,T)^{\op}\) is quasi-isomorphic to the full derived endomorphism algebra, because positive cohomology vanishes.  Hence its finite-dimensional derived category has the same image.  Apply \Cref{prop:finite-total-hom}.
\end{proof}

\begin{corollary}[Mutation stability without a heart-restriction hypothesis]\label{cor:mutation-stability}
Let \(R\) be finite-dimensional and let \(T,T'\) be compact silting generators of the same compact localisation kernel.  In particular, they may be related by a finite sequence of Aihara--Iyama mutations.  Then
\[
 T\text{ is }R\text{-bounded}
 \quad\Longleftrightarrow\quad
 T'\text{ is }R\text{-bounded}.
\]
When these conditions hold, both intrinsic dg quotients are equivalent to \(\DeltaB(\lambda)\).
\end{corollary}

\begin{proof}
This is the special case of \Cref{cor:rbounded-generator-independent} obtained by taking the two compact silting generators.  No assertion that a mutation torsion pair restricts to a finite-length heart is needed.
\end{proof}

\begin{remark}
The HRS-tilt description of silting mutation remains useful for identifying the changed heart; see Angeleri H\"ugel--Laking--\v{S}\v{t}ov\'{i}\v{c}ek--Vit\'{o}ria \cite{ALSV25}.  For bounded compatibility, the Gao--Psaroudakis characterization in \Cref{prop:finite-total-hom} yields the generator-independent statement and avoids the unresolved finite-heart restriction step.
\end{remark}


\section{Non-positive models and the heart-amplitude criterion}\label{sec:heart-amplitude}

Throughout this section, assume
\[
 \cK_\lambda=\Loc(\Kc_\lambda).
\]
In particular, this holds for a stably flat universal localisation: stable flatness gives
\[
 \cK_\lambda=\Loc\langle\Cone(\Sigma)\rangle,
 \qquad
 \Kc_\lambda=\thick\langle\Cone(\Sigma)\rangle.
\]

\begin{proposition}[Existence of a non-positive model]\label{prop:nonpositive-model}
The following are equivalent:
\begin{enumerate}[label=(\roman*)]
\item \(\cK_\lambda\simeq D(B)\) for a non-positive dg algebra \(B\);
\item \(\cK_\lambda\) has a compact silting generator;
\item \(\Kc_\lambda\) has a silting object.
\end{enumerate}
If \(T\) is such an object, one may take
\[
 B_T=\tau_{\le0}\RHom_R(T,T)^{\op}.
\]
\end{proposition}

\begin{proof}
If \(D(B)\simeq\cK_\lambda\) with \(B\) non-positive, the image of the free module \(B\) is compact, generates, and has no positive self-extensions.  Thus it is compact silting.

Under the standing hypothesis, the Thomason--Neeman theorem identifies the compact objects of \(\cK_\lambda\) with \(\Kc_\lambda\).  Hence a compact silting generator of \(\cK_\lambda\) belongs to \(\Kc_\lambda\) and thickly generates it.  Conversely, if \(T\) is a silting object of \(\Kc_\lambda\), then \(\Kc_\lambda=\thick(T)\), so the standing equality gives \(\cK_\lambda=\Loc(\Kc_\lambda)=\Loc(T)\).  Finally, for such \(T\), positive cohomology of \(\RHom_R(T,T)\) vanishes.  The good truncation is a quasi-isomorphism, and Keller--Morita identifies its derived category with \(\Loc(T)=\cK_\lambda\).
\end{proof}

Let \(T\) be a compact silting generator, put
\[
 B_T=\tau_{\le0}\RHom_R(T,T)^{\op},
 \qquad
 \Gamma_T=H^0(B_T),
\]
and let
\[
 F_T=T\Ltensor_{B_T}-:D(B_T)\xrightarrow{\sim}\cK_\lambda.
\]

\begin{definition}[Bounded compatibility]\label{def:r-bounded}
Assume \(\Gamma_T\) is left noetherian.  We call \(T\) \emph{\(R\)-bounded} if
\[
 F_T(\Dfg(B_T))=\Kb(\lambda).
\]
Over a field with \(\Gamma_T\) finite-dimensional, this reads
\(F_T(\Dfd(B_T))=\Kb(\lambda)\).
\end{definition}

By \Cref{cor:rbounded-generator-independent}, in finite dimension this is a property of the embedded kernel, not of the chosen silting generator.

\begin{lemma}\label{lem:d-bfg-heart}
If \(B\) is non-positive and \(H^0(B)\) is left noetherian, then
\[
 \Dfg(B)=\thick(H^0(B)\text{-}\modcat),
\]
where heart modules are viewed as dg \(B\)-modules via
\(B\to H^0(B)\).
\end{lemma}

\begin{proof}
The standard truncation triangle
\[
 \tau_{\le n-1}X\longrightarrow X\longrightarrow H^n(X)[-n]\longrightarrow
\]
reduces the cohomological amplitude.  Induction expresses every object of
\(\Dfg(B)\) as a finite extension of shifts of finitely generated heart modules.  The reverse inclusion holds because \(\Dfg(B)\) is thick.
\end{proof}

Define the \emph{heart bimodule}
\[
 M_T=T\Ltensor_{B_T}\Gamma_T.
\]
It is an \(R\)-\(\Gamma_T\) bimodule object.  For a finitely generated left \(\Gamma_T\)-module \(X\), derived associativity gives
\begin{equation}\label{eq:heart-realization}
 F_T(X)\simeq M_T\Ltensor_{\Gamma_T}X.
\end{equation}
For the general realization-functor perspective associated with hearts and tilting structures, see Psaroudakis--Vit\'{o}ria~\cite{PsaroudakisVitoria18}.

\begin{definition}[Fixed finite Tor-amplitude]\label{def:fixed-tor-amplitude}
The right \(\Gamma_T\)-complex \(M_T\) has \emph{finite Tor-amplitude} if there exist fixed integers \(a\le b\) such that, for every finitely generated left \(\Gamma_T\)-module \(X\),
\[
 M_T\Ltensor_{\Gamma_T}X\in\Db(R\text{-}\modcat)
 \quad\text{and}\quad
 H^n(M_T\Ltensor_{\Gamma_T}X)=0
 \quad(n\notin[a,b]).
\]
The endpoints are independent of \(X\).  Merely bounding the length of an interval whose position may vary with \(X\) is not sufficient and is not used here.
\end{definition}

\begin{proposition}[Noetherian direction]\label{prop:noetherian-heart}
Assume \(R\) and \(\Gamma_T\) are left noetherian.
\begin{enumerate}[label=(\alph*)]
\item If \(M_T\) has finite Tor-amplitude, then
\[
 F_T(\Dfg(B_T))\subseteq\Kb(\lambda).
\]
\item If, in addition,
\begin{equation}\label{eq:dagger}
 (\dagger_T)\qquad
 \RHom_R(T,X)\in\Dfg(B_T)
 \quad\text{for every }X\in\Kb(\lambda),
\end{equation}
then \(T\) is \(R\)-bounded.
\item If \(\Gamma_T\) is artinian and \(T\) is \(R\)-bounded, then \(M_T\) has finite Tor-amplitude.
\end{enumerate}
\end{proposition}

\begin{proof}
By \Cref{lem:d-bfg-heart}, \(\Dfg(B_T)\) is thickly generated by finitely generated \(\Gamma_T\)-modules.  Finite Tor-amplitude and \eqref{eq:heart-realization} place the image of every such heart module in \(\Kb(\lambda)\); exactness gives (a).

For (b), let \(X\in\Kb(\lambda)\).  The Morita counit is an isomorphism
\[
 F_T\RHom_R(T,X)\xrightarrow{\sim}X.
\]
Condition \((\dagger_T)\) places the inverse image in \(\Dfg(B_T)\), giving the reverse inclusion.

For (c), let \(L_1,\ldots,L_r\) be the finitely many simple \(\Gamma_T\)-modules.  Since \(T\) is \(R\)-bounded, each
\[
 M_T\Ltensor_{\Gamma_T}L_i\simeq F_T(L_i)
\]
lies in \(\Kb(\lambda)\).  Choose common lower and upper cohomological bounds \(a,b\) for this finite family.  Every finitely generated \(\Gamma_T\)-module has finite length.  Induction on length, using triangles induced by short exact sequences, gives the same fixed bounds \([a,b]\) for all such modules.
\end{proof}

\begin{theorem}[Finite-dimensional heart-amplitude criterion]\label{thm:heart-amplitude}
Let \(R\) be a finite-dimensional \(k\)-algebra, assume \(\cK_\lambda=\Loc(\Kc_\lambda)\), and let \(T\) be a silting object of \(\Kc_\lambda\), hence a compact silting generator of \(\cK_\lambda\).  Then the following are equivalent:
\begin{enumerate}[label=(\roman*)]
\item \(T\) is \(R\)-bounded;
\item \(M_T=T\Ltensor_{B_T}\Gamma_T\) has fixed finite Tor-amplitude over \(\Gamma_T\).
\end{enumerate}
When these conditions hold,
\[
 \DeltaB(\lambda)
 \simeq
 \Dfd(B_T)/\per(B_T).
\]
Moreover, by \Cref{cor:rbounded-generator-independent}, if the condition holds for one compact silting generator, it holds for every compact silting generator of the kernel.
\end{theorem}

\begin{proof}
The implication (i)\(\Rightarrow\)(ii) is
\Cref{prop:noetherian-heart}(c), since \(\Gamma_T\) is finite-dimensional.

Assume (ii).  Part (a) of \Cref{prop:noetherian-heart} gives
\[
 F_T(\Dfd(B_T))\subseteq\Kb(\lambda).
\]
Conversely, if \(X\in\Kb(\lambda)\), then \(X\) is represented by a bounded complex of finite-dimensional \(R\)-modules, while \(T\) is a bounded complex of finitely generated projectives.  Hence
\(\RHom_R(T,X)\) has finite-dimensional total cohomology, so it belongs to \(\Dfd(B_T)\).  The Morita counit identifies its image with \(X\).  Thus equality holds and \(T\) is \(R\)-bounded.

Finally, \(F_T\) sends \(\per(B_T)\) to
\(\thick(T)=\Kc_\lambda\).  Quotienting gives the formula.  Generator independence is \Cref{cor:rbounded-generator-independent}.
\end{proof}



\section{The one-sided finite-dimensional Kernel Dictionary}\label{sec:dictionary}

Throughout this section, let \(\lambda:R\to S\) be a homological ring epimorphism of finite-dimensional \(k\)-algebras and assume
\[
 S_R\in\per(R^{\op}).
\tag{E}
\]
Put \(F=S\Ltensor_R-\), \(S^\vee=\RHom_{R^{\op}}(S,R)\), and \(F_!=S^\vee\Ltensor_S-\).

\begin{lemma}[Dual adjoint]\label{lem:dual-adjoint}
Under the standing hypotheses, \(F_!\dashv F\), and \(F_!\) is fully faithful.
\end{lemma}

\begin{proof}
Choose a bounded resolution of \(S_R\) by finitely generated projectives.  Termwise finite-projective duality gives
\[
 \RHom_R(S^\vee,X)\simeq S\Ltensor_RX
\]
as left \(S\)-complexes, and tensor--Hom adjunction gives \(F_!\dashv F\).  The same duality gives
\[
 S\Ltensor_RS^\vee\simeq\RHom_{R^{\op}}(S,S).
\]
Homologicality makes restriction on right derived categories fully faithful, so the last complex is \(\RHom_{S^{\op}}(S,S)\simeq S\).  Thus the unit \(\id\to FF_!\) is invertible.
\end{proof}

\begin{proposition}[Compact kernel projection]\label{prop:left-projection}
Under the standing hypotheses, the inclusion \(j:\cK_\lambda\hookrightarrow D(R)\) has a left adjoint \(\ell\), defined by triangles
\[
 F_!F(X)\longrightarrow X\longrightarrow j\ell X\longrightarrow.
\]
It preserves compact objects, and \(G=\ell R\) is a compact generator satisfying
\[
 \Hom_{\cK_\lambda}(G,X[n])\cong H^n(X).
\]
As an object of \(D(R)\), \(G\) is perfect and satisfies \(\Hom(G,G[n])=0\) for \(n<-1\).  If \(S_R\) is projective, \(G\) has a projective representative in degrees \([-1,0]\), and \(\RHom_R(G,G)^{\op}\) has cohomology only in degrees \(-1,0\).
\end{proposition}

\begin{proof}
The cone of the counit lies in the kernel because applying \(F\) makes the first map an isomorphism.  For \(K\) in the kernel, adjunction gives \(\RHom_R(F_!F(X),K)=0\), whence
\[
 \RHom_{\cK_\lambda}(\ell X,K)\simeq\RHom_R(X,K).
\]
Thus \(\ell\dashv j\).  Since \(j\) preserves coproducts, \(\ell\) preserves compactness.  Taking \(X=R\) gives the Hom formula, which detects zero objects.

At \(R\), the counit triangle identifies \(jG\) with \(\Cone(S^\vee\to R)\), so it is perfect over \(R\).  Choose a projective resolution of \(S_R\) in degrees \([-p,0]\).  Its finite-projective dual represents \(S^\vee\) in degrees \([0,p]\), and the cone has no terms below \(-1\).  Taking \(X=G\) in the Hom formula gives the claimed negative self-Hom vanishing.  If \(S_R\) is projective, the cone is two-term, and the same formula also eliminates all positive self-Homs.
\end{proof}

\begin{theorem}[Finite-dimensional Kernel Dictionary]\label{thm:kernel-dictionary}
Let \(\lambda:R\to S\) be a finite stably flat universal localisation of finite-dimensional \(k\)-algebras and assume \(S_R\) has finite projective dimension.  Then the raw kernel dga \(\Lambda\) is proper and
\[
 \Psi(\Dfd\Lambda)=\Kb(\lambda).
\]
Consequently
\[
 \Dsg(\Lambda):=\Dfd(\Lambda)/\per(\Lambda)
 \xrightarrow{\sim}\DeltaB(\lambda),
\]
and the composite to \(\Defsg(\lambda)\) is fully faithful and Karoubi dense.
\end{theorem}

\begin{proof}
By \Cref{prop:finite-total-hom}, the compact cone generator and \(G=\ell R\) define the same finite-total-Hom subcategory.  For \(G\), the condition is
\[
 \sum_n\dim_kH^n(X)<\infty.
\]
Over finite-dimensional \(R\) this is equivalent to \(X\in\Db(R\text{-}\modcat)\).  Intersecting with the kernel proves the first equality.  Properness of \(\Lambda\) is automatic because the cone generator is a bounded complex of finite-dimensional projectives.  The bounded specialization in \Cref{cor:relative-dg-bounded} becomes the intrinsic quotient, and \Cref{thm:unconditional-comparison} supplies the comparison.
\end{proof}

\subsection{Why the AHKLY map is a universal localisation}
The map \(\lambda:A\to S=M_2(k)\) fixed in \Cref{sec:preliminaries} is homological because \(S_A\cong(e_1A)^{\oplus2}\) is projective.  We verify directly that it is the universal localisation at \(\sigma_L:Ae_1\to Ae_2\).  Let \(f:A\to C\) be a ring map for which \(C\otimes_A\sigma_L\) is invertible, and put \(a=f(\alpha)\).  Its inverse is right multiplication by a unique element
\(b\in f(e_2)Cf(e_1)\) satisfying
\[
 ab=f(e_1),
 \qquad
 ba=f(e_2).
\]
Since \(\alpha\beta=0\),
\[
 f(\beta)=f(e_2)f(\beta)=ba f(\beta)=b f(\alpha\beta)=0.
\]
Thus \(f(e_1),f(e_2),a,b\) form a complete set of \(2\times2\) matrix units and determine a unique homomorphism \(M_2(k)\to C\) through which \(f\) factors.  Therefore \(\lambda\) is the universal localisation at \(\sigma_L\), over an arbitrary base field.


\section{The one- and two-periodic boundary examples}\label{sec:periodic-ladder}

\subsection{The AHKLY localisation on the left-module side}
For the algebra and presentation fixed in \Cref{sec:preliminaries}, the cone of \(\sigma_L\) is quasi-isomorphic to its cokernel
\[
 U=\langle e_2,\beta\rangle.
\]
Its endomorphism algebra is \(E=k[\eps]/(\eps^2)\), and \(U\cong E\) as a left \(E\)-module.

\begin{theorem}[One-periodic AHKLY defect]\label{thm:ahkly-left}
In left modules,
\[
 \DeltaB(\lambda)\simeq\Dsg(E)\simeq\vect_k,
\]
with suspension isomorphic to the identity.  The comparison with \(\Defsg(\lambda)\) is an equivalence before idempotent completion.  The old perfect-restriction hypothesis fails because \({}_AS\) has infinite projective dimension.
\end{theorem}

\begin{proof}
The right source module \(S_A\) is projective, so \Cref{thm:kernel-dictionary} applies.  The dual-numbers algebra is self-injective local; every finite module is a direct sum of copies of its regular module and its simple \(k\).  The minimal resolution
\[
 \cdots\xrightarrow{\eps}E\xrightarrow{\eps}E\twoheadrightarrow k
\]
has period one.  Hence the stable category is the split category generated by one object fixed by suspension.  It is idempotent complete, so the comparison is an equivalence before idempotent completion.  The supplemental side-verification script gives the infinite periodic left resolution of \(S\).
\end{proof}

\subsection{The opposite right-module boundary}

This section uses right modules.  Passing to opposite rings returns to the left-module convention, and the kernel algebra below is self-opposite.  The underlying recollement example is Angeleri H\"ugel--Koenig--Liu--Yang's Example~4.13(b) \cite{AHKLY17}; the universal-localisation identification, the kernel-dga computation, and the bounded-defect conclusion below are additional.

Let \(A=kQ/I\), where
\[
 Q:\quad 1\xrightarrow{\alpha}2,
 \qquad \beta:2\to2,
\]
and
\[
 \beta^2=0,
 \qquad
 \alpha\beta=0.
\]
Let \(P_i=e_iA\), and let
\[
 \sigma=\alpha:P_2\longrightarrow P_1
\]
be left multiplication by \(\alpha\).  Put
\[
 C=\Cone(\sigma)=[P_2\xrightarrow{\alpha}P_1].
\]

\begin{proposition}\label{prop:dual-localisation}
The universal localisation of \(A\) at \(\sigma\) is
\[
 \lambda:A\longrightarrow M_2(k),
\]
given by
\[
 e_1\mapsto E_{11},\quad
 e_2\mapsto E_{22},\quad
 \alpha\mapsto E_{12},\quad
 \beta\mapsto0.
\]
It is a homological ring epimorphism.
\end{proposition}

\begin{proof}
Left multiplication by \(E_{12}\) gives an isomorphism
\[
 E_{22}M_2(k)\xrightarrow{\sim}E_{11}M_2(k)
\]
with inverse left multiplication by \(E_{21}\), so the map inverts \(\sigma\).

If \(f:A\to T\) inverts \(\sigma\), there is a unique
\(\gamma\in f(e_2)Tf(e_1)\) such that
\[
 f(\alpha)\gamma=f(e_1),
 \qquad
 \gamma f(\alpha)=f(e_2).
\]
Then
\[
 f(\beta)=f(e_2)f(\beta)
 =\gamma f(\alpha)f(\beta)=0.
\]
The elements \(f(e_1),f(e_2),f(\alpha),\gamma\) satisfy the matrix-unit relations and produce a unique factorisation through \(M_2(k)\).  Thus the map is universal.

As a right \(A\)-module, \(M_2(k)\cong P_1^{\oplus2}\).  In particular it is projective, so
\(\Tor_i^A(M_2(k),M_2(k))=0\) for \(i>0\).
\end{proof}

\begin{proposition}\label{prop:dual-dga}
The dg endomorphism algebra of \(C\) is quasi-isomorphic to the ordinary dual numbers
\[
 E=k[\eps]/(\eps^2)
\]
concentrated in degree zero.  As a left \(E\)-complex, the realisation bimodule is
\[
 C\simeq k\oplus k[1].
\]
\end{proposition}

\begin{proof}
The mapping complex \(\End_A^\bullet(C)\) has
\[
 \End^{-1}(C)=\Hom_A(P_1,P_2)=0,
\]
\[
 \End^0(C)=
 \End_A(P_2)\oplus\End_A(P_1)
 =(ke_2\oplus k\beta)\oplus ke_1,
\]
and
\[
 \End^1(C)=\Hom_A(P_2,P_1)=k\alpha.
\]
For \((ae_2+b\beta,ce_1)\in\End^0(C)\), the differential is
\[
 (ae_2+b\beta,ce_1)\longmapsto(c-a)\alpha.
\]
It is surjective.  Its kernel consists of
\((ae_2+b\beta,ae_1)\), and the class of \((\beta,0)\) generates a square-zero radical.  Thus
\[
 H^0\End_A^\bullet(C)\cong k[\eps]/(\eps^2),
 \qquad H^i=0\ (i\ne0),
\]
which gives the quasi-isomorphism.

The element \(\eps\) acts on \(C\) by the chain endomorphism
\((\beta,0)\).  As a left \(E\)-module, \(P_2\cong E\), while
\(P_1=ke_1\oplus k\alpha\) is semisimple and \(\eps\) acts by zero.  The differential sends \(e_2\mapsto\alpha\) and \(\beta\mapsto0\).  Hence
\[
 C\cong[E\twoheadrightarrow k\alpha]\oplus[0\to ke_1]
 \simeq k[1]\oplus k
\]
in \(D(E)\).
\end{proof}

\begin{theorem}[Dual-numbers boundary theorem]\label{thm:dual-numbers}
For the localisation above,
\[
 \Db_\lambda(E)=\per(E).
\]
Consequently
\[
 \DeltaB(\lambda)=0,
 \qquad
 \Dfd(E)/\per(E)=\Dsg(E)\ne0.
\]
Thus a connective ordinary kernel algebra does not determine the bounded localisation defect through its intrinsic finite-dimensional singularity category.
\end{theorem}

\begin{proof}
By \Cref{prop:dual-dga}, for \(Y\in D(E)\), bounded realisation is equivalent to bounded finite-dimensional cohomology of
\[
 Y\Ltensor_Ek,
\]
because
\(Y\Ltensor_EC\simeq(Y\Ltensor_Ek)\oplus(Y\Ltensor_Ek)[1]\).

Choose a pure-minimal semi-projective resolution \(P\to Y\).  Such resolutions exist over the left perfect ring \(E\) by Christensen--Thompson \cite[Theorem~5.7]{ChristensenThompson}; for complexes of projectives over a left perfect ring, pure-minimality agrees with the usual minimality by \cite[Theorem~3.13 and Corollary~3.14]{ChristensenThompson}.  Minimality forces
\[
 d(P^n)\subseteq(\eps)P^{n+1}
\]
for all \(n\): a unit entry in a differential matrix would split off a contractible rank-one two-term summand.  Tensoring with
\(k=E/(\eps)\) therefore kills all differentials, and
\[
 H^n(P\otimes_Ek)=P^n/(\eps)P^n.
\]
If \(Y\Ltensor_Ek\) has bounded finite-dimensional cohomology, then only finitely many \(P^n/(\eps)P^n\) are nonzero and each is finite-dimensional.  Since every \(P^n\) is free over the artinian local ring \(E\), only finitely many \(P^n\) are nonzero and each has finite rank.  Hence \(P\) is a bounded complex of finitely generated projectives and \(Y\in\per(E)\).

Conversely, perfect dg modules clearly realise boundedly.  Therefore
\(\Db_\lambda(E)=\per(E)\), and \Cref{cor:relative-dg-bounded} gives \(\DeltaB=0\).

The simple module \(k\) has the infinite periodic resolution
\[
 \cdots\xrightarrow{\eps}E\xrightarrow{\eps}E
 \xrightarrow{\eps}E\twoheadrightarrow k\to0.
\]
Thus \(k\in\Dfd(E)\setminus\per(E)\), proving the second quotient is nonzero.
\end{proof}

\begin{corollary}[Refutation of cone-only absorption]\label{cor:false-cone-absorption}
Stable flatness together with absorption of the defining cones does not imply relative singularity invariance.
\end{corollary}

\begin{proof}
Take the relative denominator to be the regular module \(M=A\).  Then
\(\Cone(\sigma)\in\Kb(\operatorname{proj}A)=\per(A)\), so the defining cone is absorbed.  The old cone-only assertion would give
\[
 \Dsg(A)\simeq\Dsg(M_2(k))=0.
\]
But the simple right module \(S_2\) has the infinite periodic resolution
\[
 \cdots\xrightarrow{\beta}P_2\xrightarrow{\beta}P_2
 \xrightarrow{\beta}P_2\twoheadrightarrow S_2\to0,
\]
so its class in \(\Dsg(A)\) is nonzero.  The failure occurs because the relevant bounded extension property is absent; after tensoring with \(M_2(k)\), the maps induced by \(\beta\) vanish and the derived tensor is unbounded.  The compact cone does not control the full bounded kernel.
\end{proof}


\subsection{Boundary with Jin--Yang--Zhou}\label{subsec:jyz-boundary}
\begin{lemma}[Fixed-side non-vacuousness]\label{lem:jyz-boundary}
The left-module application of \Cref{thm:kernel-dictionary} to the AHKLY localisation does not satisfy the downward-extension hypothesis of Jin--Yang--Zhou's Theorem~1.1~\cite{JinYangZhou}.  In AHKLY's original right-module convention, the recollement extends downward because \(i_*(S)=P_1\oplus P_1\) is compact.  Passing to opposites changes the controlling module to \({}_AS\), which is not perfect.
\end{lemma}

\begin{proof}
AHKLY compute the right-module recollement in Example~4.13(b)~\cite{AHKLY17}; their compactness criterion makes the displayed right-module object compact.  In the global left convention the same statement is the right-module recollement over \(A^{\op}\), and the object controlling downward extension is \({}_AS\).  The explicit periodic resolution in the supplemental computation has nonzero syzygies in every degree, so \({}_AS\notin\per A\).  AHKLY's correction in Section~4.4 records Example~4.13(b) as a failure of the missing two-sided finite-projective-dimension condition.  Thus the downward hypothesis fails on precisely the side where the Kernel Dictionary applies.
\end{proof}

Even when both theories apply, their roles differ: Jin--Yang--Zhou identify the quotient in a singularity sequence obtained from a downward recollement, while the Kernel Dictionary identifies the compactly generated kernel.

\subsection{Central localisations and the node}
\begin{lemma}[Central nilpotence]\label{lem:central-nilpotence}
Let \(z\in Z(R)\) act trivially on the cohomology of \(X\in\Db(R\text{-}\modcat)\).  If \(X\) has cohomology in at most \(m\) degrees, then \(z_X^m=0\) in \(D(R)\).
\end{lemma}

\begin{proof}
Induct on \(m\).  For maximal \(b\) with \(H^b(X)\ne0\), use the truncation triangle
\[
 \tau_{\le b-1}X\xrightarrow{i}X\xrightarrow{p}H^b(X)[-b]\longrightarrow.
\]
Since \(pz_X=0\), exactness gives \(a:X\to\tau_{\le b-1}X\) with \(z_X=ia\).  Centrality gives \(az_X=z_{\tau_{\le b-1}X}a\), and induction yields \(z_X^m=0\).
\end{proof}

\begin{proposition}[Central-localisation factorisation]\label{prop:central-localisation}
For a finite central localisation \(R\to R[s_1^{-1},\ldots,s_r^{-1}]\), every map from a bounded kernel object to a perfect object factors through the compact kernel.  Hence the comparison is fully faithful and Karoubi dense.
\end{proposition}

\begin{proof}
A common monomial \(u\) in the \(s_i\) annihilates the finitely generated cohomology of a bounded kernel object \(X\).  By \Cref{lem:central-nilpotence}, a power \(u^m\) is zero on \(X\).  For \(f:X\to P\) with \(P\) perfect, centrality gives \((u^m)_Pf=0\).  The multiplication triangle on \(P\) factors \(f\) through \(\Cone(u^m:P\to P)[-1]\), which is perfect and killed by localisation.
\end{proof}

\begin{theorem}[Two-periodic node defect]\label{thm:node}
For
\[
 R=k[[x,y]]/(xy)\longrightarrow R[1/x]\cong k((x)),
\]
the comparison is an equivalence before idempotent completion and
\[
 \DeltaB(\lambda)\simeq\Defsg(\lambda)\simeq\Dsg(R)
 \simeq\vect_k^{\mathbb Z/2}.
\]
Perfect restriction fails.
\end{theorem}

\begin{proof}
The strict cone is \(T=[R\xrightarrow{x}R]\), not \(R/(x)\).  Its endomorphism dga before taking an opposite has
\[
 E^{-1}=R,\quad E^0=R^2,\quad E^1=R,
\]
\[
 d^{-1}(h)=(xh,xh),\qquad d^0(a,b)=x(a-b).
\]
The comparison is full and Karoubi dense by \Cref{prop:central-localisation}.

We classify maximal Cohen--Macaulay modules directly.  This is the standard split-node classification; compare Greuel--Kn\"orrer~\cite{GreuelKnorrer85} under their algebraically closed characteristic-zero convention.  The argument below is included because the present split calculation is valid over an arbitrary field.

Write \(A=k[[x]]\), \(B=k[[y]]\), and
\[
 R=A\times_kB\subseteq\widetilde R=A\oplus B,
 \qquad
 \mathfrak c=xA\oplus yB.
\]
A finite maximal Cohen--Macaulay module \(X\) is torsion-free.  Inside \(Q(R)\otimes_RX\), put
\(\widetilde X=\widetilde R X\cong A^p\oplus B^q\).  Since
\(\mathfrak c\widetilde R\subseteq R\), one has
\[
 \mathfrak c\widetilde X\subseteq X\subseteq\widetilde X.
\]
Thus \(X\) is the inverse image of
\[
 W=X/\mathfrak c\widetilde X
 \subseteq
 \widetilde X/\mathfrak c\widetilde X
 \cong k^p\oplus k^q.
\]
The two projections of \(W\) are surjective: \(X\) generates \(\widetilde X\) over \(\widetilde R\), and Nakayama's lemma applied on each branch gives the assertion.  With
\[
 K_x=W\cap(k^p\oplus0),
 \qquad
 K_y=W\cap(0\oplus k^q),
\]
the subspace \(W\) induces an isomorphism
\[
 k^p/K_x\xrightarrow{\sim}k^q/K_y
\]
and therefore, after choosing complements and bases,
\[
 W=K_x\oplus K_y\oplus\operatorname{graph}(\id_{k^a}).
\]
Every change of basis in \(\operatorname{GL}_p(k)\times\operatorname{GL}_q(k)\) lifts to one in \(\operatorname{GL}_p(A)\times\operatorname{GL}_q(B)\).  Hence this normal form is induced by an automorphism of \(\widetilde X\), not merely by an abstract vector-space isomorphism, and taking inverse images preserves the \(R\)-module isomorphism problem.  The branch coordinates lift to \(R/(y)\) and \(R/(x)\), while each graph coordinate lifts to \(R\).  Consequently
\[
 X\cong R^a\oplus(R/(y))^b\oplus(R/(x))^c.
\]
Thus the nonfree indecomposables are \(N=R/(y)\) and \(M=R/(x)\).  One has \(\End_R(N)=k[[x]]\), and the endomorphisms factoring through a free module form the ideal \((x)\); hence \(\underline{\End}_R(N)=k\).  Similarly \(\underline{\End}_R(M)=k\), while \(\Hom_R(M,N)=\Hom_R(N,M)=0\).  The projective-cover kernels are \((x)\cong N\) and \((y)\cong M\), so \(\Omega M\cong N\) and \(\Omega N\cong M\).  Row and column reduction reduces every stable morphism to identities and zero maps.  Buchweitz's Theorem~4.4.1~\cite{Buchweitz} identifies the singularity category of this Gorenstein ring with the stable MCM category.  Under this equivalence, suspension swaps \(M\) and \(N\); hence the additive equivalence with \(\vect_k^{\mathbb Z/2}\) respects suspension and its split triangulated structure.

The object \(M\) is killed by localisation and is nonperfect.  Its suspension is \(N\), so the image of \(\DeltaB\) contains both generators; Karoubi density is therefore essential surjectivity before idempotent completion.  Finally \(k((x))\) is not finitely generated over \(R\), so it is not perfect.
\end{proof}

\section{Stratifying idempotent quotients}\label{sec:stratifying}

Let \(A\) be a finite-dimensional \(k\)-algebra, let \(e\in A\) be an idempotent, and put
\[
 B=A/AeA,
 \qquad
 C=eAe.
\]
Assume that \(AeA\) is stratifying, so the multiplication map
\[
 Ae\Ltensor_CeA\xrightarrow{\sim}AeA
\]
is an isomorphism in \(D(A^e)\), and \(A\to B\) is homological.

\begin{lemma}\label{lem:idempotent-universal}
The quotient \(A\to A/AeA\) is the universal localisation at
\(0\to Ae\).  Its compact kernel is \(\thick(Ae)\).
\end{lemma}

\begin{proof}
For a ring map \(f:A\to T\), the map
\[
 T\otimes_A0\longrightarrow T\otimes_AAe\cong Te
\]
is an isomorphism if and only if \(f(e)=0\), equivalently
\(AeA\subseteq\ker f\).  This is exactly the universal property of the quotient by \(AeA\).  The cone is \(Ae\), up to shift.
\end{proof}

Define
\[
 \cB_e^b=
 \{Y\in\Db(C\text{-}\modcat)
   \mid Ae\Ltensor_CY\in\Db(A\text{-}\modcat)\}.
\]

\begin{theorem}[Stratifying computation]\label{thm:stratifying-delta}
The functor
\[
 j_!=Ae\Ltensor_C-:\cB_e^b\longrightarrow\Kb(\lambda)
\]
is an exact equivalence with quasi-inverse
\(j^*=eA\otimes_A-\).  It sends \(\per(C)\) onto \(\thick(Ae)\).  Hence
\[
 \DeltaB(\lambda)\simeq\cB_e^b/\per(C).
\]
If \(Ae_C\) has finite Tor-amplitude, then
\[
 \cB_e^b=\Db(C\text{-}\modcat),
 \qquad
 \DeltaB(\lambda)\simeq\Dsg(C).
\]
\end{theorem}

\begin{proof}
Since \(eA_A\) is projective, \(j^*=eA\otimes_A-\) is underived.  Multiplication gives
\[
 j^*j_!(Y)
 \simeq
 (eA\otimes_AAe)\Ltensor_CY
 \cong C\Ltensor_CY\simeq Y,
\]
so \(j_!\) is fully faithful.

The stratifying triangle of \(A\)-bimodules is
\[
 Ae\Ltensor_CeA\longrightarrow A\longrightarrow B\longrightarrow.
\]
Tensoring it over \(A\) with \(X\in D(A)\) gives
\[
 j_!j^*X\longrightarrow X\longrightarrow B\Ltensor_AX\longrightarrow.
\]
If \(X\in\Kb(\lambda)\), the third term is zero, so
\(X\simeq j_!j^*X\), and \(j^*X\in\cB_e^b\).  Conversely, if
\(Y\in\cB_e^b\), then
\[
 B\Ltensor_Aj_!Y
 \simeq
 (B\Ltensor_AAe)\Ltensor_CY=0
\]
because \(Be=0\).  Thus \(j_!Y\in\Kb(\lambda)\), proving the equivalence.

Since \(j_!(C)=Ae\), exactness and full faithfulness identify
\(j_!(\per(C))\) with \(\thick(Ae)\).  Passing to quotients gives the first formula.  If \(Ae_C\) has finite Tor-amplitude, \(j_!\) sends all bounded finitely generated \(C\)-complexes to bounded \(A\)-complexes, so \(\cB_e^b=\Db(C\text{-}\modcat)\).
\end{proof}

The two source sides of the stratifying bimodules have different roles.  We record them explicitly because they are exactly what is needed for the singularity comparison.

\begin{lemma}[Bimodule perfectness]\label{lem:stratifying-perfectness}
Assume \(Ae_C\in\per(C^{\op})\).  Then \(B_A\in\per(A^{\op})\), so extension of scalars has finite Tor-amplitude.  Assume instead \({}_CeA\in\per(C)\).  Then \({}_AB\in\per(A)\), so restriction of scalars preserves perfect complexes.
\end{lemma}

\begin{proof}
Choose a bounded resolution of \(Ae_C\) by finitely generated projective right \(C\)-modules.  Tensoring it over \(C\) with the projective right \(A\)-module \(eA_A\) gives a bounded complex of finitely generated projective right \(A\)-modules representing
\(Ae\Ltensor_CeA\).  The stratifying triangle then shows
\(B_A\in\per(A^{\op})\).

Dually, resolve \({}_CeA\) by finitely generated projective left \(C\)-modules.  Tensoring with the projective left \(A\)-module \({}_AAe\) represents
\(Ae\Ltensor_CeA\) by a bounded complex of finitely generated projective left \(A\)-modules.  The same triangle gives \({}_AB\in\per(A)\).
\end{proof}

\begin{theorem}[Corner singularity as the actual kernel]\label{thm:stratifying-comparison}
Assume both
\[
 Ae_C\in\per(C^{\op}),
 \qquad
 {}_CeA\in\per(C).
\]
Then the bounded hypotheses and perfect restriction hold for
\(A\to B\), and the comparison functor is an equivalence before idempotent completion
\[
 \rho_\lambda:\DeltaB(\lambda)
 \xrightarrow{\sim}
 \Defsg(\lambda).
\]
Under the identification of \Cref{thm:stratifying-delta}, this gives an exact sequence of the displayed categories
\[
 \Dsg(C)\longrightarrow\Dsg(A)\longrightarrow\Dsg(B)
\]
and
\[
 \Defsg(\lambda)\simeq\Dsg(C).
\]
\end{theorem}

\begin{proof}
Finite-dimensionality gives noetherianity and bounded restriction on finitely generated modules.  The first perfectness assumption and
\Cref{lem:stratifying-perfectness} give finite Tor-amplitude of \(B_A\); the second gives \({}_AB\in\per(A)\), which is perfect restriction.  Thus \Cref{thm:perfect-restriction} applies.

The first assumption also gives
\(\DeltaB(\lambda)\simeq\Dsg(C)\) by
\Cref{thm:stratifying-delta}.  The singularity category of a finite-dimensional algebra is idempotent complete \cite[Corollary~2.4]{ChenRadicalSquareZero}.  Therefore \(\DeltaB(\lambda)\) is idempotent complete, and the Karoubi-dense fully faithful comparison is already essentially surjective.  The exact sequence is \Cref{thm:actual-kernel} with its kernel identified.
\end{proof}

\begin{example}[A nonzero corner defect]\label{ex:triangular}
Let \(C=k[t]/(t^3)\) and
\[
 A=T_2(C)=
 \begin{pmatrix}C&C\\0&C\end{pmatrix},
 \qquad e=e_{22}.
\]
Then \(A/AeA\cong C\), both corner bimodules are projective, and
\[
 \DeltaB(\lambda)
 \simeq
 \Defsg(\lambda)
 \simeq
 \Dsg(C)\ne0.
\]
A left \(A\)-module is a triple \((X,Y,\varphi:Y\to X)\), and
\[
 j_!(Y)=(Y,Y,\id).
\]
The unique simple \(k=C/(t)\) has periodic resolution
\[
 \cdots\xrightarrow{t^2}C\xrightarrow{t}C
 \xrightarrow{t^2}C\xrightarrow{t}C\twoheadrightarrow k,
\]
so the explicit extra bounded-kernel generator is
\(j_!(k)=(k,k,\id)\), beyond the compact generator
\(Ae=(C,C,\id)\).
\end{example}


\section{Marks-periodic self-injective Nakayama localisations}\label{sec:nakayama}

Throughout this section, \(k\) is algebraically closed.  This hypothesis is used when invoking Marks' classification and projectivity results.  The explicit path-space and matrix calculations below are characteristic-independent, but they do not by themselves extend the classification theorem to arbitrary fields.

This section uses right modules.  Equivalently, it applies the preceding left-module results to the opposite algebra.  The cyclic Nakayama algebras below are isomorphic to their opposites after reversing the cycle.

Let
\[
 A_d=\widetilde A^d_{2d}
 =k\widetilde Q_{2d}/\mathfrak r^d,
 \qquad d\ge3,
\]
where
\[
 \widetilde Q_{2d}:1\xrightarrow{\alpha_1}2\longrightarrow\cdots
 \longrightarrow2d\xrightarrow{\alpha_{2d}}1
\]
and \(\mathfrak r\) is the arrow ideal.  Put \(P_i=e_iA_d\).  Then
\[
 P_i=(S_i,S_{i+1},\ldots,S_{i+d-1})
\]
is uniserial of length \(d\).  The standard projective presentation of \(S_i\) is
\[
 P_{i+1}\xrightarrow{\alpha_i}P_i\twoheadrightarrow S_i\to0,
\]
and we write
\[
 C_i=[P_{i+1}\xrightarrow{\alpha_i}P_i]
\]
with degrees \(-1,0\).

Under this standing hypothesis, Marks' classification gives non-semisimple homological epimorphisms indexed by periodic simple collections
\begin{equation}\label{eq:marks-collection}
 \Sigma_J=\{S_j,S_{j+d}\mid j\in J\},
 \qquad
 \varnothing\ne J\subset\mathbb Z/d\mathbb Z,
 \qquad
 |J|\le d-2.
\end{equation}
We denote the corresponding universal localisation by
\(\lambda_J:A_d\to S_J\).

\subsection{Two-sided projectivity and the family comparison}

We first record a consequence of Marks' projectivity theorem which is useful beyond the subfamily computed below.

\begin{theorem}[Comparison for self-injective Nakayama epimorphisms]\label{thm:all-nakayama-comparison}
Let \(k\) be an algebraically closed field, let \(A\) be a connected finite-dimensional self-injective Nakayama \(k\)-algebra, and let
\(\lambda:A\to S\) be a finite-dimensional homological ring epimorphism occurring in Marks' classification.  Then \(\lambda\) is a universal localisation at a finite collection of maps between finitely generated projectives, and \(S\) is projective both as a left and as a right \(A\)-module.  Consequently the bounded localisation hypotheses and perfect restriction hold.  The comparison
\[
 \rho_\lambda:\DeltaB(\lambda)\longrightarrow\Defsg(\lambda)
\]
is fully faithful and induces an equivalence
\[
 \Kar\DeltaB(\lambda)\xrightarrow{\sim}\Defsg(\lambda).
\]
\end{theorem}

\begin{proof}
A connected self-injective Nakayama algebra is representation-finite, hence belongs to the classes for which the Tachikawa condition used in \cite[Corollary~3.7]{Marks15} holds.  Marks' corollary and classification theorem \cite[Theorem~5.14]{Marks15} show that every such finite-dimensional homological epimorphism is a universal localisation, that its target is again self-injective, and that \({}_AS\) is projective.  The classification uses a finite orthogonal collection, so the localisation is finite.  The opposite map
\[
 A^{\op}\longrightarrow S^{\op}
\]
is again a finite-dimensional homological ring epimorphism from a connected self-injective Nakayama algebra.  Applying the same result to the opposite map shows that \(S_A\) is projective.  Thus derived extension and restriction preserve bounded finite-dimensional complexes, while restriction sends perfect \(S\)-complexes to perfect \(A\)-complexes.  Apply \Cref{thm:unconditional-comparison}.  Finally \(\Dsg(A)\) is idempotent complete by \cite[Corollary~2.4]{ChenRadicalSquareZero}, so its thick subcategory \(\Defsg(\lambda)\) is idempotent complete.
\end{proof}

\begin{proposition}[Projectivity of the periodic-family target]\label{prop:marks-projectivity}
Assume that \(k\) is algebraically closed.  For every localisation \(\lambda_J\) in \eqref{eq:marks-collection}, the finite-dimensional target is projective both as a left and as a right \(A_d\)-module.  Consequently (B1)--(B3) and perfect restriction hold.
\end{proposition}

\begin{proof}
This is the projectivity assertion in \Cref{thm:all-nakayama-comparison}, specialised to the periodic collection \(\Sigma_J\).
\end{proof}

\begin{corollary}[Marks--Nakayama comparison]\label{cor:marks-comparison}
Assume that \(k\) is algebraically closed.  For every \(\lambda_J\) in \eqref{eq:marks-collection},
\[
 \rho_{\lambda_J}:\DeltaB(\lambda_J)
 \xrightarrow{\sim}
 \Defsg(\lambda_J)
\]
is an equivalence before idempotent completion.  Both categories are idempotent complete.
\end{corollary}

\begin{proof}
The silting object \(T_J\) and its non-positive endomorphism dga
\[
 B_J=\tau_{\le0}\RHom_{A_d}(T_J,T_J)^{\op},
 \qquad
 \Gamma_J=H^0(B_J),
\]
are constructed in \Cref{lem:nakayama-silting,prop:nakayama-heart} below.  If the line components of \(J\) have lengths \(m_1,\ldots,m_s\), then
\[
 \Gamma_J\cong
 \prod_{a=1}^s(kA_{m_a}/\rad^2)
 \times
 \prod_{a=1}^s(kA_{m_a}/\rad^2).
\]
Put \(N=2\sum_{a=1}^s m_a\).  Then
\[
 \rad(\Gamma_J)^2=0,
 \qquad
 \Gamma_J/\rad(\Gamma_J)\cong k^N.
\]
Corti{\~n}as' Proposition~2.5.2(i)~\cite{Cortinas11}, together with the comparison in Schlichting's Section~7.1~\cite{Schlichting06}, gives
\[
 K_{-1}(\Gamma_J)
 \cong K_{-1}(k^N)
 \cong K_{-1}(k)^{\oplus N}
 =0;
\]
the middle isomorphism uses additivity, and the final vanishing uses regularity of the field and Schlichting's Example~9.5.  Since \(B_J\) is non-positive, \Cref{lem:connective-negative-k} yields
\[
 K_{-1}(\Kc_{\lambda_J})
 \cong K_{-1}(\per B_J)
 \cong K_{-1}(\Gamma_J)
 =0.
\]
Thus \(\operatorname{SObs}_0(\lambda_J)=0\) by \eqref{eq:sobs-negative-k}.  Apply \Cref{thm:negative-k-vanishing,cor:negative-k-completeness}.
\end{proof}

\subsection{The complete cone-Hom table}

The path-space calculation in this subsection is valid over an arbitrary field; algebraic closedness is needed only when the result is applied to the classification-dependent family above.

For right projectives,
\[
 \Hom_{A_d}(P_a,P_b)\cong e_bA_de_a.
\]
This is one-dimensional precisely when the unique oriented path from \(b\) to \(a\) has length less than \(d\), and is zero otherwise.

The mapping complex from \(C_i\) to \(C_j\) is
\begin{equation}\label{eq:nakayama-mapping-complex}
 \Hom(P_i,P_{j+1})
 \xrightarrow{\delta^{-1}}
 \Hom(P_{i+1},P_{j+1})\oplus\Hom(P_i,P_j)
 \xrightarrow{\delta^0}
 \Hom(P_{i+1},P_j),
\end{equation}
in degrees \(-1,0,1\).  With the middle summands ordered as displayed and with the path bases fixed by orientation,
\begin{equation}\label{eq:nakayama-differentials}
 \delta^{-1}(f)=(f\alpha_i,\alpha_jf),
 \qquad
 \delta^0(a,b)=\alpha_ja-b\alpha_i.
\end{equation}
These signs are the standard cochain mapping-complex signs.

\begin{theorem}[Cone-Hom table]\label{thm:nakayama-hom-table}
Let \(m=j-i\) in \(\mathbb Z/2d\mathbb Z\), represented by
\(0\le m\le2d-1\).  The dimensions and differential matrices of
\eqref{eq:nakayama-mapping-complex} are
\[
\begin{array}{cccc}
\toprule
 m & (\dim C^{-1},\dim C^0,\dim C^1)
   & [\delta^{-1}]&[\delta^0]\\
\midrule
 0 &(0,2,1)&0&(1\ -1)\\
 1 &(0,0,1)&0&0\\
 2\le m\le d-1 &(0,0,0)&0&0\\
 d &(1,0,0)&0&0\\
 d+1 &(1,2,0)&\binom{1}{1}&0\\
 d+2\le m\le2d-1 &(1,2,1)&\binom{1}{1}&(1\ -1).\\
\bottomrule
\end{array}
\]
Their ranks and the resulting cohomology dimensions are
\[
\begin{array}{ccc}
\toprule
 m &(\operatorname{rk}\delta^{-1},\operatorname{rk}\delta^0)
   &(\dim H^{-1},\dim H^0,\dim H^1)\\
\midrule
 0 &(0,1)&(0,1,0)\\
 1 &(0,0)&(0,0,1)\\
 2\le m\le d-1 &(0,0)&(0,0,0)\\
 d &(0,0)&(1,0,0)\\
 d+1 &(1,0)&(0,1,0)\\
 d+2\le m\le2d-1 &(1,1)&(0,0,0).\\
\bottomrule
\end{array}
\]
In particular,
\[
 H^r\RHom_{A_d}(C_i,C_j)\cong k
\]
exactly for
\[
 (j-i,r)=(0,0),(1,1),(d,-1),(d+1,0),
\]
and vanishes otherwise.
\end{theorem}

\begin{proof}
All indices are cyclic modulo \(2d\).  For each of the four Hom-spaces in \eqref{eq:nakayama-mapping-complex}, existence is decided by whether the relevant oriented path has length less than \(d\).

For \(m=0\), the degree-zero term consists of the two identities and the degree-one term is generated by \(\alpha_i\).  Formula \eqref{eq:nakayama-differentials} gives the row matrix \((1\ -1)\).

For \(m=1\), only
\(\Hom(P_{i+1},P_{i+1})\) in degree one survives.  For
\(2\le m\le d-1\), every required path wraps through at least \(d\) arrows, so all terms vanish.

For \(m=d\), the degree-minus-one term is generated by the unique path of length \(d-1\) from \(i+d+1\) to \(i\); composing it with either adjacent arrow gives a path of length \(d\), hence zero.  Thus it survives in degree \(-1\).

For \(m=d+1\), the degree-minus-one path has length \(d-2\), and both of its one-arrow extensions survive.  With the chosen path bases they are the two basis vectors in degree zero, so \(\delta^{-1}\) is \(\binom11\).  There is no degree-one term.

Finally, for \(d+2\le m\le2d-1\), all three degrees occur.  The same path-concatenation calculation gives
\(\delta^{-1}=\binom11\) and
\(\delta^0=(1\ -1)\).  Their composite is zero, each has rank one, and the complex is exact.  This proves the table, including all wrap-around cases and signs.
\end{proof}

\begin{remark}[Independent computation]\label{rem:nakayama-code}
The table is characteristic-independent: the displayed nonzero matrices have rank one over every field, including characteristic two, where \((1,-1)=(1,1)\).  It was checked by an exact-matrix script over \(\mathbb Q\) for every offset for \(d=3,4,5\), and a separate script repeats the rank and cohomology checks over \(\mathbb F_2\).  Both scripts construct the path spaces, verify \(\delta^0\delta^{-1}=0\), and compare the resulting ranks with the theorem.  All checks pass, and the scripts and captured outputs accompany the paper.
\end{remark}

\subsection{Silting generators, hearts, and cancellation}

Choose a cut of the cyclic order so that each connected component of \(J\) is a line segment.  Choose a height function
\(\ell:J\to\mathbb Z\) satisfying
\[
 \ell(j+1)=\ell(j)+1
\]
inside every line component, and extend it by
\(\ell(j+d)=\ell(j)\).  Put
\[
 V_J=J\cup(J+d),
 \qquad
 T_J=\bigoplus_{i\in V_J}C_i[\ell(i)].
\]

\begin{lemma}\label{lem:nakayama-silting}
The object \(T_J\) is a compact silting generator of the compact localisation kernel.
\end{lemma}

\begin{proof}
It is a direct sum of shifts of exactly the cone generators, hence thickly generates the compact kernel.  By \Cref{thm:nakayama-hom-table}, an unshifted nonzero class has degree \(0,1,-1,0\) in the four exceptional offsets.  After shifting, its degree is
\(r+\ell(i)-\ell(j)\).  These become respectively
\[
 0,\quad 1-1=0,\quad -1,\quad 0-1=-1.
\]
Thus there are no positive self-extensions.
\end{proof}

Let \(Q_J\) be the disjoint union of the linearly oriented line quivers on the line components of \(J\) and their \(d\)-translates.  There is an arrow \(i\to i+1\) whenever the two vertices are adjacent in one of these components.

\begin{proposition}[Heart algebra]\label{prop:nakayama-heart}
Let
\[
 B_J=\tau_{\le0}\RHom_{A_d}(T_J,T_J)^{\op},
 \qquad
 \Gamma_J=H^0(B_J).
\]
Then
\[
 \Gamma_J\cong kQ_J/\rad^2.
\]
If the line components of \(J\) have lengths
\(m_1,\ldots,m_s\), then
\[
 \Gamma_J\cong
 \prod_{a=1}^s(kA_{m_a}/\rad^2)
 \times
 \prod_{a=1}^s(kA_{m_a}/\rad^2).
\]
Moreover
\[
 \gl(kA_m/\rad^2)=m-1.
\]
\end{proposition}

\begin{proof}
The degree-zero classes after shifting are the identities and the classes at offset one.  If \(a_i:T_i\to T_{i+1}\) is such an adjacent morphism, then \(a_i\) lies in
\(e_{i+1}H^0\RHom(T_J,T_J)e_i\).  After taking the opposite in the definition of \(B_J\), its class lies in
\(e_i\Gamma_Je_{i+1}\).  Thus, with the left-to-right path convention fixed in \Cref{sec:preliminaries}, these classes give precisely the arrows \(i\to i+1\) of \(Q_J\).  A composition of two consecutive arrows would be a degree-zero class at offset two, which is absent from \Cref{thm:nakayama-hom-table}; therefore all length-two paths vanish.  No relation occurs in length zero or one.

For \(\Lambda_m=k(1\to\cdots\to m)/\rad^2\), the radical of the indecomposable left projective \(\Lambda_me_r\), for \(r>1\), is the simple at \(r-1\), while the first simple is projective.  Hence
\(\operatorname{pd}S_r=r-1\) and the global dimension is \(m-1\).
\end{proof}

\begin{lemma}[Cone-to-simple calculation]\label{lem:cone-simple}
For right \(A_d\)-modules,
\[
 H^q\RHom_{A_d}(C_i,S_j)\cong k
\]
exactly when
\[
 (j-i,q)=(0,0)\quad\text{or}\quad(1,1),
\]
and it vanishes otherwise.
\end{lemma}

\begin{proof}
The Hom complex has only
\[
 \Hom^0(C_i,S_j)=\Hom(P_i,S_j),
 \qquad
 \Hom^1(C_i,S_j)=\Hom(P_{i+1},S_j).
\]
For right projectives, \(\Hom(P_a,S_j)\cong k\) exactly when \(a=j\).  In either nonzero case the other term is zero, so no differential occurs.
\end{proof}

Define
\[
 N_J=\bigoplus_{i\in V_J}S_i[\ell(i)].
\]
The next proposition constructs a strict commuting right \(\Gamma_J\)-action on a dg representative of \(N_J\).  Its image in the derived category is characterized as follows: the idempotent at \(i\) projects to the summand \(S_i[\ell(i)]\), and the arrow \(\beta_i:i\to i+1\) acts on the right by the canonical nonzero class
\[
 S_i[\ell(i)]\longrightarrow S_{i+1}[\ell(i+1)]
\]
coming from \(\Ext^1_{A_d}(S_i,S_{i+1})\).  This prescription has radical-square-zero products in the derived category: for \(d\ge3\),
\[
 \Omega^2S_i\cong S_{i+d},
\]
so \(\Ext^2(S_i,S_{i+2})=0\).

\begin{lemma}[Single-cohomology dg bimodules]\label{lem:single-cohomology-bimodule}
Let \(B\) be a non-positive dg \(k\)-algebra, let \(\Gamma\) be an ordinary \(k\)-algebra, and let \(X\) be a left \(B\)-right \(\Gamma\) dg bimodule.  If
\[
 H^q(X)=0\qquad(q\ne0),
\]
then
\[
 X\simeq H^0(X)
 \qquad\text{in }D(B\otimes_k\Gamma^{\op}),
\]
where \(H^0(X)\) is regarded as an \(H^0(B)\)-\(\Gamma\) bimodule and restricted along the canonical dg-algebra map \(B\to H^0(B)\).
\end{lemma}

\begin{proof}
Put \(C=B\otimes_k\Gamma^{\op}\).  Since \(B\) is non-positive and \(\Gamma\) is concentrated in degree zero, \(C\) is non-positive.  The good truncation of the underlying dg \(C\)-module is therefore again a dg \(C\)-module and gives a zigzag
\[
 X\longleftarrow\tau_{\le0}X\longrightarrow H^0(X).
\]
The left arrow is a quasi-isomorphism because \(H^{>0}(X)=0\), and the right arrow is a quasi-isomorphism because \(H^{<0}(X)=0\).  On \(H^0(X)\), every element of \(B^{<0}\) acts trivially for degree reasons.  If \(b=dc\in B^0\) is a boundary and \(x\) is a degree-zero cycle, then
\[
 b[x]=[(dc)x]=[d(cx)]=0.
\]
Thus the induced left action factors through \(H^0(B)\), while the right \(\Gamma\)-action is retained.  Equivalently, \(X\) lies in the heart of the standard \(t\)-structure on \(D(C)\), whose objects are \(H^0(B)\)-\(\Gamma\) bimodules.  Consequently no additional higher operation remains to be specified.
\end{proof}

Put
\[
 F_J=T_J\Ltensor_{B_J}-:D(B_J)\longrightarrow D(A_d^{\op}),
 \qquad
 G_J=\RHom_{A_d}(T_J,-).
\]
The functor \(F_J\) is the dg Morita equivalence onto the localisation kernel generated by \(T_J\).

\begin{lemma}[Morita adjunction with coefficients]\label{lem:morita-with-coefficients}
Let \(\Gamma\) be an ordinary \(k\)-algebra.  The tensor--Hom adjunction above lifts to
\[
 F_J^\Gamma=T_J\Ltensor_{B_J}-:
 D(B_J\otimes_k\Gamma^{\op})
 \rightleftarrows
 D((A_d\otimes_k\Gamma)^{\op}):G_J^\Gamma,
\]
where \(G_J^\Gamma=\RHom_{A_d}(T_J,-)\) retains the commuting right \(\Gamma\)-action.  The coefficientwise unit and counit become the ordinary unit and counit after forgetting \(\Gamma\), and the forgetful functors detect isomorphisms.  In particular, if the underlying Morita counit is an isomorphism, so is the coefficientwise counit.
\end{lemma}

\begin{proof}
Choose the displayed bounded projective representative of \(T_J\), a cofibrant replacement on the \(B_J\)-side, and compute tensor and Hom on these models.  The usual tensor--Hom formulas preserve every commuting right \(\Gamma\)-action.  In particular, evaluation satisfies
\[
 \varepsilon\bigl((t\otimes f)\gamma\bigr)
 =(f\gamma)(t)=f(t)\gamma,
\]
so the counit is right \(\Gamma\)-linear.  Forgetting \(\Gamma\) does not change the underlying complexes; it therefore detects quasi-isomorphisms and hence is conservative on the derived categories.
\end{proof}

\begin{proposition}[Bimodule cancellation]\label{prop:nakayama-cancellation}
There is an isomorphism in \(D((A_d\otimes_k\Gamma_J)^{\op})\), the derived category of complexes with commuting right \(A_d\)- and \(\Gamma_J\)-actions,
\[
 M_J:=T_J\Ltensor_{B_J}\Gamma_J
 \xrightarrow{\sim}
 N_J.
\]
Here \(\Gamma_J\) is a left \(B_J\)-right \(\Gamma_J\) bimodule through the canonical map \(B_J\to H^0(B_J)=\Gamma_J\) on the left and the regular action on the right.
\end{proposition}

\begin{proof}
We first show that \(N_J\) lies in the localisation kernel.  For every periodic pair \(i,i+d\) there is an exact sequence
\begin{equation}\label{eq:nakayama-four-term}
 0\longrightarrow S_{i+d}
 \longrightarrow P_{i+1}
 \xrightarrow{\alpha_i}P_i
 \longrightarrow S_i\longrightarrow0.
\end{equation}
By \Cref{prop:marks-projectivity}, extension of scalars is exact in the present right-module convention.  The map \(\alpha_i\) becomes an isomorphism after localisation.  Applying extension of scalars to \eqref{eq:nakayama-four-term} therefore kills both \(S_i\) and \(S_{i+d}\).  Hence every summand of \(N_J\) belongs to the kernel.

Write
\[
 \mathsf T_i=C_i[\ell(i)],
 \qquad
 \mathsf N_i=S_i[\ell(i)].
\]
For each \(i\), let
\[
 u_i\in H^0\RHom_{A_d}(\mathsf T_i,\mathsf N_i)
\]
be the class of the shifted quotient \(C_i\to S_i\).  If \(i,i+1\) are adjacent in a line component, let
\[
 v_i\in H^0\RHom_{A_d}(\mathsf T_i,\mathsf N_{i+1})
\]
be the class represented before shifting by the quotient
\(P_{i+1}\twoheadrightarrow S_{i+1}\) in degree one.  For fixed \(r\), \Cref{lem:cone-simple} shows that \(G_J(\mathsf N_r)\) has cohomology only in degree zero, with basis \(u_r\) and, when \(r-1,r\) are adjacent, \(v_{r-1}\).  Precomposition identifies this left \(\Gamma_J\)-module with \(\Gamma_Je_r\).  More precisely, \Cref{lem:single-cohomology-bimodule} with right algebra \(k\) gives an isomorphism
\[
 \theta_r:\Gamma_Je_r\xrightarrow{\sim}G_J(\mathsf N_r)
 \quad\text{in }D(B_J)
\]
whose map on cohomology sends \(e_r\) to \(u_r\) and, when present, the incoming arrow \(\beta_{r-1}\) to \(v_{r-1}\).  Since \(\mathsf N_r\) lies in the kernel, the Morita counit is an isomorphism, and
\[
 \varphi_r:=\varepsilon_{\mathsf N_r}\circ F_J(\theta_r):
 F_J(\Gamma_Je_r)\xrightarrow{\sim}\mathsf N_r
\]
gives
\begin{equation}\label{eq:nakayama-column-cancellation}
 F_J(\Gamma_Je_r)\simeq\mathsf N_r
 \qquad\text{in }D(A_d^{\op}).
\end{equation}

The regular bimodule \(\Gamma_J\) has a strict right \(\Gamma_J\)-action, so
\(M_J=F_J^{\Gamma_J}(\Gamma_J)\) is a strict dg
\((A_d^{\op},\Gamma_J)\)-bimodule.  Its underlying right \(A_d\)-object decomposes by columns, and
\(\varphi=\bigoplus_r\varphi_r\) identifies it with
\(\bigoplus_r\mathsf N_r=N_J\).  We use this very dg bimodule, together with \(\varphi\), as the strict coefficient object denoted by \(N_J\) for the remainder of the proof; its underlying right \(A_d\)-object is the displayed direct sum.  No general strictification assertion is being invoked.

To identify its action, let
\(R_{\beta_i}:\Gamma_Je_i\to\Gamma_Je_{i+1}\) be right multiplication by \(\beta_i\), and denote by \(\nu_i:\mathsf N_i\to\mathsf N_{i+1}\) the induced map under the column identifications.  Naturality of the Morita counit and the triangle identities give a commutative square in \(D(B_J)\):
\[
 \begin{array}{ccc}
 \Gamma_Je_i & \xrightarrow{R_{\beta_i}} & \Gamma_Je_{i+1}\\
 \big\downarrow\rlap{\(\scriptstyle\theta_i\)} &&
 \big\downarrow\rlap{\(\scriptstyle\theta_{i+1}\)}\\
 G_J(\mathsf N_i) & \xrightarrow{G_J(\nu_i)} & G_J(\mathsf N_{i+1}).
 \end{array}
\]
Since \(R_{\beta_i}(e_i)=\beta_i\), the map \(G_J(\nu_i)\) sends \(u_i\) to \(v_i\).  Thus \(\nu_i\) is precisely the prescribed nonzero Ext class.  Products of two arrow maps are strictly zero on this representative because they are induced by products in \(\rad^2\Gamma_J=0\).

We now verify the full bimodule identification explicitly.  Let
\[
 a_i\in H^0\RHom_{A_d}(\mathsf T_i,\mathsf T_{i+1})
\]
be the normalized adjacent class and let
\(\beta_i=a_i^{\op}\in e_i\Gamma_Je_{i+1}\).  The \(u_i\)'s and \(v_i\)'s form a \(k\)-basis of \(H^0G_J^{\Gamma_J}(N_J)\).  Precomposition on the left and the strict action on the right give
\[
 \beta_i u_{i+1}=v_i=u_i\beta_i.
\]
The vertex idempotents act by
\[
 e_ru_i=\delta_{ri}u_i,
 \qquad
 u_ie_r=\delta_{ir}u_i,
\]
and
\[
 e_rv_i=\delta_{ri}v_i,
 \qquad
 v_ie_r=\delta_{r,i+1}v_i.
\]
Every product of a radical generator with a \(v_i\) is zero.  On the left this is the absence of a degree-zero cone--Hom class at offset two; on the right it is the strict relation \(\rad^2\Gamma_J=0\).  Hence
\[
 \Phi:\Gamma_J\longrightarrow H^0G_J^{\Gamma_J}(N_J),
 \qquad
 e_i\longmapsto u_i,
 \quad
 \beta_i\longmapsto v_i,
\]
equivalently \(1\mapsto\sum_i u_i\), is an isomorphism of
\(\Gamma_J\)-\(\Gamma_J\) bimodules.  By
\Cref{lem:single-cohomology-bimodule},
\[
 G_J^{\Gamma_J}(N_J)\simeq H^0G_J^{\Gamma_J}(N_J)\simeq\Gamma_J
 \qquad\text{in }D(B_J\otimes_k\Gamma_J^{\op}).
\]
No hidden higher operation remains because this bimodule lies in the heart of the standard \(t\)-structure.  Since \(N_J\) lies in the kernel, \Cref{lem:morita-with-coefficients} shows that the coefficientwise Morita counit is an isomorphism.  Applying \(F_J^{\Gamma_J}\) in the coefficient category gives
\[
 M_J=F_J^{\Gamma_J}(\Gamma_J)
 \simeq F_J^{\Gamma_J}G_J^{\Gamma_J}(N_J)
 \xrightarrow{\ \varepsilon_{N_J}\ }N_J,
\]
and every arrow in this display is right \(\Gamma_J\)-linear.
\end{proof}

\begin{theorem}[Intrinsic Nakayama defects]\label{thm:nakayama-intrinsic}
Assume that \(k\) is algebraically closed.  For every Marks-admissible \(J\), the object \(T_J\) is \(A_d\)-bounded and
\[
 \DeltaB(\lambda_J)
 \simeq
 \Dfd(B_J)/\per(B_J).
\]
If \(J\ne\varnothing\), this category is nonzero.  Moreover
\[
 \Defsg(\lambda_J)
 =\thick_{\Dsg(A_d)}\{S_i\mid i\in V_J\},
\]
and
\[
 \Kar\bigl(\Dfd(B_J)/\per(B_J)\bigr)
 \xrightarrow{\sim}
 \Defsg(\lambda_J).
\]
\end{theorem}

\begin{proof}
By \Cref{prop:nakayama-heart}, \(\Gamma_J\) has finite global dimension.  By \Cref{prop:nakayama-cancellation}, \(M_J\) is a bounded complex.  If \(g=\gl\Gamma_J\), every finite-dimensional left \(\Gamma_J\)-module has a projective resolution of length at most \(g\).  Tensoring such a resolution with the fixed bounded complex \(M_J\) gives common lower and upper cohomological bounds.  Hence \(M_J\) has fixed finite Tor-amplitude.  The heart-amplitude theorem gives the intrinsic formula.

To prove nonvanishing, choose \(i\in V_J\).  The finite-dimensional left \(B_J\)-module \(\Gamma_Je_i\) lies in \(\Dfd(B_J)\), and
\[
 F_J(\Gamma_Je_i)=M_Je_i\cong S_i[\ell(i)].
\]
If \(\Gamma_Je_i\) were perfect over \(B_J\), its image would lie in
\(\thick(T_J)=\Kc_{\lambda_J}\subseteq\per(A_d)\).  But \(A_d\) is self-injective and \(S_i\) is nonprojective; a finitely generated module over a self-injective algebra has finite projective dimension only when it is projective.  Thus \(S_i\notin\per(A_d)\), a contradiction.

Because \(B_J\) is non-positive,
\(\Dfd(B_J)=\thick(\Gamma_J\text{-}\modcat)\).  The finite global dimension of \(\Gamma_J\) further gives
\(\Gamma_J\text{-}\modcat\subseteq\thick(\Gamma_J)\), and hence
\(\Dfd(B_J)=\thick(\Gamma_J)\).  Therefore the equality describing \(\Defsg\) follows from
\[
 \Kb(\lambda_J)=F_J(\Dfd(B_J))
 =\thick_{\Db(A_d)}(M_J)
 =\thick_{\Db(A_d)}\{S_i\mid i\in V_J\},
\]
together with the definition of \(\Defsg\).  Finally apply
\Cref{cor:marks-comparison}.
\end{proof}

\subsection{The complete \texorpdfstring{\(\widetilde A^3_6\)}{A63} comparison}\label{subsec:a63}

Under the standing algebraically closed-field convention, let \(A=A_3=\widetilde A^3_6\) and localise at
\(\Sigma=\{S_1,S_4\}\).  Marks identifies the target, up to Morita equivalence, with \(\widetilde A^2_4\) \cite[Example~5.13]{Marks15}.  Put
\[
 C_1=[P_2\xrightarrow{\alpha_1}P_1],
 \qquad
 C_4=[P_5\xrightarrow{\alpha_4}P_4],
 \qquad
 T=C_1\oplus C_4.
\]

\begin{proposition}[The kernel dga and its minimal resolutions]\label{prop:a63-dga}
The object \(T\) is compact silting.  Its non-positive dg endomorphism algebra is the dg algebra \(B\) with
\[
 B^0=ke_1\oplus ke_4,
 \qquad
 B^{-1}=ka\oplus kb,
 \qquad
 d=0,
\]
where
\[
 a\in e_4Be_1,
 \qquad
 b\in e_1Be_4,
\]
and all products of two negative arrows vanish.  Thus
\(\Gamma=H^0(B)=k\times k\).

If \(L_1,L_4\) are the simple left \(B\)-modules, then
\[
 0\to L_4[1]\to Be_1\to L_1\to0,
 \qquad
 0\to L_1[1]\to Be_4\to L_4\to0.
\]
Their totalised minimal semi-free dg resolutions have underlying graded modules
\[
 P_1^{\#}=Be_1\oplus Be_4[2]\oplus Be_1[4]\oplus Be_4[6]\oplus\cdots,
\]
\[
 P_4^{\#}=Be_4\oplus Be_1[2]\oplus Be_4[4]\oplus Be_1[6]\oplus\cdots.
\]
For \(P_1\), let \(x_n\) denote the generator of its \(n\)-th summand, of degree \(-2n\). The differential is
\[
 dx_0=0,\qquad dx_{2r+1}=a x_{2r},\qquad
 dx_{2r+2}=b x_{2r+1}\quad(r\ge0),
\]
extended by the dg module rule. For \(P_4\), interchange \(e_1,e_4\) and \(a,b\). These are totalised dg modules: the shifts include both the internal degree and the resolution degree.
In particular \(L_1,L_4\notin\per(B)\).
\end{proposition}

\begin{proof}
The case \(d=3\) of \Cref{thm:nakayama-hom-table} gives
\[
 H^*\RHom(C_i,C_i)=k[0]
 \quad(i=1,4),
\]
and one cross class in degree \(-1\) in each direction.  There are no positive classes, so \(T\) is silting.  Since no degree \(-2\) term exists, all products of the two negative classes vanish, and the displayed dg algebra is a strict model of the good truncation.

The radical of \(Be_1\) is generated by the degree \(-1\) arrow \(a\) and is isomorphic to \(L_4[1]\); similarly
\(\rad(Be_4)\cong L_1[1]\).  The two short exact sequences follow.  Each displayed differential has degree one and its square is zero because \(ab=ba=0\). Every radical basis element is the image of the next generator; the augmentation therefore has cohomology exactly the indicated simple in degree zero. Filtering by the summands gives a semi-free filtration. The resolutions are minimal because every differential lands in the graded radical.  Tensoring with \(\Gamma=B/(a,b)\) kills all differentials and leaves infinitely many nonzero graded Betti terms, so neither simple is perfect.
\end{proof}

\begin{proposition}[The bounded kernel]\label{prop:a63-bounded-kernel}
The heart bimodule is
\[
 M_T=T\Ltensor_B\Gamma\cong S_1\oplus S_4.
\]
Consequently
\[
 \Kb(\lambda)=\thick_{\Db(A)}(S_1,S_4),
 \qquad
 \DeltaB(\lambda)\simeq\Dfd(B)/\per(B)\ne0.
\]
\end{proposition}

\begin{proof}
This is the unshifted case of \Cref{prop:nakayama-cancellation}.  Since
\(\Gamma=k\times k\) is semisimple, \(M_T\) has Tor-amplitude concentrated in the cohomological range occupied by \(S_1\oplus S_4\), and the heart-amplitude theorem applies.  The standard heart modules generate \(\Dfd(B)\), so their images \(S_1,S_4\) generate the bounded kernel.  Nonvanishing also follows from \Cref{prop:a63-dga}.
\end{proof}

Write \(M(i,\ell)=P_i/\rad^\ell P_i\).  Since \(A\) is self-injective, the canonical functor identifies
\[
 \Dsg(A)\simeq\underline{\modcat}\text{-}A,
\]
the stable module category; see Buchweitz, Theorem~4.4.1~\cite{Buchweitz}.

\begin{lemma}[Stable kernel calculation]\label{lem:a63-stable}
Let
\[
 X_0=S_1,
 \qquad
 X_1=M(5,2),
 \qquad
 X_2=S_4,
 \qquad
 X_3=M(2,2).
\]
Then
\[
 X_{r+1}\cong X_r[1]
 \quad(r\in\mathbb Z/4\mathbb Z),
\]
and
\[
 \dim_k\underline{\Hom}_A(X_r,X_s[q])
 =
 \begin{cases}
 1,&s+q-r\equiv0\pmod4,\\
 0,&\text{otherwise}.
 \end{cases}
\]
The thick subcategory generated by \(S_1,S_4\) is
\[
 \add(X_0,X_1,X_2,X_3),
\]
and every distinguished triangle in it is isomorphic to a finite direct sum of trivial triangles.
\end{lemma}

\begin{proof}
For a nonprojective uniserial module,
\[
 \Omega M(i,\ell)\cong M(i+\ell,3-\ell).
\]
Solving \(\Omega Y=X\) gives
\[
 S_1[1]\cong M(5,2),
 \quad
 M(5,2)[1]\cong S_4,
 \quad
 S_4[1]\cong M(2,2),
 \quad
 M(2,2)[1]\cong S_1.
\]

A direct calculation with the four uniserial representations gives
\[
 \Hom_A(X_r,X_s)=
 \begin{cases}
 k,&r=s,\\
 0,&r\ne s.
 \end{cases}
\]
The nonzero diagonal endomorphisms do not factor through projectives, so the same formula holds stably.  Replacing \(X_s\) by its \(q\)-fold suspension yields the displayed congruence formula.

The additive subcategory generated by the four modules is closed under suspension.  After decomposing source and target into the four isotypic parts, every morphism is a direct sum of matrices over \(k\), one matrix for each \(X_r\), because all off-diagonal stable Hom-spaces vanish.  Independent changes of bases in source and target reduce each matrix to \(\operatorname{diag}(I_t,0)\).  Its cone is therefore a direct sum of cones of isomorphisms and zero maps, hence belongs to the same additive category.  Thus it is thick and contains \(S_1,S_4\).  Conversely it is contained in their thick closure.  The same normal form shows that every triangle is isomorphic to a finite direct sum of trivial triangles.
\end{proof}

By \Cref{cor:marks-comparison}, pre-completion essential surjectivity and idempotent completeness are already known for every Marks-admissible localisation; the next theorem identifies the first member explicitly as the split four-periodic category.

\begin{theorem}[Complete first comparison]\label{thm:a63-comparison}
Let \(k\) be algebraically closed.  For the localisation of \(\widetilde A^3_6\) at \(\{S_1,S_4\}\), whose target is Morita equivalent to \(\widetilde A^2_4\), the comparison functor is an equivalence before idempotent completion
\[
 \rho_\lambda:\DeltaB(\lambda)
 \xrightarrow{\sim}
 \Defsg(\lambda).
\]
Both categories are equivalent to the split four-periodic triangulated category
\[
 \vect_k^{\mathbb Z/4}
\]
of finite-dimensional \(\mathbb Z/4\)-graded vector spaces, with suspension given by cyclic degree shift.  In \(\DeltaB(\lambda)\), let
\[
 Y_0=S_1,
 \qquad Y_1=S_1[1],
 \qquad Y_2=S_4,
 \qquad Y_3=S_4[1].
\]
Their images in \(\Dsg(A)\) are represented respectively by
\[
 S_1,
 \quad M(5,2),
 \quad S_4,
 \quad M(2,2).
\]
The four objects \(Y_r\) are the graded simples under the displayed equivalence, and
\[
 \dim_k\Hom_{\DeltaB(\lambda)}(Y_r,Y_s[q])
 =\dim_k\Hom_{\Defsg(\lambda)}(\rho Y_r,\rho Y_s[q])
 =\begin{cases}
 1,&s+q-r\equiv0\pmod4,\\
 0,&\text{otherwise}.
 \end{cases}
\]
\end{theorem}

\begin{proof}
By \Cref{prop:marks-projectivity} and
\Cref{thm:unconditional-comparison}, \(\rho_\lambda\) is fully faithful.  By
\Cref{prop:a63-bounded-kernel}, the actual singularity kernel is the thick stable subcategory generated by \(S_1,S_4\).  By
\Cref{lem:a63-stable}, every object of this target is a finite direct sum of the four suspension-orbit objects, all of which lie in the triangulated image of \(\DeltaB(\lambda)\).  Hence the comparison is essentially surjective before idempotent completion.  The same lemma gives the morphism formula and identifies the split triangulated structure with \(\vect_k^{\mathbb Z/4}\).
\end{proof}



\section{The actual singularity kernel}\label{sec:actual-kernel}

Throughout this section, let \(\lambda:R\to S\) be a homological ring epimorphism satisfying \textup{(B1)--(B3)}.  In particular, the results apply to stably flat universal localisations satisfying these conditions.  Put
\[
 \cT=\Db(R\text{-}\modcat),\quad
 \cT_S=\Db(S\text{-}\modcat),\quad
 \cP=\per(R),\quad
 \cP_S=\per(S),\quad
 \cK=\Kb(\lambda).
\]
Let \(\Phi=S\Ltensor_R-\) and let \(U\) be restriction of scalars.

The boundedness assertion in the next lemma uses only \textup{(B1)--(B3)}.  Homologicality enters later through the fully faithful derived restriction functor, the localisation unit triangles, and the invertible counit.

\begin{lemma}[Bounded derived adjunction]\label{lem:bounded-adjunction}
Under \textup{(B1)--(B3)}, extension and restriction of scalars restrict to an adjunction
\[
 \Phi=S\Ltensor_R-:\Db(R\text{-}\modcat)
 \rightleftarrows
 \Db(S\text{-}\modcat):U.
\]
For every object in either bounded category, the corresponding unit or counit, and every triangle completing that morphism, lie in the asserted bounded categories.
\end{lemma}

\begin{proof}
Let \(M\) be a finitely generated left \(R\)-module.  Since \(R\) is left noetherian, \(M\) has a projective resolution by finitely generated projective left \(R\)-modules.  Tensoring with \(S\) gives a complex of finitely generated projective left \(S\)-modules.  Since \(S\) is left noetherian, every cohomology module of this complex is finitely generated, and \textup{(B2)} makes those modules vanish outside a finite range.  Hence \(\Phi M\in\Db(S\text{-}\modcat)\).  Every object of \(\Db(R\text{-}\modcat)\) is assembled from its finitely many cohomology modules by standard truncation triangles, so exactness gives \(\Phi X\in\Db(S\text{-}\modcat)\) for every such \(X\).

Restriction is exact, so \(H^n(UY)=U(H^nY)\).  Applying \textup{(B3)} to finitely generated \(S\)-modules concentrated in degree zero shows that these restricted modules are finitely generated over \(R\); hence \(UY\in\Db(R\text{-}\modcat)\).  The two bounded categories are full triangulated subcategories of their unbounded derived categories, so the unbounded adjunction restricts.  The endpoints of every unit and counit lie in the stated bounded category, and therefore so does its cone.
\end{proof}

\begin{lemma}[The compact kernel]\label{lem:compact-intersection}
By definition,
\[
 \cK\cap\cP=\Kc_\lambda.
\]
If, in addition, \(\lambda\) is a stably flat universal localisation at a set \(\Sigma\) of maps between finitely generated projective modules, then the substantive compact-generation identification is
\[
 \Kc_\lambda=\thick\langle\Cone(\Sigma)\rangle.
\]
\end{lemma}

\begin{proof}
Since \(\per(R)\subseteq\Db(R\text{-}\modcat)\), the definitions give
\[
 \cK\cap\cP
 =\bigl(\cK_\lambda\cap\Db(R\text{-}\modcat)\bigr)\cap\per(R)
 =\cK_\lambda\cap\per(R)
 =\Kc_\lambda.
\]
For the additional assertion, stable flatness gives
\(\cK_\lambda=\Loc\langle\Cone(\Sigma)\rangle\), and every displayed cone is compact.  The Thomason--Neeman theorem therefore identifies the compact objects of this localising subcategory with \(\thick\langle\Cone(\Sigma)\rangle\), which is precisely \(\Kc_\lambda\).
\end{proof}

\begin{lemma}\label{lem:preimage-perfect}
Let
\[
 \cL=\Phi^{-1}(\cP_S)
 =\{X\in\cT\mid \Phi X\in\per(S)\}.
\]
Then
\[
 \cL=\thick(\cP,\cK).
\]
\end{lemma}

\begin{proof}
The inclusion \(\thick(\cP,\cK)\subseteq\cL\) follows from
\(\Phi(\cP)\subseteq\cP_S\) and \(\Phi(\cK)=0\).

For the reverse inclusion, let \(X\in\cL\), and write \(Y=\Phi X\in\cP_S\).  By (B3), \(U(Y)\in\cT\).  Since \(Y\in\thick(S)\), one has
\(U(Y)\in\thick(U(S))\).

The localisation unit at \(R\) fits into a triangle
\[
 K_R\longrightarrow R\longrightarrow U(S)\longrightarrow K_R[1].
\]
The counit \(\Phi U\to\id_{D(S)}\) is an isomorphism, so applying \(\Phi\) shows \(\Phi K_R=0\).  By \textup{(B3)}, \(U(S)\in\cT\); since \(R\in\cT\), the triangle gives \(K_R\in\cT\), and therefore \(K_R\in\cK\).  It follows that
\(U(S)\in\thick(\cP,\cK)\), hence also \(U(Y)\in\thick(\cP,\cK)\).

Finally, the unit triangle at \(X\) is
\[
 K_X\longrightarrow X\longrightarrow U\Phi X\longrightarrow K_X[1].
\]
Here \(K_X\in\cK\), while \(U\Phi X=U(Y)\) lies in
\(\thick(\cP,\cK)\).  Thus \(X\) lies there as well.
\end{proof}

We use a standard elementary observation which is sometimes hidden in quotient notation.

\begin{lemma}\label{lem:subcategory-quotient-embedding}
Let \(\cP\subseteq\cL\subseteq\cT\) be full triangulated subcategories, with \(\cP\) thick in \(\cT\).  Then the induced functor
\[
 \cL/\cP\longrightarrow\cT/\cP
\]
is fully faithful.
\end{lemma}

\begin{proof}
A morphism in \(\cT/\cP\) between objects \(L_1,L_2\in\cL\) is represented by a roof
\[
 L_1\xleftarrow{s}Z\longrightarrow L_2,
 \qquad \Cone(s)\in\cP.
\]
Completing \(s\) to a triangle shows that \(Z\in\cL\), because both
\(L_1\) and \(\Cone(s)\) lie in \(\cL\).  Thus every roof and every common refinement may be taken inside \(\cL\), which proves fullness and faithfulness.
\end{proof}

By \Cref{lem:preimage-perfect}, the singularity defect of \Cref{def:singularity-defect} has the concrete form
\[
 \Defsg(\lambda)=\cL/\cP
\]
for a homological ring epimorphism satisfying \textup{(B1)--(B3)}.

\begin{theorem}[Actual singularity-kernel theorem]\label{thm:actual-kernel}
Let \(\lambda:R\to S\) be a homological ring epimorphism satisfying \textup{(B1)--(B3)}.  Then extension of scalars induces an exact functor
\[
 \overline\Phi:\Dsg(R)\longrightarrow\Dsg(S)
\]
whose object-kernel is identified with the fully faithful subcategory \(\Defsg(\lambda)\):
\[
 \kerob(\overline\Phi)=\Defsg(\lambda)\subseteq\Dsg(R).
\]
Moreover there is a Verdier quotient equivalence
\[
 \Dsg(R)/\Defsg(\lambda)\xrightarrow{\sim}\Dsg(S).
\]
No idempotent completion is required in either assertion.
\end{theorem}

\begin{proof}
Since \(\Phi(\per(R))\subseteq\per(S)\), the functor descends to singularity categories.  An object \(X\in\cT\) maps to zero precisely when
\(\Phi X\in\cP_S\), that is, precisely when \(X\in\cL\).  By
\Cref{lem:subcategory-quotient-embedding}, the object-kernel in the sense of \Cref{def:object-kernel} is the fully faithful quotient \(\cL/\cP=\Defsg(\lambda)\).

The third isomorphism theorem for Verdier quotients gives
\[
 (\cT/\cP)/(\cL/\cP)\simeq\cT/\cL.
\]
Because \(\cK\subseteq\cL\), another application gives
\[
 \cT/\cL\simeq (\cT/\cK)/(\cL/\cK).
\]
Under the bounded localisation equivalence of \Cref{thm:appendix-bounded-localisation}, the subcategory \(\cL/\cK\) is exactly \(\cP_S\): by definition, it consists of the objects whose images are perfect.  Hence
\[
 \cT/\cL\simeq\cT_S/\cP_S=\Dsg(S).
\]
All these are ordinary Verdier quotient equivalences.  Idempotent completion enters only later, when \(\DeltaB(\lambda)\) is compared with \(\Defsg(\lambda)\).
\end{proof}

\begin{remark}\label{rem:prior-localization}
General localisation theorems for singularity categories are known in broader recollement settings, notably for proper dg algebras and, more recently, for weakly approximable triangulated categories; see \cite{JinYangZhou,SunZhangZhang} and the comparison in \Cref{sec:literature}.  The point of \Cref{thm:actual-kernel} is the explicit formula in terms of the bounded kernel of a homological universal localisation.  The next section derives the comparison with the compact-kernel quotient from classical compact factorisation.
\end{remark}


\section{The classical compact-kernel comparison}\label{sec:rigidity}
The following is a consequence of compact factorisation and the localisation-of-subcategories criterion. No algebraicity or dg enhancement is needed.

\begin{theorem}[Classical compact-kernel comparison]\label{thm:abstract-rigidity}
Let \(\cT\) be a compactly generated triangulated category, let \(\cP=\cT^c\), and let \(\cC\subseteq\cP\) be thick and essentially small. Put \(\cL=\Loc(\cC)\). If \(\cC\subseteq\cK\subseteq\cL\), with \(\cK\) thick and essentially small, and \(\cA=\thick(\cP,\cK)\), then
\[
 \rho:\cK/\cC\longrightarrow\cA/\cP
\]
is fully faithful and induces an equivalence after idempotent completion.
\end{theorem}

\begin{proof}
Compact generation gives \(\cL\cap\cP=\cC\), hence \(\cK\cap\cP=\cC\). Moreover every map \(P_0\to X\), with \(P_0\in\cP\) and \(X\in\cL\), factors through an object of \(\cC\). This is compact factorisation: apply \cite[Theorem~7.2.1(1)]{Krause10} at \(\alpha=\aleph_0\), or the compact-generation theorem in \cite[Section~4.4]{Neeman01}.

Apply \cite[Lemma~4.7.1(1)]{Krause10} inside \(\cA\), with denominator \(\cP\) and subcategory \(\cK\). Its hypothesis is precisely that every map from an object of \(\cP\) to an object of \(\cK\) factors through \(\cP\cap\cK=\cC\), as just proved. It follows that \(\cK/\cC\to\cA/\cP\) is fully faithful. Its essential image is triangulated and thickly generates the target, because \(\cA=\thick(\cP,\cK)\) and \(\cP\) becomes zero. Passing to idempotent completions therefore makes it essentially surjective as well.
\end{proof}

The factorisation in this proof has source compact and target in the kernel. The other orientation, recorded in \Cref{sec:comparison}, is useful under one-sided perfection hypotheses but is not needed here. In the algebraic setting, \Cref{sec:obstruction} gives Efimov's enhanced Calkin explanation and the vanishing of the derived factorisation obstruction.

\begin{theorem}[Smashing-localisation comparison]\label{thm:smashing-comparison}
Let \(F:\cT\to\mathcal U\) be a smashing localisation of a compactly generated triangulated category.  Suppose its kernel \(\cL\) is generated, as a localising subcategory, by a set of compact objects of \(\cT\).  Put
\[
 \cC=\cL\cap\cT^c.
\]
For every essentially small thick subcategory \(\cK\subseteq\cL\) containing \(\cC\), the canonical functor
\[
 \cK/\cC\longrightarrow
 \thick(\cT^c,\cK)/\cT^c
\]
is fully faithful and induces an equivalence after idempotent completion.
\end{theorem}

\begin{proof}
Choose a set \(\mathcal G\subseteq\cT^c\) which generates \(\cL\).  The Thomason--Neeman compact-generation theorem~\cite[Theorem~4.4.9]{Neeman01} identifies
\[
 \cL^c=\cL\cap\cT^c=\thick(\mathcal G)=\cC
\]
and gives \(\cL=\Loc(\cC)\).  Thus every hypothesis of \Cref{thm:abstract-rigidity} is satisfied with \(\cP=\cT^c\) and \(\cA=\thick(\cT^c,\cK)\).  Applying that theorem gives the asserted full faithfulness and Karoubi equivalence.
\end{proof}

\begin{corollary}[Homological ring epimorphisms with compactly generated kernel]\label{cor:homological-epi-compact-kernel}
Let \(R\to S\) be a homological ring epimorphism.  Assume its kernel \(\cL\subseteq D(R)\) is generated by compact \(R\)-complexes.  Then the conclusion of \Cref{thm:smashing-comparison} holds for every essentially small thick \(\cK\subseteq\cL\) containing \(\cL\cap\per R\).  In particular, when \(R\) is left noetherian and
\[
 \cK=\cL\cap\Db(R\text{-}\modcat),
\]
the bounded compact-kernel defect compares fully faithfully with its ambient singular defect and is an equivalence after idempotent completion.
\end{corollary}

\begin{proof}
A homological ring epimorphism has fully faithful derived restriction, so derived extension is a smashing localisation.  The compact-generation hypotheses are those of \Cref{thm:smashing-comparison}; the final assertion takes the displayed bounded thick subcategory as \(\cK\).
\end{proof}

\begin{theorem}[Comparison without one-sided perfection hypotheses]\label{thm:unconditional-comparison}
Let \(R\) be left noetherian and let \(\lambda:R\to S=R_\Sigma\) be a stably flat universal localisation at an arbitrary set of maps between finitely generated projective left modules.  Then
\[
 \rho_\lambda:\DeltaB(\lambda)\longrightarrow\Defsg(\lambda)
\]
is fully faithful,
\[
 \Kar\DeltaB(\lambda)\xrightarrow{\sim}\Kar\Defsg(\lambda),
\]
and \(\Obs(\lambda)\simeq0\).
\end{theorem}

\begin{proof}
Take \(\cT=D(R)\), \(\cP=\per R\), \(\cC=\Kc_\lambda\), and \(\cK=\Kb(\lambda)\).  Stable flatness gives
\[
 \cK_\lambda=\Loc\langle\Cone(\Sigma)\rangle=\Loc(\cC),
\]
and the compact part of this localising category is \(\cC\).  Hence every hypothesis of \Cref{thm:abstract-rigidity} holds and its comparison is \(\rho_\lambda\).
\end{proof}

\begin{corollary}[Singularity-kernel dictionary]\label{cor:unconditional-kernel}
Under the hypotheses of \Cref{thm:unconditional-comparison}, assume in addition \textup{(B1)--(B3)}.  Then \(\rho_\lambda\) identifies \(\DeltaB(\lambda)\) fully faithfully with \(\kerob(\Dsg(R)\to\Dsg(S))\), up to direct summands.  Before idempotent completion, the quotient of \(\Dsg(R)\) by this object-kernel is \(\Dsg(S)\).
\end{corollary}

\begin{proof}
Combine \Cref{thm:unconditional-comparison} with the actual-kernel theorem in \Cref{sec:actual-kernel}.
\end{proof}

\begin{remark}\label{cor:no-beta}
In particular, full faithfulness cannot fail for a finite stably flat universal localisation. The supplemental AHKLY calculations check the concrete one-object duality maps of the enhanced appendix.
\end{remark}


\subsection{What can remain before completion}
Write \(\cE\) for the common idempotent completion, and let \(H_\Delta\subseteq H_{\mathrm{Def}}\subseteq K_0(\cE)\) be the images of the two quotient \(K_0\)-groups. The possible failure of essential surjectivity is measured by
\[
 \operatorname{SObs}_0(\lambda):=H_{\mathrm{Def}}/H_\Delta
 \cong\ker\!\left(K_{-1}(\Kc_\lambda)
 \longrightarrow K_{-1}(\Kb(\lambda))\oplus K_{-1}(R)\right).
\]
The proof compares two Schlichting localisation sequences and is given in \Cref{thm:negative-k-obstruction}. This group vanishes throughout the finite-dimensional silting range (\Cref{cor:finite-dimensional-precompletion}), including all the constructions of this paper. No example with nonzero summand obstruction is known from the present arguments; the precise existence question is \Cref{question:negative-k-counterexample}.

\section{Relation to prior work}\label{sec:literature}

\subsection{Prescribed dg kernels}
The construction in \Cref{thm:realisation} belongs to the established setting of matrix-ring universal localisation. Sheiham \cite[Theorem~2.4 and Example~2.5(1)]{Sheiham06} gives explicit localisations of triangular matrix rings at a map between column projectives, including the passage from an upper triangular ring to a full matrix ring. Morita contexts and exact-context localisations are studied in \cite{GreenPsaroudakis14,GaoPsaroudakis17,ChenXiRecollI,ChenXi19}. The claim here is the prescribed identification of the kernel dg algebra, by a strict dg map, together with two-sided target projectivity, t-exactness, and full profile. The matrix-ring mechanism by itself does not assert that package.

Marks--Vit\'oria \cite[Theorem~3.3]{MarksVitoria15} prove that a ring epimorphism whose target is finitely presented as a left source module of projective dimension at most one is homological if and only if it is a universal localisation. In particular, their recognition theorem applies to the projective targets constructed here. Our t-exact converse instead imposes a condition on a marked compact kernel generator and forces that generator to be two-term; it does not assume a projective-dimension bound on the target. Its universal-localisation conclusion should be read alongside this earlier recognition result.

Minimal \(A_\infty\)-models and their uniqueness are standard \cite{Keller01Ainf,LefevreHasegawa03}; crossed extensions are classified by \(HH^3\) in \cite{BauesMinian02}. Finite strict models for proper connective dg algebras over arbitrary fields are supplied by Raedschelders--Stevenson's appendix to Goodbody \cite[Theorem~A.3]{Goodbody24}. Accordingly, \Cref{thm:finite-crossed} is an explicit finite construction from a Hochschild class, rather than a new finite-model existence theorem. The ordinary-ring kernel realisation, including its non-formal example, is the separate assertion.

Kalck--Yang \cite{KalckYangI,KalckYangII} and Booth \cite{Booth21} describe derived quotients associated with idempotents. Their Schur-functor recollements are a different realisation setting from extension along an ordinary ring epimorphism. The compact projection in \Cref{prop:left-projection} and the negative Calabi--Yau argument isolate a boundary between these settings; \Cref{ex:sphere-schur} makes it explicit. No general converse for Schur kernels is proved here.

\subsection{Detectors and changing profiles}
The finite-total-Hom characterisation is recorded in \cite[Lemma~4.1]{GaoPsaroudakis18}; intrinsic homological finiteness is further developed in \cite{KuznetsovShinder25}. Far-away orthogonality and Brown--Comenetz detection provide the framework of \cite{KostasPsaroudakisVitoria25}. The canonical detector in \Cref{thm:detector} applies the familiar duality \(\RHom_R(X,DR)\simeq D_kX\) to the embedded kernel. Its purpose is to identify the additional data determining ambient boundedness.

One-point extensions are classical \cite{Ringel84}, and lifting recollements along them is studied in \cite{LinLin08}. \Cref{thm:ambient} computes the extended universal localisation and its unchanged marked kernel, proves the all-degree stable-flatness criterion, and identifies the new profile as an intersection of detector cuts. The degree-zero concentration test comes from \cite{NicolasSaorin09}. Chen's triangular singularity equivalence \cite[Theorem~4.1(1)]{ChenSchur09} supplies the ambient singularity category for the ordinary-kernel examples. Marks' Nakayama classification \cite{Marks15} and the relation of universal localisation to silting \cite{MarksStovicek19} are used with the stated module-side and field hypotheses.

\subsection{Classical comparison and its refinements}
The triangulated compact-kernel comparison is a classical consequence, not a new theorem requiring dg methods. Krause's compact factorisation theorem \cite[Theorem~7.2.1(1)]{Krause10} supplies the source-compact orientation of his Lemma~4.7.1(1), giving \Cref{thm:abstract-rigidity} without algebraicity. The opposite orientation of the quotient criterion is also standard \cite{Keller96,Krause10,SunZhangZhang}. Efimov's Calkin exactness \cite[Proposition~1.4]{Efimov17} is the enhanced mechanism retained in \Cref{sec:obstruction}.

Singularity localisation in recollements is developed in \cite{JinYangZhou,SunZhangZhang}. The bounded hypotheses in \Cref{thm:actual-kernel} identify the displayed ambient quotient with the actual object-kernel of a singularity functor. Schlichting localisation, the completeness criteria discussed in \cite{KalckThesis,PavicShinder21}, and connective invariance \cite{BGT13} are the inputs to \Cref{app:negative-k}. The simultaneous kernel formula compares two nested dense quotients, rather than supplying a new negative-\(K\) localisation sequence. Its obstruction vanishes for all examples in this paper.

\subsection{Support applications}
The support intersection criterion and realisability over complete intersections are from \cite{AvramovBuchweitz00,BerghJorgensenEasy}; thick subcategories are classified by support in \cite{Stevenson14}. Hyperplane avoidance is identified with matrix factorisations in \cite[Theorem~3.5]{BerghJorgensenMF}; higher-codimension constructions appear in \cite{BurkeWalker15}. For \(p\)-groups the inputs include \cite{Carlson83,Carlson84,AvruninScott82,Quillen71}. The application here is the realisation of these marked subcategories as bounded defects of ordinary finite-dimensional algebras, together with the no-generator consequence. It does not classify arbitrary abstract equivalences between the resulting categories, as \Cref{ex:support-symmetry} illustrates.

\section{Scope and open questions}\label{sec:limitations}

The t-exact realisation problem is settled by \Cref{cor:texact-classification}; projective-right-target kernel categories are classified up to dg Morita equivalence in \Cref{cor:projective-target-classes}. These conclusions do not classify arbitrary full-profile markings: \Cref{ex:marking-boundary} has a connective marking with nonzero \(H^{-2}\). Ordinary homological epimorphisms with infinite right target projective dimension, without t-exactness, and the general converse for Schur kernels remain outside the result.

The profile problem retains the marking. Coordinate symmetries already give abstract equivalences between distinct cuts (\Cref{ex:support-symmetry}); the open question is to classify such equivalences and identify invariants beyond these symmetries. We make no general Boolean-lattice assertion about the thick subcategories of an avoidance category. In particular, when \(Z=\varnothing\), the whole singularity category includes objects of full support; arguments using only closed-point supports for a nonempty excluded subset of \(\PP^1\) do not apply unchanged.

The support classification requires the pair-support criterion \ref{S:pair}. For arbitrary finite groups or self-injective algebras satisfying \textup{(Fg)}, that criterion must be checked or replaced by a suitable pair-support theory. The two-block example after \Cref{thm:pgroup} shows why individual supports alone are insufficient. The complete Nakayama profile lattices are computed only in ranks \(3,4\). The relation between higher multiplication and profile posets, and the possible dg algebras of non-silting presentation cones, remain the questions \Cref{q:profile-m3,q:nonsilting}.

The heart-amplitude converse over general noetherian rings needs the finite-generation condition \((\dagger_T)\). No general checkable replacement is supplied. Likewise, the multi-object finite-type dictionary for arbitrary infinite sets of localisation maps and the telescope property for all self-injective Nakayama algebras are not established here.

The negative-\(K\) summand obstruction vanishes in every example constructed here. Producing a localisation with nonzero obstruction remains \Cref{question:negative-k-counterexample}; the simultaneous-kernel formula is not an existence result. A further question is whether asymmetric eventual Ext-vanishing can produce two genuinely incomparable marked defects by explicit extension and coextension constructions. The phenomena in \cite{HunekeJorgensen03,JorgensenSega04} motivate this question, but no such localisation pair or witnesses to both failures of containment are provided here.

\subsection*{Statements and declarations}
\paragraph{Funding.} This research was self-funded by the author. No external funding was received.
\paragraph{Author contribution.} Marouane Lamharzi Alaoui is the sole author and is responsible for the mathematical arguments, manuscript, and computational supplement.
\paragraph{Code availability.} The exact verification scripts, captured outputs, and reproduction instructions are supplied in \texttt{anc/supplement/}; see \Cref{app:computation}. The computational files carry an MIT licence. No external datasets were used.



\appendix

\section{A known factorisation criterion and perfect restriction}\label{sec:comparison}

Continue with
\[
 \cC=\cK\cap\cP=\Kc_\lambda,
 \qquad
 \cL=\thick(\cP,\cK).
\]
The inclusion \(\cK\subseteq\cL\) induces the natural exact functor
\begin{equation}\label{eq:rho}
 \rho_\lambda:
 \DeltaB(\lambda)=\cK/\cC
 \longrightarrow
 \Defsg(\lambda)=\cL/\cP.
\end{equation}
The essential image always thickly generates the target, but this says nothing by itself about morphisms.

\subsection{The standard factorisation criterion}

\begin{lemma}[Conservativity of the comparison]\label{lem:rho-conservative}
For any triangulated category \(\cT\) and thick subcategories \(\cK,\cP\), the functor
\[
 \rho:\cK/(\cK\cap\cP)
 \longrightarrow
 \thick(\cK,\cP)/\cP
\]
reflects zero objects, and hence is conservative.
\end{lemma}

\begin{proof}
An object \(K\in\cK\) becomes zero in the target if and only if
\(K\in\cP\), since \(\cP\) is thick.  This is equivalent to
\(K\in\cK\cap\cP\), which is exactly the condition that \(K\) be zero in the source.  An exact functor which reflects zero objects reflects isomorphisms by applying this to cones.
\end{proof}

\begin{lemma}[Full plus conservative implies faithful]\label{lem:full-conservative}
Let \(F:\mathcal A\to\mathcal B\) be an exact full functor of triangulated categories.  If \(F\) is conservative, then \(F\) is faithful.  We include the short argument because it separates faithfulness from roof surjectivity.
\end{lemma}

\begin{proof}
Let \(f:X\to Y\) satisfy \(F(f)=0\), and complete it to a triangle
\[
 X\xrightarrow{f}Y\xrightarrow{g}Z\longrightarrow X[1].
\]
The image triangle is split, so \(F(g)\) has a retraction
\(r:F(Z)\to F(Y)\).  Fullness lifts \(r\) to a map \(p:Z\to Y\).  Then
\(F(pg)=\id_{F(Y)}\), so \(F(pg)\) is an isomorphism.  Conservativity implies that \(pg\) is an isomorphism.  Consequently \(g\) is split monic, the original triangle is split, and \(f=0\).
\end{proof}

The full-faithfulness assertion in the next result is the standard localisation-of-subcategories criterion; see \cite[the dual of Lemma~10.3]{Keller96}, \cite[Lemma~4.7.1(2)]{Krause10}, and \cite[Lemma~2.6(1)]{SunZhangZhang}.  This is the opposite orientation from the compact-factorisation proof of \Cref{thm:abstract-rigidity}. We retain it here for the perfect-restriction application.  The thick-generation and idempotent-completion statements are formal consequences for the displayed target.

\begin{proposition}[Known factorisation criterion with Karoubi-density addendum]\label{thm:factorisation-comparison}
Let \(\cK,\cP\) be thick subcategories of a triangulated category \(\cT\), and put \(\cC=\cK\cap\cP\).  Assume
\begin{equation}\label{eq:factorization-condition}
 \text{every morphism }K\longrightarrow P,
 \quad K\in\cK,
 \quad P\in\cP,
 \quad\text{factors through an object of }\cC.
\end{equation}
Then
\[
 \rho:\cK/\cC\longrightarrow\thick(\cK,\cP)/\cP
\]
is fully faithful.  Its essential image is triangulated and thickly generates the target.  Consequently
\[
 \Kar(\cK/\cC)
 \xrightarrow{\sim}
 \Kar\bigl(\thick(\cK,\cP)/\cP\bigr).
\]
If the source is idempotent complete, or more generally if the essential image is closed under direct summands, then \(\rho\) is an equivalence before idempotent completion.
\end{proposition}

\begin{proof}
We first prove fullness.  Let \(K_1,K_2\in\cK\), and represent a morphism between their images in \(\cT/\cP\) by a roof
\[
 K_1\xleftarrow{s}Z\xrightarrow{g}K_2,
 \qquad \Cone(s)\in\cP.
\]
Choose a triangle
\[
 Z\xrightarrow{s}K_1\xrightarrow{u}P\longrightarrow Z[1]
\]
with \(P\in\cP\).  By \eqref{eq:factorization-condition}, write
\(u=ba\) with
\[
 K_1\xrightarrow{a}C\xrightarrow{b}P,
 \qquad C\in\cC.
\]
Complete \(a\) to a triangle
\[
 Z'\xrightarrow{s'}K_1\xrightarrow{a}C\longrightarrow Z'[1].
\]
The identity of \(K_1\) together with \(b:C\to P\) extends, by the morphism axiom for triangles, to a morphism of triangles.  In particular there is a map \(t:Z'\to Z\) with
\(st=s'\).  Since \(K_1,C\in\cK\), the object \(Z'\) belongs to \(\cK\), and \(\Cone(s')\cong C\in\cC\).  The roof
\[
 K_1\xleftarrow{s'}Z'\xrightarrow{gt}K_2
\]
represents the same target morphism and is a roof in \(\cK/\cC\).  Thus \(\rho\) is full.

By \Cref{lem:rho-conservative}, \(\rho\) is conservative.  It is therefore faithful by \Cref{lem:full-conservative}.  The essential image is triangulated because \(\rho\) is exact and fully faithful.  Since
\(\thick(\cK,\cP)/\cP\) is the thick closure of the image of \(\cK\), the image of \(\rho\) becomes essentially surjective after adjoining direct summands.  This yields the Karoubi equivalence.  If the image is already thick, no completion is needed.  In particular, a fully faithful functor from an idempotent-complete source has image closed under direct summands: a target idempotent lifts by fullness, splits in the source, and its splitting maps to the given target splitting.
\end{proof}

\begin{remark}
No novelty is claimed for the full-faithfulness criterion.  The proof is retained to expose the internal-roof construction in the direction used here and to separate fullness from faithfulness via \Cref{lem:rho-conservative,lem:full-conservative}; the Karoubi-density conclusion is formal for this target.
\end{remark}

\subsection{Perfect restriction}

\begin{theorem}[Perfect-restriction comparison theorem]\label{thm:perfect-restriction}
Let \(\lambda:R\to S\) be a homological ring epimorphism satisfying \textup{(B1)--(B3)}, and assume
\[
 U(\per(S))\subseteq\per(R),
 \quad\text{equivalently }{}_RS\in\per(R).
\]
Then \(\rho_\lambda\) in \eqref{eq:rho} is fully faithful and
\[
 \Kar\DeltaB(\lambda)
 \xrightarrow{\sim}
 \Kar\Defsg(\lambda).
\]
If \(\DeltaB(\lambda)\) is idempotent complete, then \(\rho_\lambda\) is an equivalence before completion.
\end{theorem}

\begin{proof}
We verify \eqref{eq:factorization-condition}.  Let
\(K\in\Kb(\lambda)\), \(P\in\per(R)\), and let \(f:K\to P\).  Complete the localisation unit at \(P\) to a triangle
\[
 K_P\longrightarrow P\xrightarrow{\eta_P}U\Phi P\longrightarrow K_P[1].
\]
The object \(K_P\) lies in the unbounded kernel.  It is bounded because both \(P\) and \(U\Phi P\) are bounded.  Moreover \(\Phi P\in\per(S)\), and perfect restriction gives \(U\Phi P\in\per(R)\).  Hence \(K_P\in\per(R)\), so
\[
 K_P\in\Kb(\lambda)\cap\per(R)=\Kc_\lambda
\]
by \Cref{lem:compact-intersection}.

Adjunction and \(\Phi K=0\) give
\[
 \Hom_{D(R)}(K,U\Phi P)
 \cong
 \Hom_{D(S)}(\Phi K,\Phi P)=0.
\]
Therefore \(\eta_Pf=0\).  Exactness of \(\Hom(K,-)\) on the unit triangle implies that \(f\) factors through \(K_P\in\Kc_\lambda\).  The hypotheses of \Cref{thm:factorisation-comparison} are satisfied, and the conclusion follows.
\end{proof}

\begin{corollary}[Vanishing criterion]\label{cor:comparison-vanishing}
Under the hypotheses of \Cref{thm:perfect-restriction}, the following are equivalent:
\[
 \DeltaB(\lambda)=0,
 \qquad
 \Defsg(\lambda)=0,
 \qquad
 \Dsg(R)\xrightarrow{\sim}\Dsg(S).
\]
Here the final equivalence holds before idempotent completion.
\end{corollary}

\begin{proof}
Full faithfulness gives \(\DeltaB=0\) if and only if its essential image is zero.  Since that image thickly generates \(\Defsg\), this is equivalent to \(\Defsg=0\).  The last equivalence follows from \Cref{thm:actual-kernel}.
\end{proof}

\medskip
\noindent\textbf{Secondary applications and computations.}
The structural comparison package is complete at this point.  

\section{The derived factorisation obstruction}\label{sec:obstruction}

\begin{proposition}[Compatible \(h\)-flat Morita models]\label{prop:compatible-dg-models}
Let \(\cT\) be an algebraic compactly generated triangulated category, let \(\cP=\cT^c\), and let \(\cC\subseteq\cP\) be essentially small and thick.  Put \(\cL=\Loc(\cC)\).  If \(\cC\subseteq\cK\subseteq\cL\) is essentially small and thick and \(\cA=\thick(\cP,\cK)\), there are small \(h\)-flat dg categories
\[
 \cC_{\mathrm{dg}}\subseteq\cK_{\mathrm{dg}}\subseteq\cA_{\mathrm{dg}},
 \qquad
 \cC_{\mathrm{dg}}\subseteq\cP_{\mathrm{dg}}\subseteq\cA_{\mathrm{dg}},
\]
with full dg inclusions, whose homotopy categories are triangulated and idempotent complete, together with equivalences
\[
 H^0(\cC_{\mathrm{dg}})\simeq\cC,\quad
 H^0(\cK_{\mathrm{dg}})\simeq\cK,\quad
 H^0(\cP_{\mathrm{dg}})\simeq\cP,\quad
 H^0(\cA_{\mathrm{dg}})\simeq\cA.
\]
Writing \(i:\cC_{\mathrm{dg}}\to\cP_{\mathrm{dg}}\) for the inclusion and \(j:\cL\to\cT\) for the ambient functor, there are triangle equivalences
\[
 \Theta_{\cC}:D(\cC_{\mathrm{dg}})\xrightarrow{\sim}\cL,
 \qquad
 \Theta_{\cP}:D(\cP_{\mathrm{dg}})\xrightarrow{\sim}\cT
\]
and a natural isomorphism
\[
 \Theta_{\cP}\circ\mathbf L i_!\xrightarrow{\sim}j\circ\Theta_{\cC},
\]
where \(\mathbf L i_!\) is derived extension of scalars.  Dg Yoneda identifies
\[
 H^0(\cE_{\mathrm{dg}})\xrightarrow{\sim}\per(\cE_{\mathrm{dg}})
 \qquad(\cE=\cC,\cP).
\]
Consequently
\[
 H^0\!\left(\operatorname{Mod}_{\mathrm{cof}}(\cE_{\mathrm{dg}})/Y(\cE_{\mathrm{dg}})\right)
 \simeq D(\cE_{\mathrm{dg}})/\per(\cE_{\mathrm{dg}}).
\]
Compatible cofibrant \(h\)-flat replacement changes none of these derived categories, perfect subcategories, or Verdier quotients.
\end{proposition}

\begin{proof}
By algebraicity and compact generation, Keller's dg Morita theorem gives a small dg category \(\mathcal R\) and an equivalence \(D(\mathcal R)\simeq\cT\) \cite{Keller94}.  Replace \(\mathcal R\) by a cofibrant \(h\)-flat model.  In the dg category of cofibrant right \(\mathcal R\)-modules choose representatives of a small skeleton of \(\cA\), using the same representatives for objects in \(\cC,\cK,\cP\).  This is possible because these categories are essentially small.  Moreover
\[
 \cL\cap\cT^c=\cC
\]
by Thomason--Neeman, so \(\cK\cap\cP=\cC\) and the object sets can be chosen compatibly.

Let \(\cA_0\) be the full dg category on the selected modules, with representatives chosen also for all shifts, cones, and existing direct summands in the four thick subcategories.  Apply one identity-on-objects semi-free \(h\)-flat resolution to \(\cA_0\), and take the full dg subcategories on the four selected object sets.  Because the ambient resolution is made once, all inclusions remain full and compatible.  The restrictions are quasi-equivalences with the original full dg subcategories, so the models are pretriangulated and idempotent complete up to quasi-equivalence; their homotopy categories identify with \(\cC,\cK,\cP,\cA\) with the induced triangulated structures.  Standard cofibrant dg-category replacement and quasi-equivalence invariance give the asserted invariance statements; see \cite{Tabuada05,Keller94}.

Let \(\mathcal U_{\cP}\) be the \(\cP_{\mathrm{dg}}\)-\(\mathcal R\) evaluation bimodule given by the selected objects and let \(\mathcal U_{\cC}\) be its restriction.  The realisation functors
\[
 -\Ltensor_{\cP_{\mathrm{dg}}}\mathcal U_{\cP}
 \quad\text{and}\quad
 -\Ltensor_{\cC_{\mathrm{dg}}}\mathcal U_{\cC}
\]
are equivalences onto \(\cT\) and \(\cL\), respectively: the first object set is all of \(\cT^c\), while the second compactly generates \(\cL\).  Derived coend associativity and Yoneda identify the bimodule for the first functor after extension along \(i\) with \(\mathcal U_{\cC}\).  This proves the compatible square.

Finally, the Yoneda images of \(H^0(\cC_{\mathrm{dg}})\) and \(H^0(\cP_{\mathrm{dg}})\) are triangulated and closed under direct summands, and they contain all representables.  They are therefore exactly the perfect modules.  The final identifications follow from the Drinfeld quotient construction \cite{Drinfeld04}.
\end{proof}

For the rest of this section we suppress the dg subscript and work with the compatible full \(h\)-flat models of \Cref{prop:compatible-dg-models}.  Let \(\cA\) be generated by the full dg subcategories \(\cK,\cP\), and put \(\cC=\cK\cap\cP\).  For \(X,Y\in\cK\) and a full dg subcategory \(\cB\), define
\[
 \Fact_{\cB}(X,Y)=
 \cA(-,Y)|_{\cB}\Ltensor_{\cB}\cA(X,-)|_{\cB}.
\]
This is the derived complex of round trips \(X\to B\to Y\), including the higher bar terms.

\begin{definition}[Derived factorisation obstruction]\label{def:obs}
For \(X,Y\in\cK\), put
\[
 \Obs_{\cK,\cP}(X,Y)=
 \Cone\!\left(\Fact_{\cC}(X,Y)\longrightarrow\Fact_{\cP}(X,Y)\right).
\]
\end{definition}

\begin{theorem}[Obstruction triangle]\label{thm:obs-triangle}
There is a functorial triangle of morphism complexes
\[
 (\cK/\cC)(X,Y)\longrightarrow(\cA/\cP)(X,Y)
 \longrightarrow\Obs_{\cK,\cP}(X,Y)[1]\longrightarrow.
\]
Consequently the quotient functor is quasi-fully faithful if and only if the obstruction is acyclic for every pair \(X,Y\).
\end{theorem}

\begin{proof}
For a full dg subcategory \(\cB\), the bar model of the Drinfeld quotient~\cite[\S1.3.1]{Drinfeld04} gives a triangle
\[
 \Fact_{\cB}(X,Y)\longrightarrow\cA(X,Y)
 \longrightarrow(\cA/\cB)(X,Y)\longrightarrow.
\]
Apply it to \(\cC\subset\cK\) and to \(\cP\subset\cA\).  The middle terms agree because \(\cK\) is full in \(\cA\).  The inclusion \(\cC\subset\cP\) produces a morphism of triangles.  The octahedral axiom identifies the cone of the map between quotient Hom complexes with the shift by one of the cone of \(\Fact_{\cC}\to\Fact_{\cP}\).  This is the displayed triangle.  Quasi-full faithfulness is equivalent to the first map being a quasi-isomorphism for every pair.
\end{proof}

For a localisation, take
\[
 \cK=\Kb(\lambda),\qquad \cC=\Kc_\lambda,
 \qquad \cP=\per R,
 \qquad \cA=\thick(\cP,\cK),
\]
and denote the resulting bifunctor by \(\Obs(\lambda)\).

\begin{corollary}[Comparison criterion]\label{cor:obs-comparison}
The comparison \(\rho_\lambda\) is fully faithful if and only if \(\Obs(\lambda)(X,Y)\) is acyclic for all \(X,Y\in\Kb(\lambda)\).  When this holds, \(\rho_\lambda\) is an equivalence after idempotent completion.
\end{corollary}

\begin{proof}
Apply \Cref{thm:obs-triangle}.  The target is the thick closure of the image because \(\cA=\thick(\cP,\cK)\); adjoining direct summands therefore makes a fully faithful comparison essentially surjective.
\end{proof}

If \(T\) compactly generates \(\cC\) and \(E=\RHom_R(T,T)\) before taking an opposite, then
\begin{equation}\label{eq:one-object-obs}
 \Obs(\lambda)(X,Y)\simeq
 \Cone\!\left[
 \RHom_R(T,Y)\Ltensor_E\RHom_R(X,T)
 \longrightarrow
 Y\Ltensor_{R^{\op}}\RHom_R(X,R)
 \right].
\end{equation}
This is the finite linear-algebra form used below.

\subsection{Calkin explanation of the comparison}
For a small \(h\)-flat dg category \(\cC\), let \(\Mod\cC\) be the dg category of cofibrant right dg modules and let
\[
 \Calk_{\cC}=\Mod\cC/Y(\cC)
\]
be the Calkin dg quotient.  If its homotopy category is triangulated and idempotent complete as in \Cref{prop:compatible-dg-models}, dg Yoneda identifies the homotopy category of this quotient with \(D(\cC)/\per(\cC)\).  For cofibrant modules \(M,N\), the Hom complex is the cone of evaluation
\[
 N\Ltensor_{\cC}M^\vee\longrightarrow\RHom_{\cC}(M,N),
 \qquad M^\vee=\RHom_{\cC}(M,\cC).
\]

\begin{lemma}[Calkin exactness for a full inclusion]\label{lem:calkin-full}
If \(F:\cC\to\cP\) is quasi-fully faithful, derived extension of scalars induces a quasi-fully faithful dg functor
\[
 \Calk_F:\Calk_{\cC}\longrightarrow\Calk_{\cP}.
\]
\end{lemma}

\begin{proof}
Write \(F^*\) for extension and \(F_*\) for restriction.  Quasi-full faithfulness and derived Yoneda give, for a cofibrant \(\cC\)-module \(M\),
\[
 M^\vee\xrightarrow{\sim}F_*((F^*M)^\vee).
\]
Derived coend associativity gives the projection formula
\[
 N\Ltensor_{\cC}F_*L\xrightarrow{\sim}F^*N\Ltensor_{\cP}L.
\]
Taking \(L=(F^*M)^\vee\) identifies the two factorisation complexes.  Derived adjunction identifies
\[
 \RHom_{\cC}(M,N)\xrightarrow{\sim}
 \RHom_{\cP}(F^*M,F^*N).
\]
The evaluation triangles form a morphism of triangles whose first two vertical arrows are quasi-isomorphisms; hence so is the third, the Calkin Hom map.
\end{proof}

This is the quasi-full-faithfulness part of Efimov's Proposition~1.4~\cite{Efimov17}.  It explains the enhancement of the classical comparison.


To recover \Cref{thm:abstract-rigidity} with these models, use the quasi-fully faithful inclusion \(\cC_{\mathrm{dg}}\to\cP_{\mathrm{dg}}\) of \Cref{prop:compatible-dg-models}. The resulting Calkin functor identifies, on homotopy categories, with \(\cL/\cC\to\cT/\cP\). Restriction to the full subcategories \(\cK/\cC\) and \(\cA/\cP\) gives the comparison: any roof between their objects whose left cone belongs to the denominator has its middle object in the same subcategory. Thus the dg comparison is quasi-fully faithful, and \Cref{thm:obs-triangle} makes its obstruction acyclic. This is an enhanced explanation of the classical factorisation proof, not an additional triangulated novelty claim.

\section{The summand obstruction in negative \texorpdfstring{\(K\)}{K}-theory}\label{app:negative-k}\label{subsec:summand-negative-k}
Throughout this appendix, \(R\) is left noetherian and \(\lambda:R\to R_\Sigma\) is a stably flat universal localisation. Let \(\cE\) be the common idempotent completion and put
\[
 H_\Delta=\im\bigl(K_0(\DeltaB(\lambda))\to K_0(\cE)\bigr),\qquad
 H_{\mathrm{Def}}=\im\bigl(K_0(\Defsg(\lambda))\to K_0(\cE)\bigr).
\]

\begin{proposition}[Thomason summand obstruction]\label{prop:k0-obstruction}
The quotient
\[
 \operatorname{SObs}_0(\lambda)=H_{\mathrm{Def}}/H_\Delta
\]
vanishes if and only if \(\rho_\lambda\) is essentially surjective before idempotent completion.
\end{proposition}

\begin{proof}
Via the fully faithful comparison, regard \(\DeltaB(\lambda)\) and \(\Defsg(\lambda)\) as dense full triangulated subcategories of their common idempotent completion \(\cE\), with the former contained in the latter.  Thomason's Theorem~2.1~\cite{Thomason97} identifies dense triangulated subcategories of \(\cE\) with subgroups of its triangulated Grothendieck group \(K_0(\cE)\).  The two subcategories correspond to \(H_\Delta\subseteq H_{\mathrm{Def}}\).  Thus they coincide exactly when \(H_\Delta=H_{\mathrm{Def}}\), equivalently when \(\operatorname{SObs}_0(\lambda)=0\).
\end{proof}

\paragraph{K-theoretic conventions.}
For every category occurring below, \(K_{-n}\) means Schlichting's group \(\mathrm{IK}_{-n}\) of the Frobenius model specified in \Cref{lem:negative-k-models}.  Maps on these groups are induced by the displayed maps of Frobenius pairs.  By Schlichting's Corollary~1.8~\cite{Schlichting06}, the groups are invariant under derived equivalences induced by maps of models.  In degree zero,
\[
 \mathrm{IK}_0(\mathfrak X)=K_0\bigl(\Kar D(\mathfrak X)\bigr),
\]
where the right-hand side is the triangulated Grothendieck group, not the split Grothendieck group.  Thus the \(K_0\) used here is the same Grothendieck group as in Thomason's Theorem~2.1.  The finite-complex Frobenius model of \(\per R\), Schlichting's Corollary~1.8, and the comparison of Section~7.1 identify
\[
 K_{-1}(\per R)\cong K_{-1}(R),
\]
where the right-hand side is the classical Bass group.  By \Cref{lem:negative-k-models}, the categories \(\cC,\cK,\cP,\cA\) below are idempotent complete.  Their degree-zero terms are therefore ordinary triangulated Grothendieck groups; only the two Verdier-quotient terms require idempotent completion in Schlichting's convention.

\begin{lemma}[Compatible Frobenius models for the two quotient rows]\label{lem:negative-k-models}
Assume the hypotheses of \Cref{thm:unconditional-comparison}, and put
\[
 \cC=\Kc_\lambda,\qquad
 \cK=\Kb(\lambda),\qquad
 \cP=\per(R),\qquad
 \cA=\thick(\cP,\cK).
\]
There is a commutative diagram of exact sequences of Frobenius pairs
\[
\begin{array}{ccccc}
 (\cF_{\cC},\cF_0)&\longrightarrow&(\cF_{\cK},\cF_0)&\longrightarrow&(\cF_{\cK},\cF_{\cC})\\
 \big\downarrow&&\big\downarrow&&\big\downarrow\\
 (\cF_{\cP},\cF_0)&\longrightarrow&(\cF,\cF_0)&\longrightarrow&(\cF,\cF_{\cP})
\end{array}
\tag{\(*\)}\label{eq:frobenius-pair-square}
\]
whose derived diagram is
\[
\begin{array}{ccccc}
 \cC&\longrightarrow&\cK&\longrightarrow&\cK/\cC\\
 \big\downarrow&&\big\downarrow&&\big\downarrow\rho_\lambda\\
 \cP&\longrightarrow&\cA&\longrightarrow&\cA/\cP.
\end{array}
\tag{\(**\)}\label{eq:derived-quotient-square}
\]
The first two terms in both rows of \eqref{eq:frobenius-pair-square} use the same acyclic subcategory \(\cF_0\), and
\[
 \cF_{\cC}=\cF_{\cK}\cap\cF_{\cP}.
\]
Moreover, \(\cC,\cK,\cP,\cA\) are idempotent complete.
\end{lemma}

\begin{proof}
Since \(R\) is left noetherian, \(R\text{-}\modcat\) is abelian and hence idempotent complete.  Balmer--Schlichting's theorem, in the form recalled in Schlichting's Section~5.5~\cite{BalmerSchlichting01,Schlichting06}, therefore shows that \(\Db(R\text{-}\modcat)\) is idempotent complete.  The categories \(\cC,\cK,\cP\), and \(\cA\) are thick subcategories of \(\Db(R\text{-}\modcat)\); hence they are idempotent complete.

The category \(R\text{-}\Mod\) is Grothendieck abelian, and
\[
 \cA\subseteq\Db(R\text{-}\modcat)\subseteq D(R\text{-}\Mod)
\]
is essentially small.  After choosing a small skeleton, Schlichting's construction in Section~5.6~\cite{Schlichting06} supplies a Frobenius pair \((\cF,\cF_0)\) and a triangle equivalence
\[
 D(\cF,\cF_0)\xrightarrow{\sim}\cA.
\]

Apply Schlichting's construction in Section~5.2 first to the strictly full triangulated subcategories \(\cK\) and \(\cP\) of \(D(\cF,\cF_0)\).  Thus \(\cF_{\cK}\) and \(\cF_{\cP}\) are the full subcategories of objects whose images are isomorphic to objects of \(\cK\) and \(\cP\), respectively, and
\[
 (\cF_{\cK},\cF_0),\qquad (\cF_{\cP},\cF_0)
\]
are Frobenius pairs modelling \(\cK\) and \(\cP\).  Apply the same construction to \(\cC\subseteq\cK\) inside the model \((\cF_{\cK},\cF_0)\).  It gives the sub-pair \((\cF_{\cC},\cF_0)\) and the quotient pair
\[
 \mathfrak Q_{\cC}=(\cF_{\cK},\cF_{\cC}),
\]
with derived category \(\cK/\cC\).  The lower quotient pair is
\[
 \mathfrak Q_{\cP}=(\cF,\cF_{\cP}),
\]
with derived category \(\cA/\cP\).

By Lemma~3.2,
\[
 \cC=\cK\cap\cP.
\]
Since all four triangulated subcategories are strictly full, an object of \(\cF\) belongs to \(\cF_{\cC}\) exactly when its image belongs to both \(\cK\) and \(\cP\).  Hence \(\cF_{\cC}=\cF_{\cK}\cap\cF_{\cP}\).  In particular, the inclusion \(\cF_{\cK}\subseteq\cF\) sends \(\cF_{\cC}\) into \(\cF_{\cP}\), and therefore induces the right-hand vertical map in \eqref{eq:frobenius-pair-square}.  Schlichting's Section~5.2 identifies both rows as exact sequences of Frobenius pairs and identifies their derived rows with \eqref{eq:derived-quotient-square}.
\end{proof}

\begin{theorem}[Negative-\(K\) realisation of the summand obstruction]\label{thm:negative-k-obstruction}
Assume the hypotheses of \Cref{thm:unconditional-comparison}, and use the notation of \Cref{lem:negative-k-models}.  Fix
\[
 \cE=\Kar\Defsg(\lambda)
\]
and identify \(\Kar\DeltaB(\lambda)\) with \(\cE\) through the equivalence \(\Kar(\rho_\lambda)\) of \Cref{thm:unconditional-comparison}.  Let
\[
 i_{\cK}:\cC\hookrightarrow\cK,
 \qquad
 i_{\cP}:\cC\hookrightarrow\cP
\]
be the inclusions.  There is a commutative diagram with exact rows
\[
\begin{array}{ccccccccc}
 K_0(\cC)&\longrightarrow&K_0(\cK)&\xrightarrow{\ \bar q_{\cK,*}\ }&K_0(\cE)&\xrightarrow{\ \partial_{\cC}\ }&K_{-1}(\cC)&\xrightarrow{\ i_{\cK,*}\ }&K_{-1}(\cK)\\[1mm]
 \big\downarrow&&\big\downarrow&&\big\Vert&&\big\downarrow i_{\cP,*}&&\big\downarrow\\[1mm]
 K_0(\cP)&\longrightarrow&K_0(\cA)&\xrightarrow{\ \bar q_{\cA,*}\ }&K_0(\cE)&\xrightarrow{\ \partial_{\cP}\ }&K_{-1}(\cP)&\longrightarrow&K_{-1}(\cA),
\end{array}
\tag{\(\dagger\)}\label{eq:negative-k-two-rows}
\]
where
\[
 \bar q_{\cK}:\cK\longrightarrow\DeltaB(\lambda)
 \longrightarrow\Kar\DeltaB(\lambda)\xrightarrow{\sim}\cE
\]
and
\[
 \bar q_{\cA}:\cA\longrightarrow\Defsg(\lambda)
 \longrightarrow\Kar\Defsg(\lambda)=\cE.
\]
All unlabeled arrows in \eqref{eq:negative-k-two-rows} are induced by the displayed inclusions or quotient functors.  The boundary maps satisfy
\[
 \partial_{\cP}=i_{\cP,*}\partial_{\cC}.
\tag{\(\ddagger\)}\label{eq:boundary-naturality}
\]
Moreover,
\[
 H_\Delta=\ker\partial_{\cC},
 \qquad
 H_{\mathrm{Def}}=\ker\partial_{\cP},
\tag{\(\star\)}\label{eq:H-as-boundary-kernels}
\]
and the assignment
\[
 \Theta_\lambda:
 \operatorname{SObs}_0(\lambda)\longrightarrow K_{-1}(\cC),
 \qquad
 x+H_\Delta\longmapsto\partial_{\cC}(x)\quad(x\in H_{\mathrm{Def}}),
\]
is a canonical isomorphism
\[
 \boxed{
 \operatorname{SObs}_0(\lambda)
 \xrightarrow{\sim}
 \ker\!\left(
 K_{-1}(\Kc_\lambda)
 \xrightarrow{\ (i_{\cK,*},i_{\cP,*})\ }
 K_{-1}(\Kb(\lambda))\oplus K_{-1}(R)
 \right).}
\tag{\(\heartsuit\)}\label{eq:sobs-negative-k}
\]
Consequently,
\[
 \operatorname{SObs}_0(\lambda)=0
 \quad\Longleftrightarrow\quad
 (i_{\cK,*},i_{\cP,*})\text{ is injective}.
\tag{\(\clubsuit\)}\label{eq:sobs-injectivity}
\]
This equality requires only \(H_\Delta=H_{\mathrm{Def}}\).  By contrast, the stronger hypothesis \(K_{-1}(\Kc_\lambda)=0\) forces
\[
 H_\Delta=H_{\mathrm{Def}}=K_0(\cE).
\tag{\(\spadesuit\)}\label{eq:stronger-k0-density}
\]
\end{theorem}

\begin{proof}
Apply Schlichting's Theorem~11.10~\cite{Schlichting06} to the two exact sequences of Frobenius pairs in \eqref{eq:frobenius-pair-square}.  The theorem is functorial for maps of exact sequences of Frobenius pairs, so it gives a morphism between the two localisation sequences.  Before identifying their middle terms, the middle vertical map is
\[
 \mathrm{IK}_0\bigl(\Kar(\rho_\lambda)\bigr):
 K_0\bigl(\Kar\DeltaB(\lambda)\bigr)
 \longrightarrow
 K_0\bigl(\Kar\Defsg(\lambda)\bigr).
\]
It is an isomorphism by \Cref{thm:unconditional-comparison}; we identify both groups with \(K_0(\cE)\).  This gives \eqref{eq:negative-k-two-rows}.  Since \(\cC,\cK,\cP,\cA\) are idempotent complete by \Cref{lem:negative-k-models}, the other degree-zero terms in the diagram are their ordinary triangulated Grothendieck groups.  Functoriality in Schlichting's Theorem~11.10 gives
\[
 \partial_{\cP}\circ
 \mathrm{IK}_0\bigl(\Kar(\rho_\lambda)\bigr)
 =i_{\cP,*}\circ\partial_{\cC};
\]
after the preceding identification, this is \eqref{eq:boundary-naturality}.

The triangulated Grothendieck group is right exact on Verdier quotients.  Hence
\[
 K_0(\cK)\twoheadrightarrow K_0(\cK/\cC),
 \qquad
 K_0(\cA)\twoheadrightarrow K_0(\cA/\cP).
\]
It follows that
\[
 H_\Delta=\im\bigl(K_0(\cK)\to K_0(\cE)\bigr),
 \qquad
 H_{\mathrm{Def}}=\im\bigl(K_0(\cA)\to K_0(\cE)\bigr).
\]
Exactness of the two rows of \eqref{eq:negative-k-two-rows} at \(K_0(\cE)\) proves \eqref{eq:H-as-boundary-kernels}.

Since \(H_\Delta=\ker\partial_{\cC}\), the formula defining \(\Theta_\lambda\) is well defined and injective.  By \eqref{eq:boundary-naturality}, its image is
\[
 \partial_{\cC}(\ker\partial_{\cP})
 =\im\partial_{\cC}\cap\ker i_{\cP,*}.
\]
Exactness of the upper row at \(K_{-1}(\cC)\) gives
\[
 \im\partial_{\cC}=\ker i_{\cK,*}.
\]
Therefore
\[
 \im\Theta_\lambda
 =\ker i_{\cK,*}\cap\ker i_{\cP,*}
 =\ker(i_{\cK,*},i_{\cP,*}).
\]
For completeness, let \(\eta\) lie in this last kernel.  Exactness of the upper row gives \(x\in K_0(\cE)\) with \(\partial_{\cC}(x)=\eta\).  Then
\[
 \partial_{\cP}(x)=i_{\cP,*}(\eta)=0,
\]
so \(x\in H_{\mathrm{Def}}\), and \(\Theta_\lambda(x+H_\Delta)=\eta\).  Thus \(\Theta_\lambda\) is surjective onto the simultaneous kernel.  Finally, the finite-complex comparison recalled in the conventions identifies \(K_{-1}(\cP)=K_{-1}(\per R)\) with \(K_{-1}(R)\), proving \eqref{eq:sobs-negative-k}.

The equivalence \eqref{eq:sobs-injectivity} follows immediately.  If \(K_{-1}(\cC)=0\), then \(\partial_{\cC}=0\), so \eqref{eq:H-as-boundary-kernels} gives \(H_\Delta=K_0(\cE)\).  Since
\[
 H_\Delta\subseteq H_{\mathrm{Def}}\subseteq K_0(\cE),
\]
this proves \eqref{eq:stronger-k0-density}.
\end{proof}

\begin{remark}[The logical asymmetry]\label{rem:logical-asymmetry}
The exact criterion for essential surjectivity is the injectivity in \eqref{eq:sobs-injectivity}, or equivalently the equality \(H_\Delta=H_{\mathrm{Def}}\).  The vanishing of \(K_{-1}(\Kc_\lambda)\) is only a sufficient condition.  It proves the stronger statement that \(\DeltaB(\lambda)\) is dense in the whole idempotent completion at the level of \(K_0\), as in \eqref{eq:stronger-k0-density}.  Essential surjectivity may therefore hold even when \(K_{-1}(\Kc_\lambda)\ne0\), provided that the simultaneous kernel in \eqref{eq:sobs-negative-k} is zero.
\end{remark}

\begin{remark}[The boundary in Schlichting's construction]\label{rem:schlichting-boundary}
We record the object-level construction from Schlichting's Section~1.5~\cite{Schlichting06}.  Apply the flasque and suspension functors to the upper exact sequence in \eqref{eq:frobenius-pair-square}.  An object \(e\) of the idempotent completion of \(D(\mathfrak Q_{\cC})\) is lifted, along the flasque localisation, to an object \(B\) in the flasque model of \((\cF_{\cK},\cF_0)\).  Its image in the suspension model becomes zero in the suspension of \(\mathfrak Q_{\cC}\), and therefore lies in the suspension model of \((\cF_{\cC},\cF_0)\).  The boundary is
\[
 \partial_{\cC}[e]=[B_{\mathrm s}]
 \in K_0\bigl(D(S(\cF_{\cC},\cF_0))\bigr)
 =K_{-1}(\cC).
\]
Schlichting proves independence of the lift and compatibility with triangle relations.  Naturality in \eqref{eq:boundary-naturality} is not deduced from choices in this construction; it is the functoriality assertion of his Theorem~11.10 applied to \eqref{eq:frobenius-pair-square}.
\end{remark}

\begin{lemma}[Connective invariance]\label{lem:connective-negative-k}
Let \(B\) be a cohomologically non-positive dga and put \(\Gamma=H^0(B)\).  For every \(n\ge1\), the canonical map \(B\to\Gamma\) induces an isomorphism
\[
 K_{-n}(\per B)\xrightarrow{\sim}K_{-n}(\Gamma).
\]
\end{lemma}

\begin{proof}
Let \(HB\) be the Eilenberg--Mac Lane \(A_\infty\)-ring spectrum associated with \(B\).  Since \(B\) is cohomologically non-positive, \(HB\) is connective and
\[
 \pi_0(HB)=H^0(B)=\Gamma.
\]
The category of perfect \(B\)-modules identifies with the category of compact \(HB\)-modules.  Blumberg--Gepner--Tabuada's Theorem~9.53~\cite{BGT13} states that, for a connective \(A_\infty\)-ring spectrum \(A\), the map
\[
 A\longrightarrow H\pi_0A
\]
induces isomorphisms
\[
 \pi_iK(A)\xrightarrow{\sim}\pi_iK(\pi_0A)
 \qquad(i\le0).
\]
Applied to \(A=HB\), this gives the asserted isomorphism.

Alternatively, Bass delooping gives a three-step proof.  The map \(K_0(A)\to K_0(\pi_0A)\) is an isomorphism for every connective ring spectrum \(A\).  Laurent extension preserves connectivity and satisfies \(\pi_0(A[t_1^{\pm1},\ldots,t_r^{\pm1}])=\pi_0(A)[t_1^{\pm1},\ldots,t_r^{\pm1}]\).  Applying the first statement to all iterated Laurent extensions and inducting through the Bass fundamental exact sequence yields the isomorphism in every negative degree.
\end{proof}

\begin{remark}[Compatibility of nonconnective models]\label{rem:bgt-schlichting-compatibility}
Blumberg--Gepner--Tabuada's proof of Theorem~9.53 invokes the argument of Schlichting's Theorem~11.7, which compares the nonconnective construction used there with Schlichting's delooping.  In the dga model, Schlichting's Section~5.9 defines \(\mathrm{IK}(B)\) from finite cell modules, and Corollary~1.8 gives quasi-isomorphism invariance.  Thus the connective invariance used in \Cref{lem:connective-negative-k} is compatible with the Schlichting groups and dg models used elsewhere in this subsection.
\end{remark}

\begin{lemma}[Artinian vanishing]\label{lem:artinian-negative-k}
If \(\Gamma\) is an artinian ring, then
\[
 K_{-n}(\Gamma)=0
 \qquad(n\ge1).
\]
\end{lemma}

\begin{proof}
Let \(J=\rad\Gamma\).  The ideal \(J\) is nilpotent and \(\Gamma/J\) is semisimple artinian.  Corti{\~n}as defines
\[
 K_{-n}(A)=K_0(\Sigma^nA)
\]
using Karoubi's suspension, and Proposition~2.5.2(i)~\cite{Cortinas11} states that, in non-positive degrees, these functors are additive, nilinvariant, and \(M_\infty\)-stable.  These groups agree with Schlichting's groups by Schlichting's Section~7.1.  Hence
\[
 K_{-n}(\Gamma)\xrightarrow{\sim}K_{-n}(\Gamma/J).
\]
The ring \(\Gamma/J\) is regular, so its negative groups vanish by Schlichting's Example~9.5.  The result follows.
\end{proof}

\begin{lemma}[Regular-noetherian vanishing]\label{lem:regular-noetherian-negative-k}
Let \(\Gamma\) be left noetherian, and assume that every finitely generated left \(\Gamma\)-module has finite projective dimension.  Then
\[
 K_{-n}(\Gamma)=0
 \qquad(n\ge1).
\]
\end{lemma}

\begin{proof}
The resolution hypothesis gives a triangle equivalence
\[
 \Kb(\Gamma\text{-}\operatorname{proj})
 \xrightarrow{\sim}
 \Db(\Gamma\text{-}\modcat).
\]
After choosing a small skeleton, the abelian category \(\Gamma\text{-}\modcat\) is small and noetherian.  Schlichting's Theorem~9.3~\cite{Schlichting06} therefore gives vanishing of all its negative \(K\)-groups.  Corollary~1.8 transfers this vanishing across the displayed derived equivalence.
\end{proof}

\begin{theorem}[Pre-completion equivalence criteria]\label{thm:negative-k-vanishing}
Under the hypotheses of \Cref{thm:unconditional-comparison}, each of the following conditions implies that
\[
 \rho_\lambda:\DeltaB(\lambda)\xrightarrow{\sim}\Defsg(\lambda)
\]
is an equivalence before idempotent completion.
\begin{enumerate}[label=\textup{(\arabic*)}]
\item The map
\[
 K_{-1}(\Kc_\lambda)
 \longrightarrow
 K_{-1}(\Kb(\lambda))\oplus K_{-1}(R)
\]
is injective.
\item \(K_{-1}(\Kc_\lambda)=0\).
\item The small stable category underlying a dg enhancement of \(\Kc_\lambda\) admits a bounded \(t\)-structure.
\item The category \(\Kc_\lambda\) admits a silting object \(T\), and, for
\[
 B_T=\tau_{\le0}\RHom_R(T,T)^{\op},
 \qquad
 \Gamma_T=H^0(B_T),
\]
one has \(K_{-1}(\Gamma_T)=0\).
\item The category \(\Kc_\lambda\) admits a silting object \(T\) for which \(\Gamma_T\) is artinian.
\item The category \(\Kc_\lambda\) admits a silting object \(T\) for which \(\Gamma_T\) is left noetherian and every finitely generated left \(\Gamma_T\)-module has finite projective dimension.
\end{enumerate}
\end{theorem}

\begin{proof}
Condition~\textup{(1)} is equivalent to \(\operatorname{SObs}_0(\lambda)=0\) by \Cref{thm:negative-k-obstruction}.  \Cref{prop:k0-obstruction} gives essential surjectivity, while \Cref{thm:unconditional-comparison} gives full faithfulness.  Condition~\textup{(2)} implies condition~\textup{(1)}.

Under condition~\textup{(3)}, the dg enhancement supplied by \Cref{prop:compatible-dg-models} determines a small stable enhancement of \(\Kc_\lambda\).  Antieau--Gepner--Heller's Theorem~1.1~\cite{AGH19} gives
\[
 K_{-1}(\Kc_\lambda)=0,
\]
so condition~\textup{(2)} applies.

Assume condition~\textup{(4)}.  The Morita equivalence of \Cref{prop:nonpositive-model} sends the free module \(B_T\) to \(T\), and therefore restricts on compact objects to
\[
 \per(B_T)\xrightarrow{\sim}\thick(T)=\Kc_\lambda.
\]
More explicitly, \Cref{prop:compatible-dg-models} supplies a dg enhancement of the compact kernel.  Since \(T\) classically generates it, dg Morita theory replaces this enhancement by the one-object dga \(A_T=\RHom_R(T,T)^{\op}\).  The silting condition gives \(H^{>0}(A_T)=0\), so the good truncation \(B_T=\tau_{\le0}A_T\to A_T\) is a quasi-isomorphism.  Schlichting's Section~5.9 supplies the finite-cell Frobenius model for \(\per(B_T)\), and Corollary~1.8 gives invariance under this quasi-isomorphism and the induced derived equivalence.  Therefore \Cref{lem:connective-negative-k} gives
\[
 K_{-1}(\Kc_\lambda)
 \cong K_{-1}(\per B_T)
 \cong K_{-1}(\Gamma_T)=0.
\]
Thus condition~\textup{(2)} again applies.  Conditions~\textup{(5)} and~\textup{(6)} follow from \Cref{lem:artinian-negative-k,lem:regular-noetherian-negative-k}.
\end{proof}

\begin{corollary}[Finite-dimensional silting range]\label{cor:finite-dimensional-precompletion}
Let \(R\) be a finite-dimensional \(k\)-algebra and let \(\lambda\) be a stably flat universal localisation.  If \(\Kc_\lambda\) admits a silting object, then
\[
 \rho_\lambda:\DeltaB(\lambda)\xrightarrow{\sim}\Defsg(\lambda)
\]
is an equivalence before idempotent completion.
\end{corollary}

\begin{proof}
Let \(T\) be a silting object.  Since \(T\) is a bounded complex of finite-dimensional projective \(R\)-modules,
\[
 \Gamma_T=\End_{\Kc_\lambda}(T)^{\op}
\]
is finite-dimensional over \(k\), hence artinian.  Apply condition~\textup{(5)} of \Cref{thm:negative-k-vanishing}.
\end{proof}

\begin{corollary}[Idempotent completeness under compact-kernel vanishing]\label{cor:negative-k-completeness}
Under the hypotheses of \Cref{thm:unconditional-comparison}, assume
\[
 K_{-1}(\Kc_\lambda)=0.
\]
In particular, this conclusion applies under conditions \textup{(2)--(6)} of \Cref{thm:negative-k-vanishing}.  Then \(\DeltaB(\lambda)\) and \(\Defsg(\lambda)\) are idempotent complete, and
\[
 \rho_\lambda:\DeltaB(\lambda)\xrightarrow{\sim}\Defsg(\lambda)
\]
is an equivalence before idempotent completion.
\end{corollary}

\begin{proof}
Apply Schlichting's Remark~1.9~\cite{Schlichting06} to the upper exact sequence of Frobenius pairs in \eqref{eq:frobenius-pair-square}.  The categories \(\cC\) and \(\cK\) are already idempotent complete, so the Verdier quotient appearing there is precisely
\[
 \cK/\cC=\DeltaB(\lambda).
\]
The vanishing of \(K_{-1}(\cC)\) therefore implies that \(\DeltaB(\lambda)\) is idempotent complete.  The equivalence follows from \Cref{thm:negative-k-vanishing}; hence \(\Defsg(\lambda)\) is idempotent complete as well.
\end{proof}

\begin{remark}[Equivalence without completeness]\label{rem:injectivity-without-completeness}
Condition~\textup{(1)} of \Cref{thm:negative-k-vanishing} alone proves that \(\rho_\lambda\) is an equivalence, but it does not force either quotient to be idempotent complete.  This is another form of the asymmetry in \Cref{rem:logical-asymmetry}: injectivity kills the relative discrepancy between the two dense subcategories, whereas vanishing of the whole compact-kernel group controls the idempotent completeness of the upper Verdier quotient.
\end{remark}

\begin{theorem}[Residue of a nonzero summand obstruction]\label{thm:negative-k-residue}
Under the hypotheses of \Cref{thm:unconditional-comparison}, suppose
\[
 \operatorname{SObs}_0(\lambda)\ne0.
\]
Then:
\begin{enumerate}[label=\textup{(\arabic*)}]
\item \(K_{-1}(\Kc_\lambda)\ne0\).  More precisely, there is a nonzero class
\[
 0\ne\eta\in K_{-1}(\Kc_\lambda)
\]
whose images in both \(K_{-1}(\Kb(\lambda))\) and \(K_{-1}(R)\) vanish.
\item The compact kernel admits no silting object \(T\) for which \(K_{-1}(\Gamma_T)=0\).  In particular, it admits no silting object with artinian heart algebra, and none with regular noetherian heart algebra in the sense of \Cref{lem:regular-noetherian-negative-k}.
\item If \(R\) is finite-dimensional over a field, then \(\Kc_\lambda\) admits no silting object.
\end{enumerate}
\end{theorem}

\begin{proof}
Assertion~\textup{(1)} is the contrapositive form of \eqref{eq:sobs-negative-k}.  If a silting object as in assertion~\textup{(2)} existed, the proof of condition~\textup{(4)} in \Cref{thm:negative-k-vanishing} would give \(K_{-1}(\Kc_\lambda)=0\), contradicting assertion~\textup{(1)}.  The artinian and regular-noetherian cases follow from \Cref{lem:artinian-negative-k,lem:regular-noetherian-negative-k}.  Finally, if \(R\) is finite-dimensional, the heart algebra of every silting object in \(\Kc_\lambda\subseteq\per(R)\) is finite-dimensional, hence artinian.  Assertion~\textup{(2)} therefore excludes every silting object.
\end{proof}

\begin{remark}\label{rem:no-unrestricted-silting-obstruction}
The implication
\[
 K_{-1}(\Kc_\lambda)\ne0
 \quad\Longrightarrow\quad
 \Kc_\lambda\text{ has no silting object}
\]
is false without a condition on the heart algebra.  For an arbitrary ring \(\Gamma\), the category \(\per(\Gamma)\) has the silting object \(\Gamma\), whether or not \(K_{-1}(\Gamma)\) vanishes.  The unconditional exclusion in \Cref{thm:negative-k-residue} is therefore restricted to the finite-dimensional setting.
\end{remark}

\begin{question}[Counterexample problem]\label{question:negative-k-counterexample}
Does there exist a left noetherian ring \(R\), a stably flat universal localisation
\[
 \lambda:R\longrightarrow R_\Sigma,
\]
and a nonzero class
\[
 0\ne\eta\in K_{-1}(\Kc_\lambda)
\]
such that
\[
 \eta\longmapsto0
 \quad\text{in both}\quad
 K_{-1}(\Kb(\lambda))
 \quad\text{and}\quad
 K_{-1}(R)?
\tag{\(\diamondsuit\)}\label{eq:counterexample-problem}
\]
Equivalently, does there exist a stably flat universal localisation for which
\[
 \operatorname{SObs}_0(\lambda)\ne0?
\]
\end{question}

No such example is constructed here.  The problem is sharper than producing a compact kernel with nonzero negative \(K\)-theory: the same nonzero class must die under both ambient maps in \eqref{eq:counterexample-problem}.  By \Cref{thm:negative-k-residue}, a finite-dimensional example would necessarily have a compact kernel with no silting object.  More generally, a candidate must either have no silting object or have only silting hearts with nonzero \(K_{-1}\), and it must still satisfy the simultaneous-annihilation condition.

\begin{remark}[Scope]\label{rem:negative-k-scope}
The formula \eqref{eq:sobs-negative-k} is stated for the enhanced category \(\Kc_\lambda\) and applies to an arbitrary set \(\Sigma\).  It does not identify this group with the negative \(K\)-theory of a one-object kernel dga unless a one-object dg Morita model has been constructed.  For infinite \(\Sigma\), a silting generator need not exist, and the multi-object intrinsic computation remains open.  The vanishing criteria are sufficient, not necessary: the node in \Cref{thm:node} has an equivalence before completion by direct calculation although its displayed cone generator is not silting.  The explicit artinian-heart computation for the Marks family is given in the proof of \Cref{cor:marks-comparison}.
\end{remark}

\section{Bounded localisation and the corrected absorption theorem}\label{app:bounded-localisation}

We include the bounded quotient argument to make the hypotheses used in the main text explicit.

\begin{theorem}[Bounded localisation equivalence]\label{thm:appendix-bounded-localisation}
Assume \(\lambda:R\to S\) is a homological ring epimorphism satisfying \textup{(B1)--(B3)}.  Then extension of scalars induces an exact equivalence
\[
 \Db(R\text{-}\modcat)/\Kb(\lambda)
 \xrightarrow{\sim}
 \Db(S\text{-}\modcat).
\]
\end{theorem}

\begin{proof}
The unbounded adjunction is
\[
 \Phi=S\Ltensor_R-:D(R)\rightleftarrows D(S):U.
\]
By \Cref{lem:bounded-adjunction}, this adjunction restricts to the bounded finitely generated subcategories.  Homologicality makes \(U\) fully faithful and the counit \(\Phi U\to\id\) an isomorphism.

Let \(q:\Db(R\text{-}\modcat)\to\Db(R\text{-}\modcat)/\Kb(\lambda)\).  For every bounded \(X\), complete the unit to
\[
 K_X\longrightarrow X\longrightarrow U\Phi X\longrightarrow K_X[1].
\]
Applying \(\Phi\) and using the counit shows \(\Phi K_X=0\); boundedness of the other two terms puts \(K_X\) in \(\Kb(\lambda)\).  Therefore \(qX\cong qU\Phi X\).  The counit gives the inverse natural isomorphism on the \(S\)-side.  Thus the induced quotient functor is an equivalence.
\end{proof}

For a finitely generated module \(M\), write
\[
 \Dsg(R,M)=\Db(R\text{-}\modcat)/\thick_{\Db(R)}(M).
\]
This definition avoids any implicit claim that a homotopy category of complexes with terms in \(\add M\) embeds fully in the derived category.

\begin{theorem}[Correct bounded absorption theorem]\label{thm:correct-absorption}
Assume the hypotheses of \Cref{thm:appendix-bounded-localisation}.  Let \(M\) be a finitely generated \(R\)-module and assume
\[
 \Tor_i^R(S,X)=0
 \qquad(i>0,\ X\in\add M).
\]
If
\[
 \Kb(\lambda)\subseteq\thick_{\Db(R)}(M),
\]
then extension of scalars induces an exact equivalence
\[
 \Dsg(R,M)
 \xrightarrow{\sim}
 \Dsg(S,S\otimes_RM).
\]
\end{theorem}

\begin{proof}
The Tor vanishing identifies \(S\Ltensor_RM\) with the ordinary module \(S\otimes_RM\).  Exactness implies
\[
 \Phi\bigl(\thick(M)\bigr)
 \subseteq
 \thick(S\otimes_RM).
\]
Conversely, the latter thick subcategory is generated by the image of \(M\), so it is exactly the thick closure of the image of \(\thick(M)\) under the bounded localisation equivalence.

Since \(\Kb(\lambda)\subseteq\thick(M)\), the third isomorphism theorem gives
\[
 \Db(R)/\thick(M)
 \simeq
 \bigl(\Db(R)/\Kb(\lambda)\bigr)
 \Big/
 \bigl(\thick(M)/\Kb(\lambda)\bigr).
\]
Use \Cref{thm:appendix-bounded-localisation} and identify the denominator with \(\thick(S\otimes_RM)\).
\end{proof}

\begin{corollary}[Defect absorption]\label{cor:defect-absorption}
Under the hypotheses above, suppose
\[
 \Kc_\lambda\subseteq\thick(M)
\]
and the induced functor
\[
 \DeltaB(\lambda)
 \longrightarrow
 \Db(R\text{-}\modcat)/\thick(M)
\]
is zero.  Then \(\Kb(\lambda)\subseteq\thick(M)\), so
\Cref{thm:correct-absorption} applies.
\end{corollary}

\begin{proof}
The quotient functor on \(\Kb(\lambda)\) factors through
\(\Kb(\lambda)/\Kc_\lambda=\DeltaB(\lambda)\).  Its induced functor is zero by assumption, hence every bounded kernel object maps to zero in the quotient by \(\thick(M)\).
\end{proof}

\section{Computational verification protocol}\label{app:computation}

The companion directory \texttt{anc/supplement/} contains five groups of Python scripts in the subdirectories \texttt{part1\_comparison}, \texttt{part2\_fixed\_kernel}, \texttt{part3\_two\_term} \texttt{part4\_optimality} and \texttt{part5\_revision}, their captured outputs, a README recording the interpreter and package versions, an MIT licence for the computational files, and a \texttt{SHA256SUMS} manifest covering every file except the manifest itself.  The top-level runner \texttt{run\_all\_checks.py} executes all fifteen distinct certificates in a temporary copy and compares their outputs with the captured files.  All arithmetic is exact, over \(\mathbb Q\) or a prime field; the two scripts \texttt{verify\_nakayama\_hom\_table.py} and \texttt{verify\_a63\_stable\_kernel.py} use SymPy for exact rational matrices, and every other script uses only the standard library.  The scripts are verification aids; the proofs in the text do not depend on executing them.

\subsection*{Comparison and Nakayama computations}
The file \texttt{verify\_nakayama\_hom\_table.py} implements the path-space calculation in \Cref{thm:nakayama-hom-table}: right modules over \(k\widetilde Q_{2d}/\mathfrak r^d\), cohomological cones \([P_{i+1}\to P_i]\) in degrees \(-1,0\), the middle mapping-complex terms ordered as in \eqref{eq:nakayama-mapping-complex}, and differentials \(\delta^{-1}(f)=(f\alpha_i,\alpha_jf)\), \(\delta^0(a,b)=\alpha_ja-b\alpha_i\).  For every offset and for \(d=3,4,5\) it constructs the two matrices over \(\mathbb Q\), verifies that their composite is zero, computes their ranks exactly, and compares the cohomology dimensions with the theorem; \texttt{verify\_nakayama\_hom\_table\_char2.py} repeats the checks over \(\mathbb F_2\).  The file \texttt{verify\_a63\_stable\_kernel.py} represents every uniserial module \(M(i,\ell)\) of \(\widetilde A^3_6\) as a cyclic quiver representation, solves the arrow-commutativity equations for Hom-spaces, computes the subspace of maps factoring through the six projectives, and verifies the stable Hom formula of \Cref{lem:a63-stable} for shifts \(-8\le q\le8\).  Three further files check the two module-side orientations of the AHKLY presentation (\texttt{verify\_ahkly\_sides.py}), the one-object duality maps of \eqref{eq:one-object-obs} (\texttt{verify\_perfect\_complex\_duality.py}), and the factorisation criterion on small examples (\texttt{verify\_factorisation\_toys.py}).

\subsection*{Fixed-kernel computations}
\texttt{check\_7dim.py}: for \(C=k[x,y]/(x^2,y^2)\), the kernels of \(x\) and \(y\), the Tor groups of \Cref{ex:seven}, the stable endomorphism ring of \(C/(y)\), and the infinite projective dimension of \(k_A\); its variable \texttt{N} denotes the test object \(X=C/(y)\).
\texttt{verify\_dg\_boundary.py}: the totalised semi-free resolutions of \(L_1\) and \(L_4\) over the dg algebra of \Cref{thm:dg-boundary}, the check \(d^2=0\), the cohomology, the nonzero degrees of \(\Ext^q_B(L_1,L_1)\) and \(\Ext^q_B(L_4,L_1)\) for \(0\le q\le32\), and an independent stable-Ext computation over \(A_3\).
\texttt{verify\_nakayama\_profiles.py}: all stable Hom spaces between the uniserial modules of \(A_3\) and \(A_4\) over \(\mathbb Q\), the check of \Cref{lem:stable-hom}(ii) on \(144\) and \(576\) entries, and every realising module of \Cref{prop:m3} over a full shift period.
\texttt{verify\_formal\_realisation.py}: for five pairs \((\Gam,U)\), including a noncommutative \(\Gam\) and a nonprojective bimodule, associativity of the ring \eqref{eq:R-delta} with \(\delta=0\), multiplicativity of \(\lambda\), the two module decompositions, and the strict dg map on products.
\texttt{verify\_group\_block\_obstruction.py}: over \(\mathrm{GF}(3)\), the block decomposition of \(k[C_3\times C_2]\), the periodic resolutions of the two simples, and the vanishing of all cross-Ext groups used after \Cref{thm:pgroup}.
The runner \texttt{run\_checks.py} re-executes these scripts and compares each output byte for byte with the captured file.

\subsection*{Two-term kernel computations}
\texttt{verify\_nonformal.py}, with the helpers \texttt{exact.py} and \texttt{m3.py}: for \Cref{ex:nonformal}, the associativity, bimodule, Peiffer and Leibniz identities of the nine-dimensional dg algebra; the multiplicativity of \(\lambda_\delta\) on all basis pairs; the decompositions \(S_{R}\cong(pR)^{\oplus2}\) and \({}_RS\cong(Re)^{\oplus2}\) on all basis actions; the strict dg map \(D\to\End^\bullet_R(T)^{\op}\) on all products and differentials; and the homotopy transfer of \(m_3\), computed once on \(D\) and once on the presentation-cone algebra, together with the rank of the Hochschild coboundary matrix, which shows that the value \(w\) is not a coboundary.  The captured output is included.


\subsection*{Optimality and Schur-boundary computations}
The additional certificates cover the following finite calculations.  The file \texttt{verify\_connective\_obstruction.py} checks the Frobenius pairings of \Cref{ex:negative-cy-algebras}, the complete endomorphism complex, strict dg map and vertex-component cohomology of \Cref{ex:sphere-schur}, and the analogous sixteen-dimensional tensor Schur example for the exterior algebra.  It also checks the shifted two-term Hom benchmark and an eighteen-dimensional non-formal tensor control, including its Massey-product indeterminacy.  The helpers \texttt{exact.py}, \texttt{m3.py} and \texttt{verify\_nonformal.py} are copies of those in \texttt{part3\_two\_term}.  Complete basis multiplication tables for both Schur rings are supplied as JSON files.  All degrees of every finite complex are checked, with independent rank checks over \(\mathbb F_2,\mathbb F_3,\mathbb F_{101}\); the arbitrary-field arguments are the proofs in the text.  The file \texttt{verify\_marking\_boundary.py} checks the full Hom complex and multiplication of the connective marking in \Cref{ex:marking-boundary}.  Neither script purports to test all ambient rings or to infer homologicality from finitely many Tor groups.  The nonexistence assertions in \Cref{thm:texact-optimality,prop:negative-cy-obstruction} are proved mathematically.

\subsection*{Coordinate symmetry}
The file \texttt{verify\_coordinate\_swap.py} checks the unital algebra involution of \Cref{ex:support-symmetry}, the exchanged module actions, and the image--kernel identities for one repeating step of the free and Hom resolutions over \(\mathbb Q\) and \(\mathbb F_2\). The proof of the infinite resolution and of the induced equivalence of marked cuts is in the text; these finite checks do not replace it.


\clearpage
\bibliographystyle{amsplain-doi}
\bibliography{merged}
\end{document}
